\documentclass{article}

\usepackage[centertags]{amsmath}
\usepackage{amsfonts}
\usepackage{amssymb}
\usepackage{amsthm}
\usepackage{mathtools}
\usepackage{booktabs,longtable,array}
\usepackage{microtype}
\usepackage[hidelinks]{hyperref}

\hypersetup{
  pdftitle={Intrinsic Uniqueness and Reconstruction Across Mathematical Presentations},
  pdfauthor={Alex Albert},
  pdfsubject={Mathematics}
}

\newtheorem{theorem}{Theorem}[section]
\newtheorem{proposition}[theorem]{Proposition}
\newtheorem{lemma}[theorem]{Lemma}
\newtheorem{corollary}[theorem]{Corollary}

\theoremstyle{definition}
\newtheorem{definition}[theorem]{Definition}

\theoremstyle{remark}
\newtheorem{remark}[theorem]{Remark}

\newcommand{\R}{\mathbb R}
\newcommand{\N}{\mathbb N}
\newcommand{\E}{\mathbb E}
\newcommand{\dd}{\,\mathrm d}
\newcommand{\Id}{\mathrm{Id}}
\newcommand{\tr}{\operatorname{tr}}
\newcommand{\KL}{\operatorname{KL}}
\newcommand{\supp}{\operatorname{supp}}
\newcommand{\Var}{\operatorname{Var}}
\newcommand{\crr}{\operatorname{cr}}

\title{
Intrinsic Uniqueness and Reconstruction\\
Across Mathematical Presentations
}

\author{
\href{https://orcid.org/0009-0005-6981-2087}{Alex Albert}\\
Mathematical Research Institute of Physical Reality (MRIPR)\\
P\'ecs, Hungary
}

\date{15 September 2026}

\begin{document}

\maketitle

\begin{abstract}
The complete mathematical presentations established in this work are not
intrinsically distinct objects. They are exact presentations of the same
mathematical object $\Sigma$.

In a placed coordinate, $\Sigma$ is represented by
\[
\Sigma_P(r)
=
\log\!\left(
\frac{\mu^2(1+\gamma r)}{\mu+\gamma}
\right)
-\mu r,
\]
with placement recoverable exactly from local differential data. Removing
placement exposes the intrinsic decomposition
\[
H(t)=\log t-t+1,
\qquad
I(t)=t-1-\log t,
\qquad
p(t)=te^{-t}.
\]

The central result is a global reconstruction theorem. Every complete
identifying presentation established here reconstructs $\Sigma$ uniquely
within its stated candidate class, and $\Sigma$ reconstructs that
presentation exactly. Consequently, the apparent separateness of these
complete presentations is representational rather than mathematical.

The reconstruction is proved across differential, group, convex,
probabilistic, operator-spectral, formal-series, topological,
matrix-geometric, spatial, and arithmetic presentations. Exact inverse
maps and uniqueness theorems are accompanied by explicit countermodels
showing where incomplete data cease to identify $\Sigma$. A relative
irredundancy theorem separates intrinsic structure, placement, and
externally retained context.

The complete development has been formalized in Lean~4 with no admitted
proofs.
\end{abstract}

\section{The intrinsic object and placement}
\label{final:root}

The mathematical object $\Sigma$ is given first in a placed coordinate.
Its placement is recovered exactly, and removing it exposes the
intrinsic decomposition used throughout the paper.

\begin{definition}[The placed presentation of $\Sigma$]
\label{final:C0}
For $\mu>0$, $0<a<1$, put $P=(\mu,a)$, $\gamma=\mu/a$,
$b=-a/\mu$, and $r_*=(1-a)/\mu$. The placed presentation is
\[
 \Sigma_P(r)=\log\!\left(\frac{\mu^2(1+\gamma r)}{\mu+\gamma}\right)-\mu r,
 \qquad r>b.
\]
The intrinsic and placed coordinates, their units, and the parameter
marks are part of these definitions.
\end{definition}

\begin{theorem}[Exact placed reconstruction]
\label{final:C0-placement}
The map $t=\mu r+a$ bijects $(b,\infty)$ with $(0,\infty)$ and gives
\[
 \Sigma_P(r)=\log t-t+\log\frac\mu{1+a}+a.
\]
Inside the placed family, at every point of that domain, placement
is recovered by
\[
 \mu=\sqrt{-\Sigma_P''}-\Sigma_P',\qquad
 \gamma=\frac{\sqrt{-\Sigma_P''}}{1-r\sqrt{-\Sigma_P''}},
 \qquad a=\frac\mu\gamma.
\]
Removing this recovered placement and the value at $r_*$ exposes
the placement-free intrinsic form
\[
 H(t):=\Sigma_P\!\left(\frac{t-a}{\mu}\right)-\Sigma_P(r_*)
      =\log t-t+1,\qquad t>0.
\]
Thus
\begin{equation}
 \Sigma_P(r)=H(\mu r+a)+\log\frac\mu{1+a}+a-1.\label{final:placed}
\end{equation}
Its convex negative $I(t)=t-1-\log t$ and normalized density
$p(t)=te^{-t}$ complete the intrinsic decomposition $(H,I,p)$ of
$\Sigma$, with the exact identities
\begin{align}
 I&=-H,\qquad p=e^{H-1},\qquad H=1+\log p.\label{final:spine}
\end{align}
The closure satisfies $\Sigma_P'(r_*)=0$, $\Sigma_P''(r_*)=-\mu^2$.
The density $p$ has mass one on $(0,\infty)$, whereas
\[
 e^{\Sigma_P(r)}\dd r=\frac{p(t)\dd t}{(1+a)e^{-a}}.
\]
Consequently $e^{\Sigma_P}$ has mass one on $r\ge0$ and mass
$e^a/(1+a)$ on its whole analytic domain.
The presentation $\Sigma_P$ reconstructs both its intrinsic decomposition
$(H,I,p)$ and its placement $P$ exactly; conversely, $(H,I,p)$ together
with $P$ reconstructs $\Sigma_P$ exactly.
\end{theorem}
\begin{proof}
The affine map has positive slope and the stated inverse
$r=(t-a)/\mu$. Substituting $\gamma=\mu/a$ gives
$\mu^2(1+\gamma r)/(\mu+\gamma)=\mu t/(1+a)$, proving
\eqref{final:placed}, with $t=1$ at $r_*$. The logarithm/exponential identities give
\eqref{final:spine}. Integration by parts yields
$\int_a^\infty te^{-t}dt=(1+a)e^{-a}$, including $a=0$.
The change of variables proves both normalization assertions.
Set $q=\gamma/(1+\gamma r)>0$. Differentiation gives
$\Sigma_P'=q-\mu$ and $\Sigma_P''=-q^2$. Solving these equations
gives the recovery formulas, whose denominator is
$(1+\gamma r)^{-1}>0$. At $r_*$ one has $q=\mu$.
\end{proof}

\section{Calibrated analytic reconstruction}
\label{final:core}

\begin{theorem}[Curvature, Riccati and exact flow]
\label{final:C1}
For twice differentiable $F:(0,\infty)\to\R$ with
$F(1)=F'(1)=0$, each of the following conditions identifies $F=H$:
\[
 F''(t)=-t^{-2},\qquad
 F''(t)=-(F'(t)+1)^2.
\]
Equivalently, the marked exact flow
\[
 q(t+s)=\frac{q(t)}{1+s q(t)},\qquad q(1)=1,
\]
for every pair $t,t+s>0$ identifies $q(t)=1/t$; with
$F'=q-1$ and $F(1)=0$ it identifies $F=H$.
More generally $F''(r)=-(r-b)^{-2}$ on $r>b$ identifies
$F(r)=\log(r-b)+Ar+B$, and two independent affine data fix $A,B$.
\end{theorem}
\begin{proof}
Subtract the known logarithmic primitive and use the mean value
theorem twice. For the Riccati equation put $q=F'+1$ and
$W(t)=(tq(t)-1)e^{F(t)+t}$. Direct differentiation using
$q'=-q^2$ gives $W'=0$. The slope anchor gives $W(1)=0$,
so positivity of the exponential gives $tq(t)=1$ everywhere.
Integrating with the value anchor recovers $F$.
For exact flow set $t=1$, $s=v-1$ directly.
All forward assertions follow by differentiating $H$.
Affine perturbations show why the two integration constants are not
fixed by curvature alone.
\end{proof}

\begin{theorem}[Recentring transports its own regularity and anchors]
\label{final:C2}
Let $F$ be differentiable on $(0,\infty)$, with $F''(1)=-1$.
Then
\[
 F(av)-F(a)-aF'(a)(v-1)=F(v)\quad(a,v>0)
\]
holds if and only if $F=H$. No separate affine anchors or global
second-differentiability hypothesis is needed.
\end{theorem}
\begin{proof}
Taking $v=1$ gives $F(1)=0$. Differentiate with respect to $v$ at
one to obtain $F'(1)=0$. At arbitrary $v$ differentiation gives
$a(F'(av)-F'(a))=F'(v)$. The difference quotient at $v=1$ then
gives existence of $F''(a)$ and $a^2F''(a)=F''(1)=-1$.
Theorem~\ref{final:C1} applies. Substitution verifies the converse.
If the scale is not calibrated, every constant multiple of $H$
satisfies the recentering equation.
\end{proof}

\begin{theorem}[Normalized group logarithm and cocycle]
\label{final:C3}
On $D=(-1,\infty)$ define $x\star y=x+y+xy$. This is a
commutative group with identity zero and inverse $-x/(1+x)$.
A function $\ell:D\to\R$ satisfying
\[
 \ell(x\star y)=\ell(x)+\ell(y),\qquad \ell'(0)=1
\]
is uniquely $\ell(x)=\log(1+x)$, where differentiability is
initially required only at zero. Equivalently the differentiable-at-zero
cocycle clause
\[
 h(x\star y)=h(x)+h(y)-xy,\qquad h'(0)=0
\]
identifies $h(x)=\log(1+x)-x=H(1+x)$.
\end{theorem}
\begin{proof}
The identity $1+x\star y=(1+x)(1+y)$ proves closure, associativity,
commutativity and the inverse formula. A homomorphism has $\ell(0)=0$.
For fixed $x$, write $x\star y=x+(1+x)y$, divide its homomorphism
identity by $(1+x)y$, and let $y\to0$. It gives
$\ell'(x)=1/(1+x)$ everywhere. The anchor and the mean value theorem
give the unique primitive. Adding the identity function to the cocycle
turns it into exactly this homomorphism equation. Substitution proves
the forward statements and the cocycle compatibility identity.
\end{proof}

\begin{theorem}[Typed completion equation]
\label{final:C3-completion}
Suppose $h$ is differentiable on $D$, $h'(D)\subset D$ and
$\ell=\mathrm{id}+h$. Then
\[
 \ell(h'(z))=-\ell(z)\quad(z\in D)
\]
holds if and only if, for some $K>0$,
\[
 h(z)=K\log(1+z)-z-\frac K2\log K.
\]
Either anchor $h(0)=0$ or $h'(0)=0$ selects $K=1$.
Injectivity of $\ell$ and twice differentiability of $h$ follow from
the stated domain condition and equation.
\end{theorem}
\begin{proof}
Put $q=h'$. The domain condition gives $\ell'=1+q>0$, so
$\ell$ is strictly increasing and its inverse on its range is
continuous. The equation $q=\ell^{-1}\circ(-\ell)$ makes $q$
continuous. The inverse derivative theorem now makes $q$ continuously
differentiable. Applying the completion equation twice and using
injectivity gives $q(q(z))=z$. Differentiation of the first equation
then gives $(1+z)q'(z)=-(1+q(z))$. Hence
$(1+z)(1+q(z))=K>0$, and integration gives
$h=K\log(1+z)-z+C$. Substitution gives $2C=-K\log K$.
Every displayed function has the requisite derivative range and
satisfies the equation. Its two anchor values are
$-(K/2)\log K$ and $K-1$, each zero exactly when $K=1$.
\end{proof}

\begin{proposition}[Automorphisms, integer series and their calibration]
\label{final:C3-automorphisms}
Every continuous, measurable, monotone or differentiable automorphism
of the fixed group $(D,\star)$ is
\[
 F_c(x)=(1+x)^c-1,\qquad c\in\R\setminus\{0\}.
\]
The calibration $F_c'(0)=1$ selects the identity. Without regularity,
conjugates of arbitrary additive bijections of $\R$ give further
automorphisms. Over a characteristic-zero field the formal
automorphisms have the same formula, with nonzero field coefficient
$c$. The integer series is
$[n]_\star(x)=(1+x)^n-1$; it describes repeated group addition,
and $[m]_\star\circ[n]_\star=[mn]_\star$.
\end{proposition}
\begin{proof}
Conjugate the automorphism by $\log(1+x)$ and $e^u-1$. It becomes
an additive bijection. Each stated regularity condition makes an
additive real function linear: integer and rational linearity follows
algebraically, and continuity follows from measurability or local
boundedness in the other cases. Rational approximation then gives
$u\mapsto cu$, $c\ne0$. A non-linear additive bijection constructed
from a rational basis gives the unrestricted countermodel.
In the formal case, conjugation by the normalized formal logarithm
gives an additive series. Comparing the mixed monomial $u^{n-1}v$
in its additivity identity annihilates every coefficient of degree
$n>1$, since $n$ is invertible. The integer formula follows by the
multiplicative coordinate $1+x$.
\end{proof}

\begin{theorem}[Differentiable Bregman reconstruction]
\label{final:C4}
For a differentiable $\Phi:(0,\infty)\to\R$, let
$B_\Phi(t,s)=\Phi(t)-\Phi(s)-\Phi'(s)(t-s)$.
Simultaneous scale invariance $B_\Phi(ct,cs)=B_\Phi(t,s)$ for all
$c,t,s>0$ holds if and only if
\[
 \Phi(t)=A I(t)+B(t-1)+C.
\]
Thus $\Phi(1)=\Phi'(1)=0$ and
$\Phi'(2)-\Phi'(1)=1/2$ identify $\Phi=I$ without separately
assuming second differentiability or convexity. Alternatively,
for an anchored differentiable $\Phi$ the exact symmetric slice
\[
 B_\Phi(t,1)+B_\Phi(1,t)=t+t^{-1}-2
\]
identifies $\Phi=I$. In particular $B_I(t,s)=I(t/s)$.
\end{theorem}
\begin{proof}
Differentiate scale invariance only in its first argument, taking
$s=1$. For $g=\Phi'-\Phi'(1)$ one obtains
$c(g(ct)-g(c))=g(t)$. Swapping $c,t$ and subtracting gives
$g(c)(1-1/t)=g(t)(1-1/c)$. Fixing $c\ne1$ yields
$g(t)=A(1-1/t)$; integration gives the asserted affine family.
The three calibrations fix its constants. For the symmetric slice,
direct cancellation gives
$(t-1)(\Phi'(t)-\Phi'(1))=t+t^{-1}-2$.
For $t\ne1$ divide by $t-1$ to recover $\Phi'=1-1/t$;
the prescribed derivative at one supplies the remaining point.
Integrate using $\Phi(1)=0$. Direct substitution verifies the forward
identities. Varying $A,B,C$ separately proves the necessity of the
three scalar calibrations within the scale-invariant family.
\end{proof}

\begin{theorem}[Exact Legendre target with either regularity condition]
\label{final:C5}
Extend $I$ by $+\infty$ off $(0,\infty)$. Its convex conjugate is
\[
 I^*(\theta)=\begin{cases}-\log(1-\theta),&\theta<1,\\
 +\infty,&\theta\ge1.\end{cases}
\]
If an extended-real function $F$ has this exact full conjugate and
is either lower semicontinuous on $(0,\infty)$ or convex, then
$F=I$. It is unnecessary to assume both regularity conditions.
Without either, the converse is false.
\end{theorem}
\begin{proof}
For $\theta<1$, the maximizer of $\theta t-I(t)$ is
$t=(1-\theta)^{-1}$, giving the displayed value. The supremum
diverges for $\theta\ge1$. Fenchel's inequality for a candidate
with this conjugate gives $F\ge I^{**}=I$, where the last equality
also follows directly from the just computed maximizers. In particular
$F=+\infty$ off the positive ray. For $t_0>0$, choose
$\theta_0=1-1/t_0$. The function
$(1-\theta_0)t-1-\log t$ is coercive at both endpoints and has
unique minimizer $t_0$. A maximizing sequence in the definition of
$F^*(\theta_0)$ must therefore converge to $t_0$, with its $F$-values
converging to $I(t_0)$. Lower semicontinuity gives the opposite
inequality $F(t_0)\le I(t_0)$. If $F$ is instead convex, these
sequences show its effective domain is dense in the positive ray;
convexity and bracketing by two effective-domain points make it finite
everywhere there. A finite convex function on an open interval is
continuous, so the same argument applies. Finally raise $I(2)$ by
one and leave every other value unchanged. Nearby points approximate
each supremum, so the conjugate is unchanged while the function is
neither convex nor lower semicontinuous at two.
\end{proof}

\begin{theorem}[Self-concordance and the marked Hessian metric]
\label{final:C6}
Let $\Phi\in C^3(0,\infty)$ with $\Phi''>0$,
$\Phi(1)=\Phi'(1)=0$ and $\Phi''(1)=1$. The signed equality
$\Phi'''=-2(\Phi'')^{3/2}$ identifies $\Phi=I$.
The unsigned equality $|\Phi'''|=2(\Phi'')^{3/2}$ also suffices
on the entire positive ray. The corresponding Hessian metric is
$dt^2/t^2$, with distance $|\log(t/s)|$ and
\[
 I(t/s)+I(s/t)=4\sinh^2\!\left(\frac{|\log(t/s)|}{2}\right).
\]
The metric in the marked affine coordinate together with the affine
anchors reconstructs $I$; its bare distance does not fix orientation.
The self-concordance inequality alone does not identify $I$, and
$I$ is not a finite-parameter self-concordant barrier on this domain.
\end{theorem}
\begin{proof}
For the signed equality, $((\Phi'')^{-1/2})'=1$, so calibration gives
$\Phi''=t^{-2}$ and the anchors give $I$. In the unsigned case the
continuous nonzero third derivative has constant sign. The other sign
would give $(\Phi'')^{-1/2}=2-t$, impossible on all $t>0$.
The coordinate $u=\log t$ makes the metric Euclidean, proving the
distance formula; expanding $2\cosh u-2$ proves the symmetrization.
Inversion $t\mapsto1/t$ preserves the distance but reverses the
directed ratio. The quadratic $\Phi(t)=(t-1)^2/2$ satisfies all three
anchors, has positive second derivative, and satisfies the usual
self-concordance inequality since its third derivative is zero,
yet differs from $I$. Moreover,
$(I')^2/I''=(t-1)^2$ is unbounded, refuting a finite barrier parameter.
The affine-equivalent function $-\log t$ does have barrier parameter one.
\end{proof}

\begin{theorem}[Calibrated projective and Schwarzian data]
\label{final:C7}
An injective real-valued function $f$ on $(0,\infty)$ preserving
all cross-ratios and satisfying
$f(1)=0$, $f(2)=-1/2$, $f(3)=-2/3$ equals $1/t-1$.
Consequently $f=F'$ and $F(1)=0$ identify $F=H$.
Alternatively, if $f\in C^3$, $f'\ne0$,
\[
 \frac{f'''}{f'}-\frac32\left(\frac{f''}{f'}\right)^2=0,
 \quad f(1)=0,\quad f'(1)=-1,\quad f''(1)=2,
\]
then the same conclusion holds. Projectivity or zero Schwarzian
without calibration does not select this member.
\end{theorem}
\begin{proof}
With cross-ratio convention
$(x_1-x_3)(x_2-x_4)/((x_1-x_4)(x_2-x_3))$, fix the first three
arguments at $1,2,3$. The cross-ratio is injective fractional-linear
in the fourth argument. Comparing $f$ with $1/t-1$ identifies every
other point; the marks identify the three excluded ones. For the
Schwarzian version put $w=f''/f'$. Its equation is $w'=w^2/2$ with
$w(1)=-2$, whence $w=-2/t$ by ODE uniqueness. Integrate twice
with the two remaining marks to obtain $f'= -t^{-2}$ and
$f=t^{-1}-1$. Integrating the anchored primitive yields $H$.
Without the marks other fractional-linear functions remain.
\end{proof}

\begin{proposition}[Derivative, discrete and local-data boundaries]
\label{final:C8}
The placed derivatives satisfy
\[
 \Sigma_P^{(n)}(r)=(-1)^{n-1}(n-1)!
       \left(\frac\gamma{1+\gamma r}\right)^n,
 \qquad n\ge2.
\]
All orders beyond two add no information once curvature is known.
A lone prescribed derivative of order $n$ identifies a function only
modulo a polynomial of degree at most $n-1$.
For $\rho(r)=C(1+\gamma r)e^{-\mu r}$, $C>0$, its derivative
of order $n\ge1$ has the single zero
$r_n=n/\mu-1/\gamma$. Within this shape family, two consecutive
indexed zeros determine $\mu,\gamma$, and mass determines $C$.
The discrete curvature is
\[
 \Sigma_P(r+1)-2\Sigma_P(r)+\Sigma_P(r-1)
 =\log\left(1-\frac1{(r-b)^2}\right),\quad r-1>b.
\]
It determines lattice values from two starting values, not an arbitrary
real interpolation. A complete germ identifies a global function in a
specified connected real-analytic class, but not in a general smooth class.
\end{proposition}
\begin{proof}
Differentiate the logarithm inductively. Repeated integration proves
the polynomial freedom and Theorem~\ref{final:C1} proves the
second-derivative sufficiency. Leibniz's rule gives
\[
 \rho^{(n)}=C(-\mu)^{n-1}e^{-\mu r}
       \bigl[n\gamma-\mu(1+\gamma r)\bigr],
\]
proving the exact zero. Differences of consecutive zeros recover
$1/\mu$ and then the intercept $-1/\gamma$. Products of the three
logarithmic arguments give the discrete formula; recurrence gives
the lattice-value statement. A nonzero multiple of $\sin(2\pi r)$
preserves all integer samples and changes interpolation. A sufficiently
small compactly supported smooth bump away from an observed point
preserves its entire jet. Two disjoint bumps with cancelling integrals
preserve normalization of a positive density, and can be supported away
from the germ, the mode and endpoints. Analytic continuation instead
uses the identity theorem on the common connected domain.
\end{proof}

\begin{theorem}[Exact all-order derivative-zero sets do not identify $\Sigma$]
\label{final:Z4}
For every real $k>1$ the function
\[
 f_k(t)=(k-1)k^k\frac{t}{(t+k)^{k+1}},\qquad t>0,
\]
is a positive real-analytic probability density, vanishing at both
endpoints, and for every integer $n\ge1$,
\[
 f_k^{(n)}(t)=(-1)^n(k-1)k^k(k)_n
          \frac{t-n}{(t+k)^{k+n+1}}.
\]
Here $(k)_n=k(k+1)\cdots(k+n-1)$. Thus every derivative has
exactly the same single simple positive zero $n$ as $p^{(n)}$,
but $f_k\ne p$. In particular $4t/(t+2)^3$ shows that these
derivative-zero data do not characterize $p$, even among normalized
real-analytic positive densities.
\end{theorem}
\begin{proof}
Substitution $t=ks$ and subtraction of two elementary power integrals
give
\[
 \int_0^\infty \frac{t\,dt}{(t+k)^{k+1}}
 =k^{1-k}\int_0^\infty\frac{s\,ds}{(1+s)^{k+1}}
 =\frac{k^{-k}}{k-1}.
\]
This proves normalization; positivity and real analyticity follow
on the positive ray. The formula is true at order zero if $(k)_0=1$.
Its induction step is the exact identity
\[
 \frac{d}{dt}\frac{t-n}{(t+k)^{k+n+1}}
 =-(k+n)\frac{t-(n+1)}{(t+k)^{k+n+2}}.
\]
Every factor except $t-n$ is nonzero on $t>0$, proving exactness
and simplicity of the zero set, including the derivative sign pattern.
The algebraic tail separates $f_k$ from $te^{-t}$.
The counterexample is not entire and is not globally log-concave:
$(\log f_k)''=-t^{-2}+(k+1)/(t+k)^2>0$ when
$t>1+\sqrt{k+1}$. These additional conditions are not premises of
the refuted characterization. Under $t=\mu r+a$, change variables
and normalize on the selected half-line; multiplication by a positive
constant leaves every placed derivative zero $(n-a)/\mu$ unchanged.
\end{proof}

\section{Probability laws and identifying transforms}
\label{final:probability}

Throughout this section we use the intrinsic decomposition of $\Sigma$
in its marked positive coordinate:
\[
 H(t)=\log t-t+1,\qquad I(t)=t-1-\log t,\qquad
 p(t)=t e^{-t}.
\]
The symbol $\overline F$ denotes a survival function.
All Gamma parameters are shape and
\emph{rate}, in that order.  Equality of probability laws identifies
densities almost everywhere; a pointwise density assertion uses their
specified continuous representatives.  Placement parameters do not
occur in these intrinsic statements.

\begin{proposition}[Uniqueness principles for the transforms used below]
\label{final:P0-uniqueness}
Finite Borel measures on $\mathbb R$ are determined by their
characteristic functions.  Finite measures on $[0,\infty)$ are
determined by their Laplace transforms on $[0,\infty)$, and finite
measures on $[0,1]$ by all their nonnegative integer moments.
Two probability laws with moment generating functions finite on
an open interval about zero and equal there are equal.
\end{proposition}
\begin{proof}
For the compact-interval assertion, equality on monomials gives
equality on polynomials, then on continuous functions by uniform
polynomial approximation, and finally on Borel sets by uniqueness
of finite regular measures.  Pushing a measure on $[0,\infty)$ forward
by $x\mapsto e^{-x}$ reduces Laplace uniqueness to this assertion:
integer Laplace values give all moments of the pushforward on $[0,1]$.

For characteristic functions, subtract the two finite measures.
Convolution with a centered Gaussian of variance $\varepsilon>0$
has density whose Fourier transform is the vanishing characteristic
function times $e^{-\varepsilon\xi^2/2}$.  Fourier inversion, justified
by the integrable Gaussian factor and Fubini's theorem, makes this
convolution zero.  Gaussian approximate identities converge against
every continuous function of compact support, so the original signed
measure is zero.  Finally an exponential moment gives a holomorphic
bilateral transform on a vertical strip about the imaginary axis.
The identity theorem extends equality near zero to that strip;
characteristic-function uniqueness applies on the imaginary axis.
\end{proof}

\begin{theorem}[Gamma law and complete transform characterizations]
\label{final:P1}
Let $\mu_\Gamma(dt)=\mathbf 1_{(0,\infty)}(t)p(t)\,dt$.
It is the probability law $\operatorname{Gamma}(2,1)$, with topological
support $[0,\infty)$ and no atom at zero.  Its transforms are
\begin{align}
 \int e^{-zt}\,\mu_\Gamma(dt)&=(1+z)^{-2},
       &&\operatorname{Re}z>-1,\label{final:eq-gamma-laplace}\\
 \int e^{i\xi t}\,\mu_\Gamma(dt)&=(1-i\xi)^{-2},
       &&\xi\in\mathbb R,\label{final:eq-gamma-characteristic}\\
 \int t^{z-1}\,\mu_\Gamma(dt)&=\Gamma(z+1),
       &&\operatorname{Re}z>-1,\label{final:eq-gamma-mellin}\\
 \int t^n\,\mu_\Gamma(dt)&=(n+1)!,
       &&n\in\mathbb N_0.\label{final:eq-gamma-moments}
\end{align}
Each of the following complete observations identifies its candidate
as $\mu_\Gamma$:
\begin{enumerate}
\item a finite positive Borel measure on $\mathbb R$ with the moments
      \eqref{final:eq-gamma-moments}, including $n=0$;
\item a probability on $[0,\infty)$ with Laplace transform
      $(1+\lambda)^{-2}$ for every $\lambda\geq0$;
\item a probability on $\mathbb R$ with characteristic function
      $(1-i\xi)^{-2}$ for every real $\xi$;
\item a probability on $(0,\infty)$ with
      $\mathbb E T^{i\xi}=\Gamma(2+i\xi)$ for every real $\xi$.
\end{enumerate}
In the first and third clauses, positive support is a conclusion.
Moreover
\[
 \mathbb ET=2,\quad \operatorname{Var}(T)=2,\quad
 \mathbb E\log T=1-\gamma_E,\quad
 \kappa_n(T)=2(n-1)!\quad(n\geq1).
\]
\end{theorem}
\begin{proof}
Euler's integral and integration by parts give
$\int_0^\infty t^a e^{-bt}\,dt=\Gamma(a+1)b^{-a-1}$
for $a>-1$, $b>0$, with its holomorphic extension on the stated
half-planes.  Taking $a=1$ proves normalization and the first two
transforms; taking $a=z$ proves the Mellin formula and the moments.
Differentiating the Gamma integral at $a=1$ is justified by an
integrable logarithmic majorant near zero and exponential decay at
infinity.  It gives $\mathbb E\log T=\Gamma'(2)=1-\gamma_E$.
The identity $\log\mathbb E e^{uT}=-2\log(1-u)$ yields the cumulants.

For the first converse the zeroth moment gives mass one.  For
$0<c<1$, positivity and monotone convergence give
\[
 \int\cosh(ct)\,\mu(dt)
 =\sum_{n=0}^\infty(2n+1)c^{2n}
 =\frac{1+c^2}{(1-c^2)^2}<\infty.
\]
Since $e^{c|t|}\leq2\cosh(ct)$, the candidate has a two-sided
exponential moment.  Dominated power-series expansion for $|u|<1$
therefore gives
\[
 \int e^{ut}\,\mu(dt)
 =\sum_{n=0}^\infty\frac{u^n(n+1)!}{n!}
 =(1-u)^{-2}.
\]
Proposition~\ref{final:P0-uniqueness} identifies the measure on the
whole real line.  Its support follows from the strictly positive
density on $(0,\infty)$ and zero density on $(-\infty,0]$.
The second and third converses are the same uniqueness proposition.
In the fourth, the supplied function is the characteristic function
of $\log T$; identify that law and push it forward by the exponential.
\end{proof}

\begin{proposition}[Exponential samples also derive support]
\label{final:P1-samples}
If a finite positive Borel measure $\nu$ on $\mathbb R$ satisfies
\[
 \int e^{-ns}\,\nu(ds)=(n+1)^{-2}\qquad(n\in\mathbb N_0),
\]
then $\nu(ds)=\mathbf 1_{(0,\infty)}(s)s e^{-s}\,ds$.
\end{proposition}
\begin{proof}
The sample at zero gives mass one.  If
$\nu(( -\infty,-\varepsilon])=m>0$, then the $n$th sample is
at least $m e^{n\varepsilon}$, a contradiction as $n\to\infty$.
Taking a countable union excludes the entire negative ray.
Push forward by $y=e^{-s}$ to $[0,1]$.  Its moments are
$(n+1)^{-2}$, which are those of $(-\log y)\,dy$ on $(0,1)$.
Compact moment uniqueness and the inverse change of variables give
the stated law, including absence of an atom at zero.
\end{proof}

\begin{proposition}[Boundaries of moment identification]
\label{final:P1-boundaries}
Removing positivity, removing the zeroth moment, or retaining only
finitely many ordinary moments invalidates the corresponding
unrestricted moment converse in Theorem~\ref{final:P1}.
\end{proposition}
\begin{proof}
Without the zeroth equation, every $\mu_\Gamma+a\delta_0$, $a>0$,
has the specified moments of positive orders.  To remove positivity,
take a nonzero real even smooth Fourier bump supported in
$[-2,-1]\cup[1,2]$ and let $h$ be its inverse Fourier transform.
The real Schwartz function $h$ has every polynomial moment zero,
because those moments are constant multiples of derivatives of its
Fourier transform at zero.  Thus the distinct finite signed measures
$\mu_\Gamma+\varepsilon h(t)\,dt$ have all the same moments.

For any fixed $m$, choose $m+2$ nonzero smooth bumps with disjoint
compact supports in $(1,2)$.  The $m+1$ linear equations making a
linear combination $h$ orthogonal to $1,t,\ldots,t^m$ have a nonzero
solution.  Disjoint supports ensure $h\neq0$.  On this compact set
$p$ has a positive minimum, so $p+\varepsilon h$ is a positive
probability density for all sufficiently small $|\varepsilon|$.
It is distinct from $p$ and retains every requested moment.
\end{proof}

\begin{theorem}[Marked Poisson recurrence, cumulants and divergences]
\label{final:P2}
A probability sequence $(\pi_n)_{n\geq0}$ satisfies
\[
 (n+1)\pi_{n+1}=\pi_n\qquad(n\geq0)
\]
if and only if $\pi_n=e^{-1}/n!$.  For $N$ with this law,
\begin{align}
 K(u)&=\log\mathbb E e^{uN}=e^u-1,\label{final:eq-poisson-cgf}\\
 \Lambda(u)&=\log\mathbb E e^{u(N-1)}=e^u-1-u=I(e^u).
       \label{final:eq-poisson-centered}
\end{align}
The complete uncentered CGF near zero identifies $N$ as
$\operatorname{Pois}(1)$.  The complete centered expression identifies
the law of $N-1$ when the shift is marked as one.  As a scalar
function it reconstructs $I(t)=\Lambda(\log t)$.

With $D$ denoting relative entropy, for $a,b>0$ one has
\begin{align}
 D(\operatorname{Pois}(a)\Vert\operatorname{Pois}(b))
     &=a I(b/a),\label{final:eq-poisson-kl}\\
 D(\operatorname{Exp}(a)\Vert\operatorname{Exp}(b))
     &=I(b/a),\label{final:eq-exp-kl}\\
 2D(\mathcal N(0,a)\Vert\mathcal N(0,b))
     &=I(a/b),\label{final:eq-gaussian-kl}
\end{align}
where $a,b$ in the last line are variances.  Consequently each full
marked slice obtained by $(a,b)=(1,t),(1,t),(t,1)$, respectively,
reconstructs $I$.  With $\Psi(t)=t\log t-t+1$,
\[
 D(\operatorname{Pois}(a)\Vert\operatorname{Pois}(b))
 =\Psi(a)-\Psi(b)-\Psi'(b)(a-b).
\]
The Legendre transform of the centered CGF is
\[
 \Lambda^*(z)=
 \begin{cases}
 (1+z)\log(1+z)-z,&z>-1,\\
 1,&z=-1,\\
 +\infty,&z<-1.
 \end{cases}
\]
\end{theorem}
\begin{proof}
The recurrence gives $\pi_n=\pi_0/n!$ by induction; summing fixes
$\pi_0=e^{-1}$.  Conversely these coefficients satisfy the recurrence.
Summing the exponential series gives
$\mathbb E e^{uN}=e^{-1}\exp(e^u)$ and both CGFs.  Their converses
follow from Proposition~\ref{final:P0-uniqueness} applied to the actual
variable whose CGF is supplied.

For Poisson laws the log likelihood ratio at $n$ is
$b-a+n\log(a/b)$; its expectation is
$b-a+a\log(a/b)$.  For exponential laws the log ratio is
$\log(a/b)+(b-a)t$, whose expectation under rate $a$ is
$\log(a/b)+(b-a)/a$.  For centered Gaussians the log ratio is
$\frac12\log(b/a)+\frac12t^2(1/b-1/a)$, whose expectation uses
$\mathbb E t^2=a$.  Substitution proves the three formulas and the
Bregman identity.

For $z>-1$, maximizing $uz-\Lambda(u)$ gives $e^u=1+z$ and
the displayed value.  At $z=-1$ the objective is $1-e^u$ with
supremum one as $u\to-\infty$.  For $z<-1$ it tends to infinity
in that limit.  Exponential tilting by $uN$ gives
$\operatorname{Pois}(e^u)$, whose mean, variance, and Fisher
information in the $u$ coordinate are all $e^u$, as follows by
differentiating $K$ and by the score $N-e^u$.
\end{proof}

\begin{proposition}[The centering mark and KL orientation are essential]
\label{final:P2-boundaries}
An unmarked centered CGF does not determine the original location,
even for a law on $\mathbb N_0$.  Reversing the displayed divergence
slices changes their scalar potential.
\end{proposition}
\begin{proof}
Let $Y\sim\operatorname{Pois}(1)$ and $N_k=Y+k$ for an integer
$k\geq1$.  Then $N_k-\mathbb EN_k=Y-1$, with centered CGF
$e^u-1-u$, although $N_k$ is not $\operatorname{Pois}(1)$.
The reversed Poisson slice is $t\log t-t+1$; the reversed exponential
slice is $\log t+t^{-1}-1=I(1/t)$, and twice the reversed Gaussian
slice has this same value.  These follow directly from
\eqref{final:eq-poisson-kl}--\eqref{final:eq-gaussian-kl} and differ
from $I(t)$ in general.
\end{proof}

\begin{theorem}[Two marked max laws characterize the Gumbel CDF]
\label{final:P3}
Let $F$ be a CDF on $\mathbb R$.  The identities
\[
 F(x+\log2)^2=F(x),\qquad F(x+\log3)^3=F(x)
       \quad(x\in\mathbb R),\qquad F(0)=e^{-1}
\]
hold if and only if
\[
 F(x)=\exp(-e^{-x}).
\]
No density, smoothness, strict monotonicity, or full-support
hypothesis is required.  This CDF has density
$g(x)=e^{-x-e^{-x}}=p(e^{-x})$, and it satisfies the analogous
max identity for every positive integer $n$.
\end{theorem}
\begin{proof}
Iteration in both directions gives
$F(k\log2)=\exp(-2^{-k})$ for every integer $k$.
Monotone bounds by adjacent such values yield $0<F(x)<1$ at every
finite $x$.  Hence $G=-\log F$ is finite, positive, and nonincreasing.
The two equations imply
\[
 G(m\log2+n\log3)=e^{-(m\log2+n\log3)}
          \qquad(m,n\in\mathbb Z).
\]
The subgroup generated by $\log2$ and $\log3$ is dense in
$\mathbb R$: otherwise their ratio would be rational, contradicting
unique factorization in the equality $2^r=3^s$ for positive integers
$r,s$.  Approximating an arbitrary $x$ from both sides by subgroup
points and using monotonicity proves $G(x)=e^{-x}$.
Substitution proves the converse, the density formula, and the
max identities.  For independent observations the CDF of their
maximum is $F^n$, so these are precisely the stated shifted max laws.
\end{proof}

\begin{proposition}[Haar weighting and Gumbel calibration boundaries]
\label{final:P3-haar}
For the marked map $\tau(x)=e^{-x}$, the Gumbel probability satisfies
\[
 \tau_*(dF)=e^{-t}\,dt=p(t)\frac{dt}{t},\qquad
 \int_0^\infty p(t)\frac{dt}{t}=1.
\]
The measure $dt/t$ is multiplicative Haar measure and has infinite
mass.  Thus the Haar-weighted probability, and its complete density
in the marked logarithmic coordinate, each reconstruct $p$.
Removing the anchor in Theorem~\ref{final:P3}, or keeping only its
base-two equation, destroys uniqueness.
\end{proposition}
\begin{proof}
The positive change-of-variables Jacobian is $|dx/dt|=1/t$.
Since $g(-\log t)=p(t)$, it gives the pushforward formula;
conversely $p(t)=g(-\log t)$ recovers the density.
Invariance of $dt/t$ under $t\mapsto at$ is immediate by
substitution, and its mass diverges at both endpoints.

Without the anchor, $F_c(x)=\exp(-c e^{-x})$, $c>0$, satisfies
both max equations.  For only the base-two equation, put
$a=\log2$ and $h(x)=1+\varepsilon\sin(2\pi x/a)$.
For nonzero $\varepsilon$ sufficiently small, $h>0$ and $h'<h$.
Then $\exp(-e^{-x}h(x))$ is a strictly increasing CDF with the
correct endpoint limits and anchor; periodicity gives the base-two
equation.  It is not the Gumbel CDF.  The original Gamma probability
$p(t)\,dt$ and the Haar-weighted probability $p(t)\,dt/t$ are
different measures.
\end{proof}

\begin{proposition}[Canonical deficit law and coordinate involution]
\label{final:P4-data}
For $T\sim\mu_\Gamma$ let $V=I(T)$.  For $v\geq0$ let
$a(v)\leq1\leq b(v)$ be the roots of $I(t)=v$; explicitly
\[
 a(v)=-W_0(-e^{-1-v}),\qquad
 b(v)=-W_{-1}(-e^{-1-v}),
\]
where $W_0$ and $W_{-1}$ are the real inverse branches of $w e^w$.
Writing $\overline F_\Gamma(t)=(1+t)e^{-t}$, the canonical deficit
CDF is zero for $v<0$ and is
\begin{equation}
 F_I(v)=\overline F_\Gamma(a(v))-\overline F_\Gamma(b(v)),
                \qquad v\geq0.\label{final:eq-deficit-cdf}
\end{equation}
It has no atom at zero and, for $v>0$, has density
\begin{equation}
 f_I(v)=e^{-1-v}
 \left(\frac{a(v)}{1-a(v)}+\frac{b(v)}{b(v)-1}\right),
 \qquad f_I(v)\sim\frac{\sqrt2}{e\sqrt v}\quad(v\downarrow0).
 \label{final:eq-deficit-density}
\end{equation}
The canonical coordinate involution is
\begin{equation}
 j_I(t)=
 \begin{cases}
 -W_{-1}(-te^{-t}),&0<t<1,\\
 1,&t=1,\\
 -W_0(-te^{-t}),&t>1.
 \end{cases}\label{final:eq-deficit-involution}
\end{equation}
It is a smooth strictly decreasing involution, with $j_I'(1)=-1$,
unique fixed point one, and $I\circ j_I=I$, $p\circ j_I=p$.
\end{proposition}
\begin{proof}
The derivative $I'(t)=(t-1)/t$ and the infinite endpoint limits
give precisely two roots at every positive level.  Rearranging
$I(t)=v$ gives $(-t)e^{-t}=-e^{-1-v}$ and proves the Lambert
formulas.  The sublevel set is $[a(v),b(v)]$, proving the CDF by
integration.  Implicit differentiation gives
$a'=a/(a-1)$ and $b'=b/(b-1)$, while
$p(a)=p(b)=e^{-1-v}$; differentiating the CDF gives its density.
The Taylor expansion $I(1+h)=h^2/2+O(h^3)$ gives
$a(v)=1-\sqrt{2v}+O(v)$ and
$b(v)=1+\sqrt{2v}+O(v)$, yielding the stated asymptotic.
The zero level is the singleton $\{1\}$, which has Gamma mass zero.

Equal-level pairing immediately gives the involution identities.
To verify smoothness at the joining point, write
$I(t)=(t-1)^2 r(t)$ near one, where $r$ is smooth, positive, and
$r(1)=1/2$.  The function
\[
 q_I(t)=\operatorname{sgn}(t-1)\sqrt{2I(t)}
\]
therefore extends smoothly with derivative one at one.  Away from
one its derivative is strictly positive by the sign of $I'$.
The endpoint limits make it a smooth increasing bijection onto
$\mathbb R$, and
$j_I=q_I^{-1}\circ(-\operatorname{id})\circ q_I$.
This proves all the remaining assertions, including $j_I'(1)=-1$.
\end{proof}

\begin{proposition}[Deficit transform, cumulants, and determinacy]
\label{final:P4-cumulants}
The maximal real moment-generating domain of $V$ is $s<1$, and
\begin{align}
 M_V(s)&=e^{-s}\frac{\Gamma(2-s)}{(1-s)^{2-s}},\label{final:eq-deficit-mgf}\\
 K_V(s)&=\gamma_Es+
 \sum_{n=2}^\infty\frac{1}{n}
       \left(\zeta(n)-\frac{1}{n-1}\right)s^n\quad(|s|<1).
       \label{final:eq-deficit-cgf}
\end{align}
Consequently
\begin{equation}
 \kappa_1(V)=\gamma_E,\qquad
 \kappa_n(V)=(n-1)!
       \left(\zeta(n)-\frac{1}{n-1}\right)\quad(n\geq2).
       \label{final:eq-deficit-cumulants}
\end{equation}
In particular $\operatorname{Var}V=\pi^2/6-1$,
$\kappa_3(V)=2\zeta(3)-1$, and $\kappa_4(V)=\pi^4/15-2$.
Every nonnegative probability law with all finite moments and these
algebraic cumulants equals the law of $V$; no separate exponential
moment assumption on that candidate is necessary.
For $a>0$ and $s<1$, the conditional transform is
\begin{equation}
 \mathbb E[e^{sI(T)}\mid T\geq a]
 =\frac{e^{-s}\Gamma(2-s,(1-s)a)}
       {(1+a)e^{-a}(1-s)^{2-s}},
 \label{final:eq-deficit-truncation}
\end{equation}
where $\Gamma(q,x)=\int_x^\infty u^{q-1}e^{-u}\,du$.
\end{proposition}
\begin{proof}
The identity $p=e^{-1-I}$ gives
\[
 M_V(s)=e^{-s}\int_0^\infty t^{1-s}e^{-(1-s)t}\,dt.
\]
For $s<1$ the Gamma substitution gives
\eqref{final:eq-deficit-mgf}.  At $s=1$ its integrand is the
positive constant $e^{-1}$, and for $s>1$ the tail diverges.
Euler's Gamma product, expanded absolutely for $|s|<1$, gives
\[
 \log\Gamma(1-s)=\gamma_Es+
 \sum_{k=1}^\infty\{-\log(1-s/k)-s/k\}
 =\gamma_Es+\sum_{n=2}^\infty\zeta(n)s^n/n.
\]
Using $\Gamma(2-s)=(1-s)\Gamma(1-s)$ and the logarithm series
in $\log M_V$ proves \eqref{final:eq-deficit-cgf}, and differentiation
gives the cumulants.  The classical evaluations
$\zeta(2)=\pi^2/6$, $\zeta(4)=\pi^4/90$ give the displayed values.

Algebraic cumulants and moments determine each other recursively:
\[
 m_0=1,\qquad
 m_{n+1}=\sum_{j=0}^n\binom nj\kappa_{j+1}m_{n-j}.
\]
This is coefficient comparison in the formal identity $M'=K'M$.
Thus a candidate has the same moments as $V$.  For any $0<s<1$,
monotone convergence of its nonnegative exponential series gives
$\mathbb E e^{sX}=\sum m_ns^n/n!=M_V(s)<\infty$.
Nonnegativity also gives finiteness for negative $s$; uniqueness
from Proposition~\ref{final:P0-uniqueness} finishes the converse.
Restricting the same integral to $[a,\infty)$ and dividing by
$\mathbb P(T\geq a)=(1+a)e^{-a}$ proves the conditional formula.
\end{proof}

\begin{theorem}[The linked deficit law and involution identify the potential]
\label{final:P4}
Let $J:(0,\infty)\to[0,\infty)$ be continuous, satisfy $J(1)=0$,
be strictly decreasing on $(0,1]$ and strictly increasing on
$[1,\infty)$, and tend to infinity at both endpoints.  Suppose
$q_J(t)=e^{-1-J(t)}$ has integral one.  Let $F_J$ be the law of
$J(T_J)$ under this same density $q_J$, and let $j_J$ exchange its
two equal-level roots in the same $t$ coordinate and fix one.
If
\[
 F_J=F_I\quad\hbox{and}\quad j_J=j_I,
\]
with the target values specified in
\eqref{final:eq-deficit-cdf} and \eqref{final:eq-deficit-involution},
then $J=I$ and $q_J=p$.  Conversely $I$ supplies these data.
The full list \eqref{final:eq-deficit-cumulants} may replace $F_J=F_I$.
More generally the complete pair $(F_J,j_J)$ uniquely determines
every potential in this candidate class.
\end{theorem}
\begin{proof}
Let $a_J(v)\leq1\leq b_J(v)$ be the inverse branches, including
$a_J(0)=b_J(0)=1$, and put $w_J(v)=b_J(v)-a_J(v)$.
The pushforward of Lebesgue measure under $J$ has locally finite
Stieltjes distribution $w_J$: the sublevel interval has exactly
that length.  Weighting by $e^{-1-J}$ therefore gives, as measures,
\[
 dF_J(v)=e^{-1-v}\,dw_J(v),\qquad
 w_J(v)=e\int_{[0,v]}e^s\,dF_J(s).
\]
Both widths vanish at zero, so there is no undetermined integration
constant.  The map $D_J(t)=t-j_J(t)$ is a continuous strictly
increasing bijection from $[1,\infty)$ to $[0,\infty)$:
$j_J$ is decreasing, fixes one, and tends to zero as $t\to\infty$.
Consequently
\[
 b_J(v)=D_J^{-1}(w_J(v)),\qquad a_J(v)=j_J(b_J(v)).
\]
The observed pair recovers both inverse branches and hence $J$
pointwise.  This proof uses Stieltjes measures and does not require
a derivative of $J$ or a density of $F_J$.
The cumulant alternative follows from
Proposition~\ref{final:P4-cumulants}, since $J(T_J)\geq0$.
\end{proof}

\begin{proposition}[Neither linked observation alone suffices]
\label{final:P4-boundaries}
Within the normalized candidate class of Theorem~\ref{final:P4},
there are noncanonical potentials with exactly $F_I$, and others
with exactly $j_I$.  Thus both canonical target observations and
their common density--potential linkage carry identifying information.
The law of $I(T)$ alone does not identify an arbitrary source law
for $T$, even when the potential $I$ itself is fixed.
\end{proposition}
\begin{proof}
Start with the canonical inverse branches $a,b$.  Choose a nonzero
smooth function $h$ supported on a compact level interval inside
$(0,\infty)$ and set $a_\varepsilon=a+\varepsilon h$,
$b_\varepsilon=b+\varepsilon h$.  On that compact interval the
strictly negative derivative of $a$ and strictly positive derivative
of $b$ are bounded away from zero, and the branches are separated
from zero and one.  Small $|\varepsilon|$ therefore preserves the
branch ranges and monotonicities.  Their inverses define an
admissible $J_\varepsilon\neq I$.  The width is unchanged; the
Stieltjes formula in Theorem~\ref{final:P4} preserves both the entire
deficit law and its total probability mass.  Joint uniqueness shows
that its involution differs.

Conversely choose nonzero nonnegative smooth level bumps $f,g$ with
disjoint compact supports in $(0,\infty)$, and set
$h_{\varepsilon,\delta}(v)=v+\varepsilon f(v)+\delta g(v)$.
Small coefficients make $h$ a strictly increasing bijection fixing
zero, with $h'(0)=1$.  The potential $h\circ I$ has exactly $j_I$.
Its mass functional
\[
 N(\varepsilon,\delta)=
 \int_0^\infty e^{-1-h_{\varepsilon,\delta}(I(t))}\,dt
\]
is smooth near zero, since the perturbations have compact level
support, and
$\partial_\delta N(0,0)=-\int p(t)g(I(t))\,dt<0$.
The implicit function theorem gives $\delta(\varepsilon)$ with
$N(\varepsilon,\delta(\varepsilon))=1$.  For nonzero small
$\varepsilon$, disjoint supports ensure that this potential differs
from $I$, and joint uniqueness makes its deficit law different.

Finally draw $V$ with law $F_I$ and set $T_+=b(V)$.  Then
$I(T_+)=V$ exactly, but $T_+\geq1$ almost surely, so its law is not
$\mu_\Gamma$.  This explicitly exhibits the lost allocation between
the two branches if the source law is not linked to $e^{-1-I}dt$.
Merely possessing some linked pair, without prescribing its two
canonical values, plainly holds for all the potentials just built.
\end{proof}

\begin{theorem}[Survival, hazard, logistic equation, and causal Green kernel]
\label{final:P5}
The Gamma survival and hazard are
\begin{equation}
 \overline F_\Gamma(x)=(1+x)e^{-x}
 =\frac{p(1+x)}{p(1)}=e^{-I(1+x)},\qquad
 h_\Gamma(x)=\frac{x}{1+x}\quad(x\geq0).
 \label{final:eq-survival}
\end{equation}
Each of the following identifying clauses recovers $p$:
\begin{enumerate}
\item an absolutely continuous probability on $[0,\infty)$ with
      positive survival at finite arguments, survival value one at
      zero, and actual hazard $x/(1+x)$ almost everywhere;
\item that same survival framework, with its marked hazard
      $v(u)=h(e^u)$ differentiable on $\mathbb R$ and satisfying
      $v'=v(1-v)$, $v(0)=1/2$;
\item a $C^2$ function $f$ on $(0,\infty)$ extending continuously
      to zero, with $f''+2f'+f=0$, $f(0)=0$, and
      $\int_0^\infty f(t)\,dt=1$;
\item a distribution $g$ on $\mathbb R$ supported in $[0,\infty)$
      with $(D+1)^2g=\delta_0$.
\end{enumerate}
In the last clause the result is
$g(t)=\mathbf 1_{[0,\infty)}(t)t e^{-t}$.
If $k(t)=\mathbf 1_{[0,\infty)}(t)e^{-t}$, then $g=k*k$.
Thus $\mu_\Gamma$ is also the law of a sum of two independent
unit-rate exponential variables and the second arrival time of a
unit-rate Poisson process.
\end{theorem}
\begin{proof}
Integration by parts gives the survival formula, and its negative
derivative divided by the survival gives the hazard.  For the first
converse, local absolute continuity and positivity imply
\[
 (\log\overline F)'=-\frac{x}{1+x}\quad\hbox{almost everywhere}.
\]
Integration from zero gives
$\log\overline F(x)=-x+\log(1+x)$; differentiating recovers the
density $p$ almost everywhere.  The function
$v(u)=(1+e^{-u})^{-1}$ solves the second clause.  Uniqueness for the
scalar differential equation, whose right-hand side is locally
Lipschitz, gives this solution on all of $\mathbb R$; the substitution
$x=e^u$ returns the first clause.

For the third clause, $(e^tf)''=0$, so
$f(t)=(A+Bt)e^{-t}$.  The endpoint gives $A=0$ and integration
gives $B=1$.  For the distributional clause, multiplication by
$e^t$ conjugates $(D+1)^2$ to $D^2$, while
$e^t\delta_0=\delta_0$.  The function $t_+=\max(t,0)$ has
distributional second derivative $\delta_0$.  The difference
$e^tg-t_+$ has second derivative zero, hence is an affine
distribution; its support in $[0,\infty)$ forces both affine
coefficients to vanish.  This proves existence and uniqueness.
Direct convolution gives
\[
 (k*k)(t)=\mathbf 1_{[0,\infty)}(t)
       \int_0^t e^{-s}e^{-(t-s)}\,ds
       =\mathbf 1_{[0,\infty)}(t)t e^{-t}.
\]
The convolution is the independent-sum density.  For a unit-rate
Poisson process the event that its second arrival exceeds $x$ is
the event of at most one arrival by $x$, with probability
$e^{-x}(1+x)$.  This gives the final construction.
\end{proof}

\begin{proposition}[Precise survival boundaries and Gamma shift uniqueness]
\label{final:P5-survival-boundaries}
Local absolute continuity and the survival anchor in
Theorem~\ref{final:P5} cannot be dropped from its almost-everywhere
hazard formulation.  For the logarithmic random variable $Y=\log T$,
the hazard with respect to its own coordinate is
\[
 h_Y(u)=e^u h_\Gamma(e^u)=\frac{e^{2u}}{1+e^u},
\]
which differs from the composition $v(u)$ in that theorem.
For a Gamma density of shape $\alpha>0$ and rate $\beta>0$, the
self-survival shift identity
\begin{equation}
 \overline F_{\alpha,\beta}(x)
 =\frac{f_{\alpha,\beta}(x+c)}{f_{\alpha,\beta}(c)}
                  \quad(x\geq0),\qquad c>0,
 \label{final:eq-gamma-self-shift}
\end{equation}
holds exactly when $\alpha=1$ with arbitrary $c$, or when
$\alpha=2$ and $c=1/\beta$.
\end{proposition}
\begin{proof}
The mixtures $r\delta_0+(1-r)\mu_\Gamma$, $0<r<1$, retain the
hazard at positive arguments and change the survival anchor.
For the regularity boundary, let $C$ be a nonconstant continuous
nondecreasing singular function on $[0,\infty)$, zero on $[0,1]$,
constant on $[2,\infty)$, and with derivative zero almost everywhere.
For instance, translate a Cantor distribution function to $[1,2]$.
Then
\[
 \overline F(x)=(1+x)e^{-x}e^{-C(x)}
\]
is a continuous decreasing survival with value one at zero and
limit zero, and satisfies the same logarithmic derivative equation
almost everywhere.  Its singular component makes it different
from $\overline F_\Gamma$.  This example concerns a bare derivative
equation without absolute continuity.  For $Y=\log T$, the density
is $e^u p(e^u)$ and the survival is
$\overline F_\Gamma(e^u)$, which gives the asserted Jacobian factor.

For \eqref{final:eq-gamma-self-shift}, the density is
$f_{\alpha,\beta}(x)=\beta^\alpha x^{\alpha-1}e^{-\beta x}/
\Gamma(\alpha)$.  If $0<\alpha<1$, the right derivative of the
survival at zero is $-\infty$, while the shifted ratio has finite
derivative, so equality is impossible.  If $\alpha=1$, the ratio
and survival are both $e^{-\beta x}$ for every $c$.
For $\alpha>1$, the survival derivative at zero is zero; equality
forces $(\alpha-1)/c-\beta=0$.  For that $c$, the logarithmic
derivative of the ratio is $-\beta x/(x+c)$, so its derivative is
$-(\beta/c)x(1+o(1))$ near zero.  The survival derivative is
$-\beta^\alpha x^{\alpha-1}(1+o(1))/\Gamma(\alpha)$.
Their equality forces $\alpha-1=1$, then $c=1/\beta$.
The Gamma shape-two survival is $(1+\beta x)e^{-\beta x}$,
which verifies the surviving case.  In particular this shift
property without a specified comparison class is not an
unrestricted characterization of $p$.
\end{proof}

\begin{proposition}[Probability Stieltjes transform and its inverse]
\label{final:P5-stieltjes}
For $z>0$,
\begin{equation}
 G_\Gamma(z)=\int_0^\infty\frac{p(t)}{t+z}\,dt
 =1-z e^z E_1(z),\qquad
 E_1(z)=\int_z^\infty\frac{e^{-u}}u\,du.
 \label{final:eq-gamma-stieltjes}
\end{equation}
The complete transform, or its values on any nonempty open interval
in $(0,\infty)$, identifies a candidate finite positive measure on
$[0,\infty)$ as $\mu_\Gamma$.
The Laplace transform $(1+\lambda)^{-2}$ is completely monotone
but is not a nonzero Stieltjes function.  Neither $p$ nor
$\overline F_\Gamma$ is completely monotone.
\end{proposition}
\begin{proof}
Split $t/(t+z)=1-z/(t+z)$ and substitute $u=t+z$ in the second
integral; this proves \eqref{final:eq-gamma-stieltjes}.
For a finite measure $\mu$, Tonelli's theorem gives
\[
 G_\mu(z)=\int_0^\infty e^{-zv}L_\mu(v)\,dv,
 \qquad L_\mu(v)=\int e^{-vt}\,\mu(dt).
\]
The Stieltjes transform is holomorphic for $\operatorname{Re}z>0$,
by local dominated differentiation.  Equality on an open positive
interval thus extends to all positive $z$.  For two such measures
let $h=L_\mu-L_\nu$, a bounded continuous function.  The finite signed
measure $e^{-v}h(v)\,dv$ has zero Laplace transform by the displayed
identity evaluated at $z+1$.  Proposition~\ref{final:P0-uniqueness}
makes it zero; continuity gives $h=0$, and a second application of
Laplace uniqueness gives $\mu=\nu$.

Complete monotonicity of $(1+\lambda)^{-2}$ follows by differentiating
its positive Laplace integral.  A Stieltjes representation is
\[
 f(\lambda)=a/\lambda+b+
 \int_{[0,\infty)}\frac{\rho(dt)}{\lambda+t},
 \qquad a,b\geq0,\quad
 \int\frac{\rho(dt)}{1+t}<\infty.
\]
If this function is nonzero, either $a>0$, $b>0$, or $\rho$ has
positive mass on some bounded interval.  In every case
$\liminf_{\lambda\to\infty}\lambda f(\lambda)>0$.
For $(1+\lambda)^{-2}$ that limit is zero, a contradiction.
Finally $p'(t)=(1-t)e^{-t}>0$ for $0<t<1$, while
$\overline F_\Gamma''(t)=(t-1)e^{-t}<0$ there.  These signs violate
the respective complete-monotonicity inequalities.
\end{proof}

\begin{proposition}[Reverse size bias, equilibrium, and marked tilting]
\label{final:P5-transforms}
Let $\mu$ be a probability on $[0,\infty)$ with finite positive
mean $m$.
\begin{enumerate}
\item Its size-biased probability is $\beta(dt)=t\mu(dt)/m$.
      If $\mu(\{0\})=0$, then
      \[
      c=\int_{(0,\infty)}t^{-1}\,\beta(dt)=1/m,
      \qquad \mu(dt)=\frac{\beta(dt)}{ct}.
      \]
      Conversely any probability $\beta$ on $(0,\infty)$ with
      $0<c<\infty$ in this formula has that unique positive inverse.
      In particular the positive size-bias inverse of
      $\mu_\Gamma$ is $\operatorname{Exp}(1)$.
\item Its equilibrium law has canonical decreasing right-continuous
      density $q_E(t)=\mathbb P(X>t)/m$ for $t\geq0$.
      If $\mathbb P(X>0)=1$, its inverse is
      \[
      m=1/q_E(0+),\qquad\mathbb P(X>t)=m q_E(t).
      \]
      In particular the full equilibrium density
      $(1+t)e^{-t}/2$ identifies $\mu_\Gamma$, without assuming a
      density for the candidate original lifetime.
\item For a probability $\mu$ on $\mathbb R$, a marked real tilt
      $\theta$ with $0<M(\theta)=\int e^{\theta t}\mu(dt)<\infty$
      gives $\mu_\theta(dt)=e^{\theta t}\mu(dt)/M(\theta)$ and has
      the inverse
      \[
      \mu(dt)=\frac{e^{-\theta t}\mu_\theta(dt)}
                    {\int e^{-\theta y}\mu_\theta(dy)}.
      \]
      For $\mu_\Gamma$, the admissible real tilts are $\theta<1$,
      and
      \[
      p_\theta(t)=(1-\theta)^2t e^{-(1-\theta)t},\quad
      \kappa_n(T_\theta)=\frac{2(n-1)!}{(1-\theta)^n}.
      \]
      Its mean is $2/(1-\theta)$ and variance
      $2/(1-\theta)^2$.
\end{enumerate}
Allowing an atom at zero destroys the first two inverses.  Forgetting
the tilt mark leaves a rate ambiguity in the third.
\end{proposition}
\begin{proof}
In the first clause, integrating $1/t$ against the definition gives
$c=\mu((0,\infty))/m=1/m$.  The inverse has total mass one and
mean $1/c$ by direct integration, and size-biasing it returns
$\beta$.  For $\beta(dt)=p(t)dt$, one has $c=1$, and division by
$t$ gives $e^{-t}dt$.

The tail integral identity $\int_0^\infty\mathbb P(X>t)dt=m$
normalizes the equilibrium density.  Positivity of the original
lifetime gives its right limit at zero as $1/m$; the displayed
inverse therefore recovers the entire survival, including all
positive atoms.  Two right-continuous monotone representatives
that agree almost everywhere agree everywhere, since a strict
difference at one point would persist on a short interval to its
right.  Thus equality of equilibrium laws suffices.  For the target
density $q_E(0+)=1/2$, the recovered survival is
$(1+t)e^{-t}$, whose negative derivative is $p(t)$.

Multiplication of the tilt formula by $e^{-\theta t}$ and
normalization proves the third inverse.  The Gamma transform gives
$M(\theta)=(1-\theta)^{-2}$ for $\theta<1$ and divergence otherwise.
Its tilted log MGF is
$-2\log(1-u/(1-\theta))$, proving the cumulants.

Every $r\delta_0+(1-r)\operatorname{Exp}(1)$, $0\leq r<1$, has
the same size bias $\mu_\Gamma$.  Similarly
$r\delta_0+(1-r)\mu_\Gamma$ has the same equilibrium law for all
$r<1$, because its tail and mean are both multiplied by $1-r$.
For the tilt boundary, a Gamma law of any rate $b>0$, tilted by
$\theta=b-r$, has the fixed output rate $r>0$.  Hence the output
without its input tilt mark does not recover $b$.
\end{proof}

\begin{corollary}[Equilibrium fixed points]
\label{final:P5-equilibrium-fixed}
A positive lifetime with finite mean $m$ equals its equilibrium law
if and only if it is exponential with rate $1/m$.
\end{corollary}
\begin{proof}
Equality to the equilibrium law makes the original distribution
absolutely continuous with density $\overline F/m$.
Consequently $\overline F'=-\overline F/m$ almost everywhere and
$\overline F(0)=1$.  Integration gives $\overline F(t)=e^{-t/m}$.
This exponential law satisfies the same equation and mean by direct
integration.
\end{proof}

\begin{theorem}[Rooted coefficients, Borel probabilities, and the inverse germ]
\label{final:P6}
There is a unique formal power series $R(z)\in z\mathbb R[[z]]$
with $R=ze^R$, and for $n\geq k\geq1$,
\begin{equation}
 [z^n]R(z)^k=\frac{k}{n}\frac{n^{n-k}}{(n-k)!}.
 \label{final:eq-rooted-coefficients}
\end{equation}
In particular
\[
 R(z)=\sum_{n\geq1}\frac{n^{n-1}}{n!}z^n,
 \qquad \operatorname{radius}(R)=e^{-1},
\]
and its compositional inverse germ is $z=t e^{-t}=p(t)$.
The marked Borel PGF and its probabilities are
\begin{equation}
 B(s)=R(s/e),\qquad
 b_n=\frac{e^{-n}n^{n-1}}{n!}\quad(n\geq1),\qquad
 B(s)=s e^{B(s)-1},\quad B(1)=1.
 \label{final:eq-borel}
\end{equation}
Its inverse germ is $s=t e^{1-t}=e p(t)$, with the factor $e$
part of the marking.  For $k$ independent Borel variables their
sum has probabilities
\begin{equation}
 [s^n]B(s)^k=
 \frac{k}{n}e^{-n}\frac{n^{n-k}}{(n-k)!}\qquad(n\geq k).
 \label{final:eq-borel-forest}
\end{equation}
The complete rooted EGF or complete Borel PGF thus determines the
inverse analytic germ.  If a candidate global density is real
analytic on the connected interval $(0,\infty)$ and agrees with
that inverse germ near zero, it equals $p$ everywhere.  Alternatively
$t e^{-t}$ specifies a canonical global extension of the germ.
\end{theorem}
\begin{proof}
Successive coefficient comparison in $R=ze^R$ determines each
coefficient from preceding ones and gives a unique formal solution
with linear coefficient one.  The inverse function theorem applied
to $t e^{-t}$ at zero gives a convergent solution germ, so its
coefficients are those of the formal solution.  For completeness,
the residue change of variables $z=u e^{-u}$ gives
\begin{align*}
 [z^n]R(z)^k
 &=\operatorname*{Res}_{u=0}
       u^{k-n-1}e^{nu}(1-u)\,du\\
 &=[u^{n-k}]e^{nu}-[u^{n-k-1}]e^{nu}
 =\frac{k}{n}\frac{n^{n-k}}{(n-k)!}.
\end{align*}
When $n=k$, the second coefficient is zero and the same formula
gives one.  The coefficient ratio for $k=1$ tends to $e$,
since it is $(1+1/n)^{n-1}$, so the radius is $e^{-1}$.
The inverse identity and Borel scaling follow by substitution.

For $0\leq s<1$, the positive-coefficient solution is the inverse
on $[0,1)$ of the strictly increasing map
$t\mapsto t e^{1-t}$ on $[0,1)$.  Indeed the local inverses agree,
and continuation along this interval is possible because its
derivative is $e^{1-t}(1-t)>0$.  Hence $B(s)\uparrow1$ as
$s\uparrow1$.  Monotone convergence of the nonnegative series
proves $\sum b_n=1$.  Products of PGFs for independent sums and
\eqref{final:eq-rooted-coefficients} give
\eqref{final:eq-borel-forest}.

For labelled rooted trees, removing the root leaves an unordered
set of rooted trees on disjoint sets of labels.  The exponential
formula gives precisely $R=ze^R$; hence their number on $n$ labels
is $n^{n-1}$.  An ordered $k$-tuple of such components has EGF
$R^k$; an unordered set of $k$ components has EGF $R^k/k!$.
These conventions are part of the respective count statements.
Finally equality of the inverse germs and the real-analytic
identity theorem prove the global converse in its stated class.
\end{proof}

\begin{theorem}[Borel total size identifies a one-ancestor iid GW law]
\label{final:P6-gw}
In a Galton--Watson tree started from one ancestor, with independent
identically distributed offspring counts and the usual rooted-tree
construction fixed, the total size has law \eqref{final:eq-borel}
if and only if the offspring law is $\operatorname{Pois}(1)$.
This determines the full tree law within that construction.
\end{theorem}
\begin{proof}
Construct a countable independent family of offspring counts indexed
by all finite words in positive integers, and retain a word as a
vertex exactly when its ancestors and the indicated child indices
are present.  For Poisson offspring the offspring PGF is
$\Phi(w)=e^{w-1}$.  The extinction probability is the increasing
limit of its iterates starting at zero, and therefore the least
fixed point of $\Phi$ in $[0,1]$.  For $0<w<1$ one has
$\log w<w-1$, so $\Phi(w)>w$; also $\Phi(0)>0$.
Its only such fixed point is one.  Extinction is therefore almost
sure, and the tree, having finite offspring counts, has finite
total size almost surely.

For any offspring PGF $\Phi$ in this iid construction, conditioning
on the root offspring and using independence of its descendant
trees gives $B(s)=s\Phi(B(s))$ for the total-size PGF when the total
size is almost surely finite.  For the forward Poisson construction
this is $B=se^{B-1}$, whose formal coefficients are the Borel
probabilities in Theorem~\ref{final:P6}.  Conversely, if the observed
total-size law is Borel, then its PGF satisfies both equations.
The function $B$ maps $(0,1)$ onto $(0,1)$, so division by $s$ gives
$\Phi(w)=e^{w-1}$ for every $0<w<1$.  Analytic coefficient uniqueness
identifies every offspring probability as $e^{-1}/n!$.
The independent product construction then fixes the tree law.
\end{proof}

\begin{proposition}[Tree and continuation boundaries]
\label{final:P6-boundaries}
A complete Borel size law does not identify an arbitrary random
rooted tree.  A complete inverse germ does not identify an arbitrary
smooth positive probability density on $(0,\infty)$, even with the
canonical maximum, strict two-branch shape, and endpoint limits.
Criticality or a leading Borel tail asymptotic alone does not
identify the complete Borel law.
\end{proposition}
\begin{proof}
Sample $N$ with probabilities $b_n$ and return either a rooted path
with $N$ vertices or a rooted star with $N$ vertices.  These laws
have the same size, but on $N=3$ their root degrees differ, and
$b_3>0$.  A Poisson GW tree also has positive probability of root
degree two, which a path never has.  Likewise the PGF product in
\eqref{final:eq-borel-forest} constructs independent components;
the sizes of an independently designated forest need not be
independent merely because component marginals have been specified.

Choose nonzero $h\in C_c^\infty((2,3))$ with $\int h=0$.
For small nonzero $\varepsilon$, set $p_\varepsilon=p+\varepsilon h$.
On the compact support both $p$ and $-p'$ have positive minima,
so positivity and strict decrease survive for small $|\varepsilon|$.
Outside the support it equals $p$.  The integral remains one, the
complete germ at zero is unchanged, and the maximum at one and
all endpoint limits remain canonical.  Nevertheless the density
is different globally.

The critical offspring law assigning probabilities $1/2,1/2$ to
zero and two children has mean one and is not Poisson; its trees
have only odd total sizes.  Finally Stirling's formula gives
$b_n\sim(2\pi)^{-1/2}n^{-3/2}$.  Altering $b_1,b_2$ by opposite
sufficiently small amounts preserves positivity, total mass, and
this exact tail asymptotic, while changing the probability law.
\end{proof}

\begin{theorem}[The calibrated maximum-entropy characterization]
\label{final:P7}
Among Lebesgue probability densities $f$ on $(0,\infty)$ satisfying
\[
 \int t f(t)\,dt=2,\qquad
 \int f(t)\log t\,dt=1-\gamma_E,
 \qquad \int f(t)|\log t|\,dt<\infty,
\]
the unique maximizer of differential entropy is $p$, up to equality
almost everywhere, and its entropy is $1+\gamma_E$.
Here the extended convention allows entropy $-\infty$ and defines
it by the relative-entropy formula below; for finite entropy this
agrees with $-\int f\log f$.
\end{theorem}
\begin{proof}
The target constraints follow from Theorem~\ref{final:P1}.
Since $\log p=\log t-t$, every candidate has the finite cross entropy
$\int f\log p=-1-\gamma_E$.  Put $z=f/p$ and define
\[
 D(f\Vert p)=\int_0^\infty p(t)
          \{z(t)\log z(t)-z(t)+1\}\,dt\in[0,\infty].
\]
The integrand is nonnegative and vanishes exactly when $z=1$.
Because both densities have mass one, this is the usual relative
entropy, and
$\mathsf h(f)=1+\gamma_E-D(f\Vert p)\leq1+\gamma_E$.
For $p$ the divergence is zero; equality for a candidate forces
$f=p$ almost everywhere.  This proves both maximality and its exact
converse.  The two constraints alone do not identify $p$: use four
disjoint compactly supported smooth bumps, solve the three linear
orthogonality equations against $1,t,\log t$, and add a sufficiently
small nonzero resulting perturbation to $p$.  Positivity and all
three equations survive, while the strict entropy inequality holds.
\end{proof}

\begin{theorem}[Unique probability convolution completion]
\label{final:P8}
In the category of probabilities on $[0,\infty)$, there is a unique
family $(\mu_r)_{r\geq0}$ with $\mu_0=\delta_0$,
$\mu_{r+s}=\mu_r*\mu_s$, and $\mu_1=\mu_\Gamma$.
It is
\begin{equation}
 \mu_r=\operatorname{Gamma}(2r,1)\ (r>0),\qquad
 L_{\mu_r}(\lambda)=(1+\lambda)^{-2r}\quad(\lambda\geq0).
 \label{final:eq-gamma-semigroup}
\end{equation}
Weak continuity in $r$ is a conclusion, not an additional premise.
The associated process with stationary independent increments is uniquely
specified in finite-dimensional law once that process category and
the starting value zero are supplied.
\end{theorem}
\begin{proof}
Euler's Gamma integral gives the displayed transforms, and their
products prove the convolution equations by Laplace uniqueness.
For uniqueness let $(\eta_r)$ be another family.  At fixed
$\lambda\geq0$, its Laplace transform is strictly positive and at
most one, so $g(r)=-\log L_{\eta_r}(\lambda)$ is a nonnegative
additive function of $r\geq0$.  It is monotone, since
$g(s)-g(r)=g(s-r)\geq0$ when $s\geq r$.  Rational additivity and
squeezing by rational times yield
$g(r)=r g(1)=2r\log(1+\lambda)$ for every real $r\geq0$.
Another application of Laplace uniqueness gives \eqref{final:eq-gamma-semigroup}.
The continuity theorem for probability Laplace transforms now
gives weak continuity, since these transforms vary continuously
and their pointwise limits are the displayed probability transforms.

For increasing times, take independent increments with the laws
indexed by the successive time differences and form their partial
sums.  Convolution ensures consistency under inserting or deleting
times; the probability extension theorem constructs the process
law, and these prescribed products determine every finite-dimensional
distribution.  In the usual subordinator path convention one takes
its right-continuous increasing version.  The time-one marginal
alone does not specify an independently supplied process: the
process $X_r=rT$ with $T\sim\mu_\Gamma$ has that marginal and
dependent increments, since two unit increments are the same
nonconstant variable $T$.
\end{proof}

\begin{proposition}[Gamma L\'evy measure, drift, and complete Bernstein structure]
\label{final:P8-levy}
The Gamma convolution semigroup has exponent and L\'evy measure
\begin{equation}
 \psi(\lambda)=2\log(1+\lambda)
 =\int_0^\infty(1-e^{-\lambda x})\frac{2e^{-x}}x\,dx,
 \qquad \nu(dx)=\frac{2e^{-x}}x\,dx.
 \label{final:eq-gamma-levy}
\end{equation}
Among probability subordinators with this exact L\'evy measure, the
following are equivalent and each identifies that semigroup:
zero drift; $\psi(\lambda)/\lambda\to0$; time-one support infimum
zero; and time-one mean two.  The L\'evy measure alone leaves arbitrary
nonnegative drift.  Moreover
\begin{equation}
 \frac{2e^{-x}}x=2\int_1^\infty e^{-sx}\,ds,
 \qquad
 \log(1+\lambda)=\int_1^\infty
                 \frac{\lambda}{s(s+\lambda)}\,ds.
 \label{final:eq-complete-bernstein}
\end{equation}
Thus $\psi$ is a complete Bernstein function and
$\psi'(\lambda)=2/(1+\lambda)$ is Stieltjes.
\end{proposition}
\begin{proof}
The integrability condition
$\int(1\wedge x)\nu(dx)<\infty$ follows from boundedness of
$2e^{-x}$ near zero and its exponential tail after division by $x$.
Differentiating the integral in \eqref{final:eq-gamma-levy} with
respect to $\lambda$ gives $2/(1+\lambda)$ by dominated convergence
on compact nonnegative $\lambda$ intervals.  Its value at zero is
zero, so integration proves the exponent formula.

 The L\'evy--Khintchine representation for a subordinator is given in
 \cite[Theorem~3.2, p.~22]{SSV}.  With this measure it is
 $k+d\lambda+2\log(1+\lambda)$, with killing $k\geq0$
and drift $d\geq0$.  Probability mass one excludes killing and
leaves $d\lambda+2\log(1+\lambda)$.  Its time-one law is
$d+T$, by multiplication of Laplace transforms and uniqueness;
its support infimum is $d$, its mean is $d+2$, and the asymptotic
ratio of the exponent to $\lambda$ is $d$.  This proves all the
equivalences and the drift counterexample.  In particular the
complete time-one law in Theorem~\ref{final:P8} already forces
zero drift.

The first identity in \eqref{final:eq-complete-bernstein} is the
elementary exponential integral; the second follows by integrating
$1/s-1/(s+\lambda)$ from one to infinity.  Differentiation of the
first on compact positive $x$ intervals shows that the L\'evy density
is completely monotone.  By definition this makes its Bernstein
exponent complete.  The derivative is the Stieltjes transform of
$2\delta_1$, proving the last assertion.
\end{proof}

\begin{corollary}[Complete integer tails identify the Gamma exponent]
\label{final:P8-samples}
Let $f$ be a Bernstein function.  If, for some integer $N\geq1$,
\[
 f(n)=2\log(1+n)\qquad\hbox{for every integer }n\geq N,
\]
then $f(\lambda)=2\log(1+\lambda)$ for every $\lambda>0$.
\end{corollary}
\begin{proof}
The target is Bernstein by Proposition~\ref{final:P8-levy}.
Apply the complete integer-tail uniqueness theorem
\ref{final:O6} to $f$ and this target.  That theorem derives its
finite compact moment measure from the L\'evy--Khintchine
representation and does not assume a separately given value at zero.
\end{proof}

\begin{proposition}[Explicit Gamma self-decomposition]
\label{final:P8-selfdecomposition}
For $0<c<1$ there exists a nonnegative compound-Poisson residual
$Y_c$, independent of a Gamma copy $T'$, such that
$T\stackrel{d}{=}cT'+Y_c$.  Its L\'evy measure, total jump intensity,
Laplace transform, and mass at zero are respectively
\begin{align}
 \nu_c(dx)&=\frac{2(e^{-x}-e^{-x/c})}{x}\,dx,
 &\nu_c((0,\infty))&=-2\log c,\label{final:eq-residual-levy}\\
 R_c(\lambda)&=\left(\frac{1+c\lambda}{1+\lambda}\right)^2,
 &\mathbb P(Y_c=0)&=c^2.\label{final:eq-residual-transform}
\end{align}
In particular the Gamma law is infinitely divisible and
self-decomposable.  Those two general properties alone do not
identify its shape or rate.
\end{proposition}
\begin{proof}
The density in \eqref{final:eq-residual-levy} is nonnegative.
Since
$e^{-x}-e^{-x/c}=\int_1^{1/c}x e^{-sx}\,ds$, Tonelli gives its
total mass as $2\int_1^{1/c}ds/s=-2\log c$.
A Poisson number of independent positive jumps with normalized law
$\nu_c/(-2\log c)$ therefore constructs $Y_c$.  Its log Laplace
transform is $-\int(1-e^{-\lambda x})\nu_c(dx)$.
Subtracting the two integrals in \eqref{final:eq-gamma-levy}, with
the change of variables for the second, gives
\[
 \int(1-e^{-\lambda x})\nu_c(dx)
 =2\log(1+\lambda)-2\log(1+c\lambda).
\]
Exponentiation proves \eqref{final:eq-residual-transform}; the
mass at zero is the no-jump probability $e^{2\log c}=c^2$.
On a product space take an independent $T'$ with law $\mu_\Gamma$.
The product of the Laplace transforms of $cT'$ and $Y_c$ is
$(1+\lambda)^{-2}$, proving the decomposition by uniqueness.
Infinite divisibility also follows from the $n$-fold convolution
of the time-$1/n$ law in Theorem~\ref{final:P8}.
The identical construction with density
$\alpha e^{-\beta x}/x$, for arbitrary $\alpha,\beta>0$,
proves both properties for every Gamma shape and rate, providing
distinct laws with those same broad properties.
\end{proof}

\begin{theorem}[One linked residual determines a law on the entire real line]
\label{final:P9}
Fix $0\leq c<1$.  Let $Y_c$ have the probability law whose Laplace
transform is \eqref{final:eq-residual-transform}, equivalently the
law of a sum of two independent variables with law
$c\delta_0+(1-c)\operatorname{Exp}(1)$.
For an arbitrary Borel probability $\mu$ on $\mathbb R$, the
following are equivalent:
\begin{enumerate}
\item $\mu=\mu_\Gamma$;
\item there are independent variables $X'$ and $Y_c$, with
      $X'\sim\mu$, such that $cX'+Y_c\sim\mu$;
\item the characteristic function of $\mu$ satisfies
      \[
      \varphi(\xi)=\varphi(c\xi)
          \left(\frac{1-ic\xi}{1-i\xi}\right)^2
                         \qquad(\xi\in\mathbb R).
      \]
\end{enumerate}
No support, density, or moment assumption on $\mu$ is needed.
The residual itself determines $c$ within this family through
$\mathbb P(Y_c=0)=c^2$, so its scale may be recovered from the
complete residual observation.  The source-law linkage is essential;
the contraction endpoint $c=1$ does not identify $\mu$.
\end{theorem}
\begin{proof}
A variable with law $c\delta_0+(1-c)\operatorname{Exp}(1)$ has
Laplace transform
$c+(1-c)/(1+\lambda)=(1+c\lambda)/(1+\lambda)$.
The sum construction proves existence of the residual and its
transform, including $c=0$.  Both summands are nonnegative and
have an atom of mass $c$ at zero, so the sum has an atom of mass $c^2$ at zero.
Its characteristic function is
$((1-ic\xi)/(1-i\xi))^2$, by the same calculation.

Independence gives the factorization in the third clause from the
second.  Conversely that factorization and a product-space
construction, followed by characteristic-function uniqueness, give
the second clause.  Iterating the factorization for any integer
$n\geq1$ telescopes:
\[
 \varphi(\xi)=\varphi(c^n\xi)
       \left(\frac{1-ic^n\xi}{1-i\xi}\right)^2.
\]
The denominators are nonzero for real $\xi$.  Since $c^n\xi\to0$
and every probability characteristic function is continuous at
zero with value one, letting $n\to\infty$ gives
$\varphi(\xi)=(1-i\xi)^{-2}$.  Theorem~\ref{final:P1} and
characteristic-function uniqueness now identify $\mu$ on all of
$\mathbb R$.  For the converse, the Gamma characteristic function
satisfies the factorization by direct cancellation, and independent
variables with the required laws can be constructed on a product
space.  This also covers $c=0$, where $Y_0$ is already Gamma.

Forgetting linkage permits the very same observed $Y_c$ alongside
an unrelated designated $X$ with either Gamma shape two or Gamma
shape three, so it identifies neither such designation.  At $c=1$
the residual is $\delta_0$ and the factorization reduces to
$\varphi(\xi)=\varphi(\xi)$ for every law.  Thus contraction and
the equation linking the original law to this precise residual
are substantive identifying premises.  Independence may be
replaced by the exact factorization as stated, but that
factorization does not follow for an arbitrary supplied coupling.
\end{proof}

\begin{remark}[Formalization]
\label{final:P-formalization-scope}
The probability development, including its measure-theoretic and
analytic statements and counterexamples, is part of the complete
Lean~4 formalization with no admitted proofs. The Bernstein
integer-tail inverse used above is Theorem~\ref{final:O6}.
\end{remark}

\section{Stein characterization and the marked Laguerre realization}
\label{final:operators}

Throughout this section, $p(t)=te^{-t}$ is the density in the intrinsic
decomposition of $\Sigma$, in the specified real coordinate $t>0$.
The complete identifying data below reconstruct $\Sigma$ with their
coordinate and operator context retained. The candidate classes for a
coordinate probability, an abstract self-adjoint operator, and a
semigroup mixing measure have their respective identifying equations.
All Hilbert spaces are complex.  We use $\mathrm{Id}$ for the identity
operator, and reserve $I$ for the scalar potential.  The arguments give
analytic proofs with explicit domains and candidate classes; numerical
audits are not used to establish these universal statements.

\begin{theorem}[The full-line weak Stein characterization]
\label{final:O1}
Let $P$ be a positive Borel probability measure on $\mathbb R$.  Then
\begin{equation}\label{final:O1-identity}
 \int_{\mathbb R}\bigl[t f'(t)+(2-t)f(t)\bigr]P(dt)=0
 \qquad\bigl(f\in C_c^\infty(\mathbb R)\bigr)
\end{equation}
if and only if
\[
 P(dt)=\mathbf 1_{(0,\infty)}(t)te^{-t}\,dt.
\]
In particular, absolute continuity, positive support, absence of an atom
at zero, and all finite moments are conclusions, not hypotheses.
\end{theorem}

\begin{proof}
The integrands in \eqref{final:O1-identity} have compact support, so no
moment assumption is needed to state the identity.  In distributions it
says $(tP)'=(2-t)P$.  On either open half-line put $Q=tP$.  There
\[
 Q'=(2/t-1)Q,\qquad
 \bigl(e^t t^{-2}Q\bigr)'=0.
\]
A distribution with zero derivative on a connected interval is a constant
distribution.  Thus the restrictions of $P$ have densities
$C_+te^{-t}$ on the positive half-line and $C_-|t|e^{-t}$ on the negative
half-line, with $C_+,C_-\geq0$.  Finiteness of $P$ forces $C_-=0$.
The only remaining possible part is $m\delta_0$, $m\geq0$.
The positive-half-line flux $t^2e^{-t}$ tends to zero at the origin;
its distributional derivative therefore has no atom there.  Substitution
of $m\delta_0$ into the equation gives $0=2m\delta_0$, so $m=0$.
Finally $\int_0^\infty te^{-t}\,dt=1$ gives $C_+=1$.
Conversely, integration by parts, using
$(t^2e^{-t})'=(2-t)te^{-t}$ and the zero endpoint flux, proves the
displayed identity.  Tests supported only in $(0,\infty)$ could not
exclude mass on $(-\infty,0]$; the full-line test class is essential
to this support conclusion.
\end{proof}

\begin{proposition}[The integrable centered Stein kernel]
\label{final:O1-kernel}
Fix the probability density $p$, and let $\tau$ be real measurable with
$\tau p\in L^1_{\mathrm{loc}}(0,\infty)$.  The equation
\[
 (\tau p)'=(2-t)p
\]
in distributions is equivalent to
\begin{equation}\label{final:O1-kernel-family}
 \tau(t)=t+C\frac{e^t}{t}\quad\hbox{almost everywhere},
 \qquad C\in\mathbb R.
\end{equation}
The locally absolutely continuous representative of $\tau p$ has both
endpoint limits equal to $C$.  Each of the following separately selects
$\tau=t$ almost everywhere: zero flux at either endpoint;
$\tau\in L^1(p\,dt)$; a finite expectation of $\tau$ under $p\,dt$; or
the finite identity $\int\tau p=\operatorname{Var}_p(t)=2$.
Nonnegativity alone permits exactly $C\geq0$.
\end{proposition}

\begin{proof}
Write $q=tp=t^2e^{-t}$.  Since $q'=(2-t)p$, the distribution
$\tau p-q$ has derivative zero and is constant.  This proves
\eqref{final:O1-kernel-family} and the assertions about its absolutely
continuous representative.  For $C\ne0$, $q+C$ tends to $C$ at infinity,
so $\int_0^\infty|\tau|p=\int_0^\infty|q+C|=\infty$; its tail also
has the constant sign of $C$, so $\int \tau\,p\,dt$ cannot be finite as
a Lebesgue integral.  For $C=0$, integration by parts gives
$\int tp=2$, $\int t^2p=6$, and hence variance $2$.
As $q>0$ and $q$ tends to zero at both endpoints, $q+C\geq0$
almost everywhere exactly when $C\geq0$.

Every value of $C$ satisfies the compact-test identity
\[
 \int\tau f'p\,dt=\int(t-2)fp\,dt
 \qquad(f\in C_c^\infty(0,\infty)),
\]
because the added term is $C\int f'=0$.  Including $f(t)=t$, with
both integrals required finite, instead gives $\int\tau p=6-4=2$.
This inverse problem fixes $P$ and varies $\tau$; it is different from
Theorem~\ref{final:O1}, which fixes the coefficient $t$ and varies $P$.
\end{proof}

\begin{theorem}[Two marked probes and the Pearson realization]
\label{final:O2}
In the conservative local second-order class
\[
 Bf=-a(t)f''-b(t)f',\qquad t>0,
\]
the two differential-expression identities
\begin{equation}\label{final:O2-probes}
 B(2-t)=2-t,\qquad
 B(t^2/2-3t+3)=t^2-6t+6
\end{equation}
hold if and only if $a(t)=t$ and $b(t)=2-t$.
The conclusion is pointwise or almost everywhere according to the
meaning of the coefficient identities.  If a positive, locally
absolutely continuous probability density $w$ is also required to obey
the zero-current Pearson equation $(aw)'=bw$, then $w=p$.
Thus these data determine the differential expression and its normalized
Hilbert weight in this specified coordinate class.
\end{theorem}

\begin{proof}
The first probe gives $b=2-t$.  The second then gives
\[
 -a-(2-t)(t-3)=t^2-6t+6,
\]
and consequently $a=t$.  Substitution proves the converse.
The Pearson equation reduces to
$tw'=(1-t)w$, so $(e^t w/t)'=0$ locally and $w=Cte^{-t}$.
Probability normalization gives $C=1$.
The probes are polynomials; they are identities of the expression,
or identities on an explicitly enlarged domain containing those
polynomials.  They are not assertions that polynomials belong to
$C_c^\infty(0,\infty)$.
\end{proof}

\begin{proposition}[The exact limits of the probe and stationary data]
\label{final:O2-boundaries}
The pair \eqref{final:O2-probes} is irredundant within the indicated
coefficient class, even if the coefficients are smooth and $a>0$.
It does not identify an arbitrary self-adjoint operator outside that
local class.  The invariant density $p$ alone does not identify a
reversible dynamics or its time scale.
\end{proposition}

\begin{proof}
With only the linear probe, set $b=2-t$ and choose any positive smooth
$a$.  With only the quadratic probe, take a nonzero
$k\in C_c^\infty(0,\infty)$ and set
\[
 b=2-t+\varepsilon k(t),\qquad
 a=t-\varepsilon k(t)(t-3).
\]
For sufficiently small $|\varepsilon|$, positivity holds on the compact
support of $k$ and elsewhere $a=t$; the quadratic identity holds
identically while the linear one changes.
For the larger self-adjoint category, use the operator and eigenbasis
constructed below.  Let $U$ fix $e_0,e_1,e_2$ and every $e_n$, $n\geq5$,
and make a nontrivial rotation, with angle in $(0,\pi/2)$, on
$\operatorname{span}\{e_3,e_4\}$.  Then $UAU^{-1}\ne A$, since the
two rotated eigenvalues are distinct, but it preserves the constant
vector and both probes.  This example has self-adjointness; positivity
preservation of the rotated semigroup is not imposed.
Finally $cA$, $c>0$, gives distinct conservative reversible semigroups
with the same invariant $p\,dt$, as follows from the construction below.
The centered integrable Stein selection fixes $c=1$; invariance alone
does not.
\end{proof}

\subsection{The actual self-adjoint domain and complete spectral resolution}

Set
\[
 \mathcal H=L^2((0,\infty),p(t)\,dt),\qquad q(t)=tp(t)=t^2e^{-t},
 \qquad \langle f,g\rangle=\int_0^\infty f\overline g\,p\,dt.
\]
In particular $\langle\cdot,\cdot\rangle$ is linear in its first
argument.  The minimal differential operator considered here is
\begin{equation}\label{final:O3-minimal}
 A_0f=-p^{-1}(qf')'=-tf''-(2-t)f',\qquad
 D(A_0)=C_c^\infty(0,\infty).
\end{equation}

\begin{theorem}[Essential self-adjointness and the full domain]
\label{final:O3}
The operator $A_0$ is densely defined, nonnegative and symmetric.
Both singular endpoints are limit point, its deficiency indices are
$(0,0)$, and its unique self-adjoint extension is
$A=\overline{A_0}=A_0^*$.  Its exact domain is
\begin{equation}\label{final:O3-domain}
 D(A)=\left\{f\in\mathcal H:
 f,qf'\in AC_{\mathrm{loc}}(0,\infty),\quad
 -p^{-1}(qf')'\in\mathcal H\right\}.
\end{equation}
Here the regularity refers to the indicated representatives of the
$L^2$ classes.  For every $f\in D(A)$,
\begin{equation}\label{final:O3-flux}
 \lim_{t\downarrow0}q(t)f'(t)=
 \lim_{t\to\infty}q(t)f'(t)=0.
\end{equation}
These limits are consequences of the domain, not additional boundary
conditions.  In particular no finite boundary value $f(0)$ is required.
\end{theorem}

\begin{proof}
Smooth compactly supported functions are dense in this weighted $L^2$
space, because the weight is positive and smooth on each compact
interior interval.  Integration by parts gives
\[
 \langle A_0f,g\rangle=\int_0^\infty qf'\overline{g'}\,dt,
\]
which proves symmetry and nonnegativity.
The definition of the adjoint says that $f\in D(A_0^*)$ precisely when
there is an $h\in\mathcal H$ with $-(qf')'=ph$ in distributions.
The right side is locally integrable.  Taking a primitive, dividing
by the positive smooth function $q$, and taking a second primitive
shows that $qf'$ and $f$ have locally absolutely continuous
representatives.  Conversely those representatives permit integration
by parts against each compact test.  Thus the adjoint is exactly the
maximal operator in \eqref{final:O3-domain}.

To determine the endpoints, two independent zero-energy solutions are
\[
 u_1(t)=1,\qquad u_2(t)=\int_1^t\frac{e^s}{s^2}\,ds.
\]
Expanding $e^s/s^2=s^{-2}+s^{-1}+O(1)$ near zero and using
l'H\^opital's rule at infinity gives
\[
 u_2(t)=-t^{-1}+\log t+O(1)\quad(t\downarrow0),\qquad
 u_2(t)\sim e^t/t^2\quad(t\to\infty).
\]
Consequently $|u_2|^2p\sim t^{-1}$ at zero and
$|u_2|^2p\sim e^t/t^3$ at infinity.  Neither is integrable.
The constant solution is square integrable near each endpoint.
The Weyl endpoint alternative therefore makes both endpoints limit
point: an endpoint is limit circle precisely when all solutions at
one, equivalently every, spectral parameter are square integrable
there.  The deficiency-index theorem gives one deficiency at each
limit-circle endpoint and zero at each limit-point endpoint.
Its hypotheses hold here because $p,q>0$, $p,q^{-1}$ are locally
integrable, and the potential is zero; see
Gesztesy--Littlejohn--Nichols~\cite[Theorems 3.6 and 3.8]{GLN}.
Their minimal domain is the closure of compactly
supported maximal-domain elements.  It has the same closure as
$C_c^\infty$ here: on compact interior intervals those elements are
in $H^2$, and smooth approximation converges in the graph norm.
This verifies the application to the actual domain
\eqref{final:O3-minimal}, including the critical logarithmic divergence
at the origin.  It gives $\overline{A_0}=A_0^*$ and deficiency indices
$(0,0)$.  Nonnegativity passes to this closure.

For the flux statement put $h=Af$.  Cauchy--Schwarz and
$\int p=1$ give $\int|h|p\leq\|h\|_{\mathcal H}$.
Thus $(qf')'=-ph$ is globally integrable and $qf'$ has finite limits
at both endpoints.  A nonzero limit $\ell$ at zero would yield
$f(t)\sim-\ell/t$, by integrating $f'=(\ell+o(1))/q$;
a nonzero limit at infinity would give
$f(t)\sim\ell e^t/t^2$.  The same weighted-square divergences exclude
both.  This proves \eqref{final:O3-flux}; in fact
\[
 q(t)f'(t)=-\int_0^t hp\,ds=\int_t^\infty hp\,ds,
 \qquad \int_0^\infty hp\,dt=0.
\]
To see why finiteness of $f(0)$ cannot be added, multiply
$\log\log(1/t)$ by a smooth cutoff supported sufficiently near zero
and equal to one there.  This function is in $\mathcal H$, and near
zero its image under the differential expression is
$O((t\log(1/t))^{-1})$.  Its weighted square integral is finite,
since $\int_0^\varepsilon dt/(t\log^2(1/t))<\infty$.
It therefore belongs to \eqref{final:O3-domain} and has no finite
value at zero.
\end{proof}

\begin{proposition}[The closed form and the conservative semigroup]
\label{final:O3-form}
The closed form of $A$ is
\[
 \mathcal E(f,g)=\int_0^\infty qf'\overline{g'}\,dt,
\]
with the exact domain
\begin{equation}\label{final:O3-form-domain}
 D(A^{1/2})=
 \left\{f\in\mathcal H:f\in AC_{\mathrm{loc}}(0,\infty),\quad
 \int_0^\infty q|f'|^2\,dt<\infty\right\}.
\end{equation}
There are no additional endpoint traces.  The semigroup $e^{-sA}$,
$s\geq0$, is positivity preserving, preserves the constant function
$1$, and preserves the probability $p\,dt$.
\end{proposition}

\begin{proof}
Denote the displayed form domain by $V$.  A Cauchy sequence in
$\|f\|_V^2=\|f\|_{\mathcal H}^2+\int q|f'|^2$ is Cauchy in
ordinary $H^1$ on every compact interior interval, where both weights
are bounded above and below by positive constants.  Its local weak
derivative equals the weighted derivative limit, proving completeness.
Truncate real and imaginary parts of $f\in V$ at height $N$.
The weak chain rule and dominated convergence show convergence in
$V$ as $N\to\infty$.  For bounded $f$, choose cutoffs which rise on
$[\varepsilon,2\varepsilon]$, fall on $[R,2R]$, and have derivatives
bounded respectively by $C/\varepsilon$ and $C/R$.  Their energy is
bounded by
\[
 C\varepsilon+\frac{C}{R^2}\int_R^{2R}t^2e^{-t}\,dt\longrightarrow0.
\]
The two derivative errors in $\chi f-f$ are
$(\chi-1)f'$ and $\chi'f$; the first tends to zero by dominated
convergence, the second by this estimate and boundedness of $f$.
The $\mathcal H$ norm also tends to zero.  On the resulting compact
interior support, mollification converges in $H^1$ and hence in $V$.
Thus $C_c^\infty$ is a form core, and the displayed form is the
closure of the minimal form.  Its represented self-adjoint operator
extends $A_0$ and must equal $A$ by Theorem~\ref{final:O3}.

The weak chain rule shows that real normal contractions preserve $V$
and decrease its energy.  Equivalently, the variational resolvent is
positivity preserving and maps $[0,1]$ into $[0,1]$, as follows by
testing its equation with the negative part and with $(u-1)^+$.
The exponential resolvent limit gives the same Markov property for
$e^{-sA}$.  Finally $1\in D(A)$ and $A1=0$, so $e^{-sA}1=1$.
Self-adjointness gives
$\int e^{-sA}f\,p\,dt=\langle e^{-sA}f,1\rangle
=\langle f,e^{-sA}1\rangle=\int fp\,dt$.
\end{proof}

\begin{theorem}[The Laguerre basis and the exact integer spectrum]
\label{final:O3-spectrum}
Define the polynomials with their marked normalization by
\begin{equation}\label{final:O3-laguerre}
 L_n^{(1)}(t)=\frac{e^t}{n!t}\frac{d^n}{dt^n}
       (e^{-t}t^{n+1})
 =\sum_{k=0}^n(-1)^k\binom{n+1}{k+1}\frac{t^k}{k!},
 \qquad n\in\mathbb N_0.
\end{equation}
Then $e_n=L_n^{(1)}/\sqrt{n+1}$ is a complete orthonormal basis of
$\mathcal H$, each $e_n\in D(A)$, and $Ae_n=ne_n$.
In particular the spectrum is exactly $\mathbb N_0$, each eigenvalue
is simple, and the resolvent is compact.  If $f=\sum c_ne_n$, then
\begin{align}
 D(A)&=\left\{f:\sum_{n=0}^\infty n^2|c_n|^2<\infty\right\},
 &Af&=\sum_{n=0}^\infty nc_ne_n,\label{final:O3-spectral-domain}\\
 D(A^{1/2})&=\left\{f:\sum_{n=0}^\infty n|c_n|^2<\infty\right\},
 &\mathcal E(f,f)&=\sum_{n=0}^\infty n|c_n|^2.
 \nonumber
\end{align}
\end{theorem}

\begin{proof}
Leibniz's formula gives the finite sum in
\eqref{final:O3-laguerre}.  Its coefficients $a_k$ satisfy
$(k+1)(k+2)a_{k+1}+(n-k)a_k=0$, so
\[
 t(L_n^{(1)})''+(2-t)(L_n^{(1)})'+nL_n^{(1)}=0.
\]
Every polynomial belongs to the maximal domain: the polynomial and
its polynomial image are in $\mathcal H$, and all local regularity
conditions hold.  These are therefore genuine eigenfunctions of $A$.
For $m<n$, integrating the Rodrigues formula $n$ times by parts gives
$\int pL_n^{(1)}t^m\,dt=0$.  The boundary terms vanish at infinity
by exponential decay and at zero because derivatives of
$e^{-t}t^{n+1}$ of order at most $n-1$ vanish to order at least two.
For $m=n$ the integral is $(-1)^n\Gamma(n+2)$.
The leading coefficient $(-1)^n/n!$ then gives
\[
 \int_0^\infty L_n^{(1)}L_m^{(1)}p\,dt=(n+1)\delta_{nm}.
\]

For completeness, let $g\in\mathcal H$ be orthogonal to every
polynomial and define
$F(z)=\int_0^\infty g(t)p(t)e^{zt}\,dt$.
Cauchy--Schwarz gives locally uniform domination of this integral
and of all its differentiated integrands on $\operatorname{Re}z<1/2$.
Thus $F$ is holomorphic there, and $F^{(n)}(0)=0$ for all $n$.
The identity theorem gives $F=0$ on that half-plane.
The finite complex measure $gp\,dt$, pushed forward by $y=e^{-t}$
and regarded as a measure on $[0,1]$, consequently has every moment
zero, since $F(-k)=0$ for $k\in\mathbb N_0$.
Uniform polynomial approximation on $[0,1]$ shows that it integrates
every continuous function to zero; measure uniqueness gives $g=0$.
This proves completeness.
The unitary coefficient map now takes $A$ to multiplication by $n$
on $\ell^2(\mathbb N_0)$, with exactly the displayed domain.
The form domain follows by taking square roots.  The diagonal
entries $(n+\alpha)^{-1}$ of any resolvent, $\alpha>0$, tend to zero,
proving compactness and the asserted spectrum.
\end{proof}

\begin{theorem}[Complete marked Laguerre orthogonality identifies the law]
\label{final:O7}
Let $P$ be a positive Borel probability measure on $\mathbb R$ for
which every polynomial is integrable.  With the specific polynomials
\eqref{final:O3-laguerre},
\[
 \int_{\mathbb R}L_n^{(1)}(t)\,P(dt)=0\quad(n\geq1)
\]
holds if and only if $P(dt)=\mathbf1_{(0,\infty)}p(t)\,dt$.
The positive support is again a conclusion.
\end{theorem}

\begin{proof}
Mass fixes moment zero.  The nonzero leading coefficient
$(-1)^n/n!$ means that the $n$th displayed equation determines moment
$n$ from all smaller moments.  The preceding Rodrigues calculation
proves these equations for $p\,dt$, whose moments are $(n+1)!$.
Induction therefore gives $\int t^nP(dt)=(n+1)!$ for all $n\geq0$.
For $0<c<1$, monotone convergence in the nonnegative even-power
series gives
\[
 \int\cosh(ct)\,P(dt)
 =\sum_{n=0}^\infty(2n+1)c^{2n}
 =\frac{1+c^2}{(1-c^2)^2}<\infty.
\]
As $e^{c|t|}\leq2\cosh(ct)$, the bilateral moment-generating
function is holomorphic on $|\operatorname{Re}z|<1$.
Its Taylor series near zero, determined by the moments, is
$\sum_{n\geq0}(n+1)z^n=(1-z)^{-2}$.
The identity theorem gives equality to the Gamma moment-generating
function throughout the strip, including the imaginary axis.
Uniqueness of characteristic functions identifies $P$ on all of
$\mathbb R$.  The converse was already proved by Rodrigues's formula.
\end{proof}

\subsection{Complete spectral data, kernel formulas, and their inverse category}

\begin{theorem}[Exact trace domains and reconstruction of the unitary class]
\label{final:O4}
For the operator $A$ above, the ordinary traces are
\begin{align}
 Z_A(s)&=\operatorname{Tr}(\mathrm{Id}+A)^{-s}=\zeta(s),
 &&\operatorname{Re}s>1,\label{final:O4-zeta}\\
 \Theta_A(\tau)&=\operatorname{Tr}e^{-\tau A}
       =\frac1{1-e^{-\tau}},
 &&\operatorname{Re}\tau>0.\label{final:O4-heat-trace}
\end{align}
These are precisely their trace-class domains.  The powers use the
real logarithms of the positive eigenvalues $n+1$; the zero eigenspace
is included through this shift.

Let $B\geq0$ be self-adjoint with compact resolvent.  Each of the
following complete data sets determines its eigenvalue multiset,
including multiplicities, and hence its unitary equivalence class:
all finite heat traces $\operatorname{Tr}e^{-\tau B}$, $\tau>0$;
all shifted zeta traces on a real half-line $s>s_0$ on which they are
finite; or all integer shifted traces
$\operatorname{Tr}(\mathrm{Id}+B)^{-k}$, $k\geq2$, with
$\operatorname{Tr}(\mathrm{Id}+B)^{-2}<\infty$.
The coordinate law of an unspecified realization is not part of this
unitary-class conclusion.
\end{theorem}

\begin{proof}
The basis of Theorem~\ref{final:O3-spectrum} gives diagonal entries
$(n+1)^{-s}$ and $e^{-\tau n}$.  Their singular values are respectively
$(n+1)^{-\operatorname{Re}s}$ and $e^{-n\operatorname{Re}\tau}$.
Their sums converge exactly on the stated half-planes and there give
the Dirichlet and geometric series in
\eqref{final:O4-zeta}--\eqref{final:O4-heat-trace}.
Analytic or meromorphic continuations outside these domains do not
become ordinary operator traces.

For the inverse, compact resolvent gives an orthonormal eigenbasis,
finite multiplicities, and eigenvalues tending to infinity if there
are infinitely many.  For any nonempty remaining multiset let its
least eigenvalue be $\lambda_0$, of multiplicity $m_0$.
If $T(\tau)$ is its remaining heat trace, then
\[
 \lambda_0=-\lim_{\tau\to\infty}\frac{\log T(\tau)}{\tau},
 \qquad m_0=\lim_{\tau\to\infty}e^{\tau\lambda_0}T(\tau).
\]
Indeed, after this rescaling the least-eigenvalue terms equal $m_0$,
and the other terms tend to zero.  Finiteness at a fixed positive
time supplies a summable dominating sequence for all larger times.
Subtract $m_0e^{-\tau\lambda_0}$ and iterate.  The identically zero
remainder detects the end of a finite spectrum, including the empty
Hilbert space.
For zeta data write $r_j=1+\lambda_j$.  The analogous formulas are
\[
 r_0=\lim_{s\to\infty}Z(s)^{-1/s},\qquad
 m_0=\lim_{s\to\infty}r_0^sZ(s).
\]
Finiteness at one exponent again dominates the rescaled tails.
Subtracting $m_0r_0^{-s}$ and repeating recovers every eigenvalue.

For integer data set $x_j=(1+\lambda_j)^{-1}$ and form the finite
positive measure
\[
 \eta_B=\sum_j x_j^2\delta_{x_j}\quad\hbox{on }[0,1],
 \qquad \int x^n\eta_B(dx)=Z_B(n+2),\quad n\geq0.
\]
Two such measures with the same moments agree on polynomials and
then on all continuous functions by uniform polynomial approximation.
They are therefore equal.  Every nonzero atom at $x$ has mass
$m x^2$, recovering the integer multiplicity $m$ and eigenvalue
$x^{-1}-1$.  There is no atom at zero from any finite eigenvalue.
This is the compact Hausdorff moment argument; it uses no interpolation
from the integer values.  Finally an eigenvalue-preserving bijection
of orthonormal eigenbases is an intertwining unitary, with the
weighted coefficient domains carried to one another.
\end{proof}

\begin{proposition}[Compact resolvent does not supply finite trace data]
\label{final:O4-convergence-boundary}
There is a nonnegative diagonal self-adjoint operator with compact
resolvent for which neither a heat trace nor a positive shifted zeta
trace is finite at any positive parameter.
\end{proposition}

\begin{proof}
On $\ell^2(\mathbb N)$ take eigenvalues
$\lambda_j=\log\log(j+e^e)$, with domain
$\{c:\sum\lambda_j^2|c_j|^2<\infty\}$.
They tend to infinity, so $(\mathrm{Id}+B)^{-1}$ is compact.
For every $\tau>0$, $(\log(j+e^e))^{-\tau}$ eventually exceeds
$1/j$, and for every $s>0$,
$(1+\log\log(j+e^e))^{-s}$ eventually exceeds $1/j$.
Both trace series diverge.  Thus the convergence hypotheses in
Theorem~\ref{final:O4} cannot be inferred from compactness alone.
\end{proof}

\begin{theorem}[The complete heat kernel relative to the invariant measure]
\label{final:O4-heat-kernel}
For $\tau>0$, set $r=e^{-\tau}$.  The kernel of $e^{-\tau A}$ relative
to $p(y)\,dy$ is
\begin{align}
 K_\tau(x,y)
 &=\sum_{n=0}^\infty\frac{r^n}{n+1}
       L_n^{(1)}(x)L_n^{(1)}(y)\label{final:O4-kernel-series}\\
 &=\frac{\exp[-r(x+y)/(1-r)]}{(1-r)\sqrt{rxy}}
       I_1\!\left(\frac{2\sqrt{rxy}}{1-r}\right),
       \qquad x,y>0.\label{final:O4-hille-hardy}
\end{align}
Here $I_1$ is the modified Bessel function, specified by
$I_1(z)=\sum_{m\geq0}(z/2)^{2m+1}/(m!(m+1)!)$.
The series converges locally uniformly in $x,y>0$ and in
$L^2(p(x)p(y)\,dx\,dy)$.
The kernel is positive and symmetric,
$\int K_\tau(x,y)p(y)\,dy=1$, and it tends to $1$ as
$\tau\to\infty$.  The Lebesgue transition density is
$K_\tau(x,y)p(y)$.
\end{theorem}

\begin{proof}
Here is a direct verification of the normalization and the closed
formula, which is the parameter-one Hille--Hardy
identity~\cite[\href{https://dlmf.nist.gov/18.18.E27}{Eq.~18.18.27}]{DLMF}.
The finite coefficients in \eqref{final:O3-laguerre} give the
generating function
\[
 G_z(y)=\sum_{n\geq0}L_n^{(1)}(y)z^n
       =(1-z)^{-2}\exp[-yz/(1-z)],\qquad |z|<1.
\]
Let $\widetilde K$ be the right side of
\eqref{final:O4-hille-hardy}.  Expanding $I_1$ gives the nonnegative
series
\[
 \widetilde K(x,y)=e^{-r(x+y)/(1-r)}
   \sum_{m=0}^\infty
   \frac{(rxy)^m}{m!(m+1)!(1-r)^{2m+2}}.
\]
For real $z$ sufficiently near zero, put
$c=(1-rz)/((1-r)(1-z))>0$.
Termwise gamma integration, justified by positivity, gives
\begin{align*}
 \int_0^\infty\widetilde K(x,y)G_z(y)p(y)\,dy
 &=\frac{e^{-rx/(1-r)}}{(1-z)^2(1-r)^2c^2}
       \exp\!\left(\frac{rx}{(1-r)^2c}\right)\\
 &=(1-rz)^{-2}\exp[-rxz/(1-rz)]=G_{rz}(x).
\end{align*}
In particular $z=0$ gives row mass one.  The kernel is symmetric and
positive, so Jensen's inequality and row mass one make its integral
operator a contraction on $\mathcal H$.  For $z$ in a sufficiently
small complex neighborhood of zero, the same exponentially dominated
integral is holomorphic.  Differentiating at zero proves that its
operator maps $L_n^{(1)}$ to $r^nL_n^{(1)}$.  Completeness of the
basis identifies it with $e^{-\tau A}$.
The spectral series is its $L^2$ kernel, since
$\sum r^{2n}<\infty$.  Cauchy's coefficient estimate for $G_z$,
on any circle $|z|=R$ with $\sqrt r<R<1$, bounds its summands
uniformly on compact sets by a convergent geometric sequence.
Thus this series is continuous and locally uniformly convergent.
Its $L^2$ equality with the continuous closed formula implies their
pointwise equality on $(0,\infty)^2$.  Finally the $m=0$ Bessel term
and the convergent series give $K_\tau(x,y)\to1$ as $r\downarrow0$.
The definition of an integral kernel relative to $p(y)\,dy$ proves
the assertion about the Lebesgue density.
\end{proof}

\begin{proposition}[Resolvents, their kernels, and the operator inverse]
\label{final:O4-resolvents}
For $\alpha>0$, the kernels relative to $p(y)\,dy$ of
$(A+\alpha\mathrm{Id})^{-1}$ and $(A+\alpha\mathrm{Id})^{-2}$ are
respectively
\begin{align*}
 R_\alpha(x,y)&=\int_0^\infty e^{-\alpha\tau}K_\tau(x,y)\,d\tau,\\
 S_\alpha(x,y)&=\int_0^\infty \tau e^{-\alpha\tau}K_\tau(x,y)\,d\tau.
\end{align*}
These identities hold as $L^2$ kernels, and the nonnegative integrals
are finite almost everywhere.  The first resolvent is Hilbert--Schmidt
but is not trace class.  The second is trace class and
\[
 \operatorname{Tr}(A+\alpha\mathrm{Id})^{-2}=\zeta(2,\alpha),
\]
where $\zeta(s,\alpha)=\sum_{n=0}^\infty(n+\alpha)^{-s}$ for
$\operatorname{Re}s>1$.  For $\alpha,\beta>0$,
\begin{equation}\label{final:O4-resolvent-difference}
 \operatorname{Tr}\bigl[(A+\alpha\mathrm{Id})^{-1}
       -(A+\beta\mathrm{Id})^{-1}\bigr]
       =\psi(\beta)-\psi(\alpha),
\end{equation}
with $\psi=\Gamma'/\Gamma$.
A complete resolvent operator, with its Hilbert realization and shift
retained, determines $A$ and its domain.  A full marked heat operator
at any fixed positive time also determines $A$.
\end{proposition}

\begin{proof}
On each eigenvector the Laplace integrals are
$\int_0^\infty e^{-(n+\alpha)\tau}\,d\tau=(n+\alpha)^{-1}$ and
$\int_0^\infty\tau e^{-(n+\alpha)\tau}\,d\tau=(n+\alpha)^{-2}$.
The corresponding strong operator integrals converge, and their
kernel integrals converge in $L^2$ because
\[
 \|K_\tau\|_{L^2(p\otimes p)}
       =(1-e^{-2\tau})^{-1/2},
\]
which is $O(\tau^{-1/2})$ near zero and bounded at infinity.
The exponential weights are therefore integrable in this norm.
Positivity and Tonelli identify these $L^2$ limits with the
nonnegative pointwise integrals almost everywhere.
The singular-value series for the resolvents are
$\sum(n+\alpha)^{-1}=\infty$ and
$\sum(n+\alpha)^{-2}<\infty$, proving the ideal statements.
The difference series is absolutely summable and equals
\[
 \sum_{n=0}^\infty
   \left(\frac1{n+\alpha}-\frac1{n+\beta}\right)
 =\psi(\beta)-\psi(\alpha),
\]
as follows by subtracting the convergent digamma expansions
$\psi(a)=-\gamma_E+\sum_{n\geq0}((n+1)^{-1}-(n+a)^{-1})$.
This also fixes the sign in \eqref{final:O4-resolvent-difference}.

For $R=(A+\alpha\mathrm{Id})^{-1}$ one has
$D(A)=\operatorname{Ran}R$ and
$A(Rh)=h-\alpha Rh$ for every $h\in\mathcal H$.
Thus the complete operator recovers the generator, whereas a scalar
resolvent trace discards its eigenvectors and coordinate realization.
For $T=e^{-\tau A}$ the Borel functional calculus similarly gives
$A=-\tau^{-1}\log T$, with its spectral domain; zero is not an
eigenvalue of $T$, and its possible accumulation there is handled by
this unbounded functional calculus.  Consequently a full kernel with
its reference measure retained recovers the corresponding operator.
\end{proof}

\begin{proposition}[Zeta and Fredholm determinants with their conventions]
\label{final:O4-determinants}
For $\alpha>0$, define the zeta determinant using the meromorphic
continuation of the Hurwitz zeta function.  Then
\begin{equation}\label{final:O4-zeta-determinant}
 \det{}_{\zeta}(A+\alpha\mathrm{Id})
 =\exp[-\partial_s\zeta(0,\alpha)]
 =\frac{\sqrt{2\pi}}{\Gamma(\alpha)}.
\end{equation}
In particular
$\det{}_{\zeta}(\mathrm{Id}+A)=\det{}'_{\zeta}(A)=\sqrt{2\pi}$,
where the prime omits the zero mode.
For every $z\in\mathbb C$ one has the entire-function identities
\begin{align}
 \det\bigl(\mathrm{Id}+z(\mathrm{Id}+A)^{-2}\bigr)
 &=\prod_{k=1}^\infty(1+z/k^2)
   =\frac{\sinh(\pi\sqrt z)}{\pi\sqrt z},
   \label{final:O4-fredholm}\\
 \det{}_2\bigl(\mathrm{Id}+z(\mathrm{Id}+A)^{-1}\bigr)
 &=\prod_{k=1}^\infty(1+z/k)e^{-z/k}
   =\frac{e^{-\gamma_E z}}{\Gamma(1+z)}.
   \label{final:O4-det2}
\end{align}
The first determinant is the ordinary trace-class Fredholm determinant;
the second is the Hilbert--Schmidt regularized determinant.
The square-root quotient denotes its entire power series, with
value $1$ at zero.  Its value is independent of the square-root choice.
The full zero multiset of either displayed determinant identifies
the eigenvalue multiset of $A$ in the corresponding nonnegative,
shifted determinant category.
\end{proposition}

\begin{proof}
The spectral series for $A+\alpha\mathrm{Id}$ is the Hurwitz zeta
series on $\operatorname{Re}s>1$.  Its continuation is regular at
zero and obeys
$\partial_s\zeta(0,\alpha)=\log\Gamma(\alpha)-\tfrac12\log(2\pi)$;
see \cite[\href{https://dlmf.nist.gov/25.11.E18}{Eq.~25.11.18}]{DLMF}.
This gives \eqref{final:O4-zeta-determinant}.  The positive eigenvalues
of $A$ and all eigenvalues of $\mathrm{Id}+A$ are both $1,2,\ldots$,
which proves the two special determinant values with the stated
zero-mode convention.
The eigenvalues of $(\mathrm{Id}+A)^{-2}$ are $k^{-2}$, a summable
sequence; those of $(\mathrm{Id}+A)^{-1}$ are $k^{-1}$, whose squares
are summable.  The determinant definitions therefore give the two
products, locally uniformly in $z$.  The identities
\cite[\href{https://dlmf.nist.gov/4.36.E1}{Eqs.~4.36.1 and
5.8.2}]{DLMF} evaluate them as shown.
Evenness of $\sinh w/w$ proves the branch statement.
For a nonnegative operator $B$ with the respective Schatten-class
shifted inverse, a nonzero eigenvalue $\kappa$ of that inverse
produces a zero at $-1/\kappa$ with its multiplicity.
The maps $\lambda\mapsto(1+\lambda)^{-2}$ and
$\lambda\mapsto(1+\lambda)^{-1}$ are injective on $[0,\infty)$.
The shift, exponent, and regularization convention therefore give
the claimed inverse from the entire zero multiset.
\end{proof}

\begin{theorem}[Finite familiar spectral invariants do not determine the spectrum]
\label{final:O4-finite-countermodels}
Fix any finite list of values of the continued shifted zeta function
at nonzero real regular points, optionally its finite part at $s=1$
and its zeta determinant.  There is a nonconstant local family of
nonnegative self-adjoint diagonal operators, differing from $A$ in
only finitely many positive eigenvalues, that retains all these
specified invariants.  Their zeta functions also retain every pole
and principal part, and the value at zero.  Thus this finite list,
including a single zeta determinant as a special case, does not
identify the complete spectrum.
\end{theorem}

\begin{proof}
Vary $N$ distinct shifted eigenvalues $r_j=1+\lambda_j>1$ and fix
the others.  The correction
\[
 \Delta Z(s)=\sum_{j=1}^N(r_j^{-s}-r_{j,0}^{-s})
\]
is entire, leaves every principal part unchanged, and has
$\Delta Z(0)=0$.  Preserving a value at a nonzero real $s_i$ imposes
the equation $\sum r_j^{-s_i}=\sum r_{j,0}^{-s_i}$.
Preserving the finite part at one imposes this same equation with
$s_i=1$.  Preserving the determinant imposes
$\sum\log r_j=\sum\log r_{j,0}$.
After duplicate equations are removed, let their number be $m$.
The gradient rows are nonzero multiples of $r_j^{-s_i-1}$,
and, if present, the additional row $r_j^{-1}$.
Their exponents are distinct.

These rows have rank $m$ at distinct positive $r_j$ whenever $N\geq m$.
Indeed, put $u_j=\log r_j$.  A nonzero linear combination of $m$
exponentials with distinct real exponents has at most $m-1$ distinct
real zeros.  Inductively, divide by its smallest exponential and
differentiate; Rolle's theorem would otherwise give at least $m-1$
zeros for a nonzero combination of $m-1$ exponentials.  The
single-exponential case starts the induction.  Thus a nonzero row
combination cannot vanish at $m$ distinct $u_j$.
Choose $N>m$.  The implicit-function theorem gives a local level
manifold of dimension $N-m>0$ through the original shifted
eigenvalues.  In small disjoint neighborhoods they remain greater
than one and distinct; changes in this manifold cannot merely
permute the spectrum.  The fixed infinite tail preserves compact
resolvent and every trace convergence domain.

For an explicit simultaneous example replace $2,3,4$ by $a,b,c$,
where
\[
 a=2+\varepsilon,\qquad b+c=7-\varepsilon,
 \qquad bc=\frac{24}{2+\varepsilon}.
\]
For sufficiently small nonzero $\varepsilon$, the quadratic roots
$b,c$ are real, distinct, greater than one, and close to $3,4$.
The sum remains $9$ and the product remains $24$.
Consequently $Z(0)$, $Z(-1)$, $Z'(0)$, the residue at one and the
zeta determinant are unchanged, while the full spectrum changes.
All eigenvalues outside this triple, including the zero eigenvalue,
remain fixed.  A determinant alone already permits the two-value
perturbation $2\mapsto2c$, $3\mapsto3/c$ for $c$ near $1$.
These countermodels concern the specified familiar invariants;
they make no claim about an arbitrarily encoded finite datum.
\end{proof}

\subsection{The fixed mixing recipe and its exact unknown-measure inverse}

\begin{theorem}[Existence and uniqueness of the linked mixing measure]
\label{final:O5}
Let $B\geq0$ be a self-adjoint operator on a Hilbert space.  The fixed
function $J(\lambda)=(1+\lambda)^{-2}$ always has the strong-operator
mixing representation
\begin{equation}\label{final:O5-J}
 \int_0^\infty e^{-sB}\,s e^{-s}\,ds
       =(\mathrm{Id}+B)^{-2}.
\end{equation}
If $B$ has a nonzero eigenvector for every $n\in\mathbb N_0$, then
the unknown finite complex Borel measure $\nu$ on $[0,\infty)$ in
\begin{equation}\label{final:O5-unknown}
 \int_{[0,\infty)}e^{-sB}\,\nu(ds)=(\mathrm{Id}+B)^{-2}
\end{equation}
is uniquely $\nu(ds)=s e^{-s}\,ds$.
Finite total variation is understood for a complex measure.  A single
eigenvector at every integer suffices; multiplicities and additional
spectrum do not affect the conclusion.

There is also a full-line formulation without initially defining
negative-time operators.  If $\nu$ is a finite positive Borel measure
on $\mathbb R$ and its finite scalar integer integrals satisfy
\begin{equation}\label{final:O5-integer-samples}
 \int_{\mathbb R}e^{-ns}\,\nu(ds)=(n+1)^{-2}
 \qquad(n\in\mathbb N_0),
\end{equation}
then $\nu(ds)=\mathbf1_{(0,\infty)}s e^{-s}\,ds$.
In this formulation mass, positive support and absence of an atom
at zero are conclusions of the supplied samples.
\end{theorem}

\begin{proof}
The scalar gamma integral gives, for every $\lambda\geq0$,
\[
 \int_0^\infty se^{-(1+\lambda)s}\,ds=(1+\lambda)^{-2}.
\]
The semigroup is a contraction, so each vector-valued integral is
absolutely integrable.  Integration against the spectral measure,
justified by bounded domination, proves
\eqref{final:O5-J} for every nonnegative self-adjoint $B$.

For uniqueness, apply \eqref{final:O5-unknown} to a nonzero
$n$-eigenvector.  This gives all scalar equations
\eqref{final:O5-integer-samples} on $[0,\infty)$.
Under $y=e^{-s}$, their left sides are the moments of the finite
complex pushforward measure on $(0,1]$, extended by zero to $[0,1]$.
The finite positive measure with density $-\log y$ on $(0,1)$ has
these same moments, since
\[
 \int_0^1y^n(-\log y)\,dy=(n+1)^{-2}.
\]
Uniform polynomial approximation on $[0,1]$ identifies the two
finite measures.  Undoing the measurable bijection $s=-\log y$
gives $\nu(ds)=se^{-s}\,ds$.  In particular the zeroth equation
recovers its mass; its lack of an atom at $s=0$ follows as well.

For a positive measure initially on all of $\mathbb R$, mass on
$(-\infty,-\varepsilon]$ would imply
\[
 (n+1)^{-2}\geq e^{n\varepsilon}
       \nu(( -\infty,-\varepsilon])
\]
for every $n$, which is impossible for positive mass.  Taking the
union over $\varepsilon=1/k$ proves support in $[0,\infty)$.
The preceding compact moment proof then applies.  No integral of
unbounded negative-time operators was assumed in this support step.
\end{proof}

\begin{theorem}[A Bernstein function is determined by every integer tail]
\label{final:O6}
Let $f,g$ be Bernstein functions, in their usual representation
category
\begin{equation}\label{final:O6-bernstein-representation}
 f(\lambda)=k_f+d_f\lambda+
   \int_{(0,\infty)}(1-e^{-\lambda s})\,\Pi_f(ds),
 \qquad \lambda\geq0,
\end{equation}
where $k_f,d_f\geq0$, $\Pi_f$ is positive, and
$\int(1\wedge s)\Pi_f(ds)<\infty$, and likewise for $g$.
For any fixed $N\in\mathbb N_0$, equality $f(n)=g(n)$ for every
integer $n\geq N$ implies $f=g$ on $[0,\infty)$, including equality
of killing, drift and L\'evy measure.
In particular the samples
\[
 f(n)=2\log(1+n)\qquad(n\geq N)
\]
identify the complete exponent $f(\lambda)=2\log(1+\lambda)$,
with $k_f=d_f=0$ and $\Pi_f(ds)=2e^{-s}ds/s$.
Knowledge of $f(A)$ on the marked integer eigenspaces supplies these
samples.  This inverse does not extend to arbitrary smooth functions.
\end{theorem}

The integer-tail uniqueness principle for Bernstein functions also appears
in Aguech and Jedidi, Corollary~4~\cite{AguechJedidi}.  The proof below is
independent and additionally records the recovery of the
L\'evy--Khintchine data in the present notation.

\begin{proof}
The finite positive measure
\[
 \rho_f=d_f\delta_1+
    (s\mapsto e^{-s})_\#[(1-e^{-s})\Pi_f(ds)]
 \quad\hbox{on }[0,1]
\]
has no atom at zero.  Its finiteness follows from
$1-e^{-s}\leq1\wedge s$.  Direct subtraction of
\eqref{final:O6-bernstein-representation} gives
\[
 f(n+1)-f(n)=\int_0^1y^n\rho_f(dy),\qquad n\geq0.
\]
Equality of the integer tails therefore identifies all moments of
$y^N\rho_f$ and $y^N\rho_g$.  Compact moment uniqueness gives
$y^N\rho_f=y^N\rho_g$.  Since the weight is strictly positive on
$(0,1]$, division on $[1/j,1]$ and passage to their increasing union
give equality there.  The absence of atoms at zero gives
$\rho_f=\rho_g$ on $[0,1]$.
The atom at one recovers the drift $d_f=d_g$.  The restriction to
$(0,1)$, the inverse coordinate $s=-\log y$, and division by the
strictly positive weight $1-e^{-s}$ recover
$\Pi_f=\Pi_g$.  With these fixed, a single common value at $N$
recovers $k_f=k_g$; this also works when $N>0$.

For the displayed target measure, the L\'evy integrability condition
holds at both endpoints.  Its exponent is zero at $\lambda=0$ and
has derivative
\[
 \frac{d}{d\lambda}\int_0^\infty
       (1-e^{-\lambda s})\frac{2e^{-s}}s\,ds
 =2\int_0^\infty e^{-(1+\lambda)s}\,ds
 =\frac2{1+\lambda}.
\]
Thus it equals $2\log(1+\lambda)$, proving the target assertion.
Functional calculus gives
$f(A)e_n=f(n)e_n$ and
\[
 D(f(A))=\left\{\sum c_ne_n:\sum|f(n)|^2|c_n|^2<\infty\right\}.
\]
Consequently equality on the marked integer eigenspaces has exactly
the required scalar meaning.  For arbitrary smooth functions the
nonzero perturbation $\varepsilon\sin^2(\pi\lambda)$ is invisible
at every nonnegative integer; adding it to the nonnegative target
also preserves nonnegativity.  It therefore gives a distinct scalar
function with the same functional calculus on $A$, outside the
Bernstein inverse category just proved.
\end{proof}

\begin{proposition}[Two different functional-calculus inverse questions]
\label{final:O6-functional-calculus}
For a fixed strictly increasing Bernstein function $f$, the full
self-adjoint operator $f(B)$ determines a nonnegative self-adjoint
$B$.  For the canonical $A$, the complete operator
$2\log(\mathrm{Id}+A)$ is equivalently determined by
$J(A)=(\mathrm{Id}+A)^{-2}$.
These statements retain the known scalar function $f$ or $J$.
\end{proposition}

\begin{proof}
Strict monotonicity and continuity give a Borel inverse on the range
of $f$.  The Borel functional calculus gives
$B=f^{-1}(f(B))$, with its exact spectral domain.  If an endpoint of
the range is a spectral limit, the spectral projection at that
unattained endpoint is zero, so the unbounded inverse is still
well-defined.  In the indicated special case, $J(A)$ has eigenvalues
$(1+n)^{-2}>0$, and
$-\log J(A)=2\log(\mathrm{Id}+A)$ on its functional-calculus domain.
Inverting a known function of an unknown operator is thus different
from Theorem~\ref{final:O6}, which identifies an unknown function
from its values on a known integer spectrum.
\end{proof}

\begin{theorem}[The Gamma shapes two and three retain all spectral and mixing data]
\label{final:O5-isospectral-boundary}
For $\alpha\in\{2,3\}$ set
\[
 p_\alpha(t)=\frac{t^{\alpha-1}e^{-t}}{\Gamma(\alpha)},
 \qquad \mathcal H_\alpha=L^2((0,\infty),p_\alpha(t)\,dt),
\]
and let $A_\alpha$ be the closure of
\[
 -tf''-(\alpha-t)f',\qquad D=C_c^\infty(0,\infty).
\]
Both endpoints are limit point in both cases.  The operators have
the same simple complete spectrum $\mathbb N_0$ and are intertwined
by a unitary carrying $1$ to $1$.  Therefore all their heat traces,
shifted zeta functions, resolvent trace functions and spectral
determinants agree, with the same conventions and domains.
Both satisfy \eqref{final:O5-unknown} with the same unique mixing
measure $\nu_2(ds)=se^{-s}\,ds$.  Their invariant coordinate laws
nevertheless differ: they are $p_2\,dt$ and
$p_3\,dt$, with means $2$ and $3$ respectively.
\end{theorem}

\begin{proof}
The divergence-form coefficients are $q_\alpha=tp_\alpha$.
Apart from an irrelevant nonzero constant, the second zero-energy
solution is $\int_1^t e^s s^{-\alpha}\,ds$.
At zero it is asymptotic to
$-t^{1-\alpha}/(\alpha-1)$; its weighted square is asymptotic to a
positive multiple of $t^{1-\alpha}$, nonintegrable for both
$\alpha=2$ and $3$.  At infinity the solution is asymptotic to
$e^t/t^\alpha$, whose weighted square is a positive multiple of
$e^t/t^{\alpha+1}$, also nonintegrable.
The constant solution is square integrable at both ends.
The maximal-domain and endpoint arguments of Theorem~\ref{final:O3}
therefore apply to each, giving their unique self-adjoint closures.
The form argument also applies, since the lower-cutoff energy is
$O(\varepsilon^{\alpha-1})$, proving conservativity and invariance
of $p_\alpha\,dt$.

The Rodrigues polynomials
\[
 L_n^{(\alpha-1)}(t)=
   \frac{t^{1-\alpha}e^t}{n!}
      \frac{d^n}{dt^n}(e^{-t}t^{n+\alpha-1})
\]
satisfy $A_\alpha L_n^{(\alpha-1)}=nL_n^{(\alpha-1)}$ because their
coefficients obey
$(k+1)(k+\alpha)a_{k+1}+(n-k)a_k=0$.
Integration by parts gives squared norm
$\Gamma(n+\alpha)/(n!\Gamma(\alpha))$ in $\mathcal H_\alpha$.
The weighted Laplace-transform proof of completeness applies without
change: $p_\alpha$ has the same exponential tail and all polynomial
moments.  Hence
\[
 e_{n,\alpha}=
 \sqrt{\frac{n!\Gamma(\alpha)}{\Gamma(n+\alpha)}}
       L_n^{(\alpha-1)}
\]
is a complete orthonormal eigenbasis, with $e_{0,\alpha}=1$.
The unitary $Ue_{n,2}=e_{n,3}$ intertwines the two operators and
fixes the distinguished constant vectors.  Every listed scalar
spectral invariant therefore agrees.  Theorem~\ref{final:O5}
gives the common mixing identity and its mixing-measure uniqueness.
Finally gamma integration gives
$\int t p_\alpha(t)\,dt=\alpha$, proving that the coordinate
probabilities differ despite all those agreements.
\end{proof}

\begin{corollary}[Canonical mixing construction is not unmarked coordinate recovery]
\label{final:O5-canonical-versus-identification}
The fixed exponent-two recipe $J$ constructs the Gamma$(2,1)$ mixing
measure for every nonnegative self-adjoint operator.  Integer
eigenvalues make that mixing measure uniquely recoverable from the
restricted operator identity.  They do not identify the stationary
coordinate law of an arbitrary realization, even when the constant
vector is retained.  To identify that candidate law one must retain
either its explicit equality to the unknown mixing measure in
\eqref{final:O5-unknown}, the local Pearson realization of
Theorem~\ref{final:O2}, or equivalent coordinate information.
\end{corollary}

\begin{proof}
Existence and mixing uniqueness are Theorem~\ref{final:O5}; the two
realizations in Theorem~\ref{final:O5-isospectral-boundary} are a
counterexample to the unmarked coordinate inverse, even under the
conjunction of all the stated spectral and mixing data.
For contrast, a fully marked package
$(\mathcal H,A,1,M_t)$ retains the multiplication operator
$M_tf(t)=tf(t)$ on $\{f:tf\in\mathcal H\}$, and recovers its
coordinate probability by
\[
 P(E)=\langle\mathbf1_E(M_t)1,1\rangle
\]
for every Borel $E$.  A unitary intertwining this multiplication
operator as well as fixing $1$ would preserve that probability,
so the preceding unitary cannot do so.
Furthermore, the different retained recipe
$J_a(\lambda)=(1+\lambda)^{-a}$, $a>0$, has mixing measure
$s^{a-1}e^{-s}\,ds/\Gamma(a)$ by the same gamma integral.
Spectral data do not choose this exponent.

In particular let $F$ forget a stationary coordinate realization
and retain its spectrum, and let $G_J$ construct the canonical
Gamma$(2,1)$ realization using the fixed recipe.  Then
\[
 G_J(F(p_3,A_3))=(p_2,A_2)\ne(p_3,A_3).
\]
Thus this construction is not a left inverse to forgetting on the
class of stationary coordinate realizations.  An identification
graph can use the linked unknown-measure theorem or retain the
fixed coordinate template in its nodes; a cycle through an
unmarked spectrum and canonical selection alone does not prove
identification of its original candidate.
\end{proof}

\section{Complete characteristic series and topology}
\label{final:series}

The complete characteristic series below recover the intrinsic
decomposition of $\Sigma$. The coordinate in every formal statement
is marked. Unless
otherwise specified, coefficients lie in a field $K$ of characteristic
zero.  The same identities hold in a commutative $\mathbb Q$-algebra
when every scalar used as a denominator is a unit.  For a formal series,
$[u^n]$ denotes coefficient extraction and $\operatorname{Res}_u$ is
the coefficient of $u^{-1}$ in a Laurent differential.  A quotient
such as $u/(1-e^{-u})$ means the inverse of the unit
$(1-e^{-u})/u$; no inverse of the nonunit $u$ is being asserted.

\begin{theorem}[The normalized Todd tower and all its twists]
\label{final:F1}
There is a unique series $Q(u)=\sum_{j\geq0}q_ju^j$ satisfying
\[
 q_0=1,\qquad [u^n]Q(u)^{n+1}=1\quad(n\geq1).
\]
It is
\[
 Q_T(u)=\frac{u}{1-e^{-u}}.
\]
Equivalently, one may impose the coefficient equations for every
$n\geq0$, since their degree-zero equation is exactly $q_0=1$.
For this series, for every $n\geq0$,
\begin{equation}\label{final:todd-twists}
 [u^n]e^{ku}Q_T(u)^{n+1}
   =\binom{n+k}{n}
   =\frac{(k+1)(k+2)\cdots(k+n)}{n!}.
\end{equation}
This is a polynomial identity in $k$ over $\mathbb Q$, including the
empty-product convention for $n=0$, and hence holds at every $k\in K$.
In particular it holds for all positive and negative integer twists.
Every individual positive-degree equation is indispensable within the
displayed tower with $q_0=1$.
\end{theorem}

\begin{proof}
For $n\geq1$, expanding the product of $n+1$ copies of $Q$ gives
\[
 [u^n]Q^{n+1}=(n+1)q_n+P_n(q_1,\ldots,q_{n-1}).
\]
The only occurrences of $q_n$ choose degree $n$ from one factor and
degree zero from the others.  As $n+1$ is invertible, this equation
determines $q_n$ from the preceding coefficients.  This proves
uniqueness, provided existence is established.

Put $z=1-e^{-u}$, whose compositional inverse is $u=-\log(1-z)$.
The formal change-of-coordinate rule for residues gives
\begin{align*}
 [u^n]e^{ku}Q_T^{n+1}
 &=\operatorname{Res}_u\frac{e^{ku}}{(1-e^{-u})^{n+1}}\,du\\
 &=\operatorname{Res}_z z^{-n-1}(1-z)^{-k-1}\,dz
   =[z^n](1-z)^{-k-1}
   =\binom{n+k}{n}.
\end{align*}
For completeness, residue substitution by a coordinate
$u=\phi(z)=z+O(z^2)$ follows by linearity from Laurent monomials.
If $j\ne-1$, $\phi^j\phi'$ is the derivative of
$\phi^{j+1}/(j+1)$ and has zero residue; if $j=-1$, write
$\phi=zv$, $v(0)=1$, to obtain
$\phi'/\phi=z^{-1}+v'/v$, with residue one.  Only finitely many
negative Laurent powers are involved, and nonnegative powers cannot
contribute a residue.  The final coefficient identity follows from
the binomial series, or by solving
$(1-z)F'=(k+1)F$, $F(0)=1$, coefficient by coefficient.
Taking $k=0$ proves existence and all tower equations.

If the equation at a positive degree $N$ is removed, retain the
uniquely determined coefficients below $N$, choose any $q_N$ different
from its Todd value, and apply the same recursion at every $n>N$.
The resulting series satisfies every retained equation and differs
from $Q_T$.  A finite initial tower leaves every later coefficient
unconstrained.  Finally, if $q_0$ is only required to be a unit and the
degree-zero equation is absent, the coefficient of $q_n$ becomes
$(n+1)q_0^n$.  Thus each chosen unit $q_0$ gives a recursively determined
tower; normalization cannot be silently dropped.
\end{proof}

For comparison, arbitrary powers have the separate finite coefficient
formula
\begin{equation}\label{final:todd-arbitrary-power}
 [u^n]Q_T(u)^k=
 \sum_{\substack{m_1,\ldots,m_n\geq0\\\sum j m_j=n}}
 (k)_{\underline M}\prod_{j=1}^n\frac{q_j^{m_j}}{m_j!},
 \qquad M=\sum_{j=1}^n m_j,
\end{equation}
where $(k)_{\underline M}=k(k-1)\cdots(k-M+1)$ and the
empty product is one.  Indeed, apply the formal binomial expansion
to $Q_T=1+(Q_T-1)$ and then the multinomial theorem.  This formula
does not replace the exponent $n+1$ and exponential twist in
\eqref{final:todd-twists}.

\begin{theorem}[Reversible characteristic-series formulas]
\label{final:F2}
Define the normalized series
\[
 Q_A(u)=\frac{u/2}{\sinh(u/2)},\qquad
 Q_L(u)=\frac{u}{\tanh u},\qquad
 Q_y(u)=Q_T((1+y)u)-yu.
\]
Their constant terms are one.  With $y$ fixed and marked and
$1+y\ne0$, each of $Q_T,Q_A,Q_L,Q_y$, and each of their multiplicative
inverses, determines the others in the marked formal coordinate.
Explicit reconstruction formulas are
\begin{gather}
 Q_T=e^{u/2}Q_A,\qquad Q_L(u)=Q_T(2u)-u,
 \qquad Q_T(u)=Q_L(u/2)+u/2,\label{final:TAL}\\
 e^u=\frac{Q_T(u)}{Q_T(u)-u},\qquad
 e^{2u}=\frac{Q_L(u)+u}{Q_L(u)-u},\label{final:exponential-recovery}\\
 Q_T(v)=Q_y\left(\frac{v}{1+y}\right)+\frac{yv}{1+y}.
 \label{final:chi-inverse}
\end{gather}
For a reconstruction from $Q_A$ alone, set
\[
 b=\frac{u}{Q_A(u)},\qquad
 w=\frac{b+\sqrt{b^2+4}}{2},\qquad Q_T=wQ_A,
\]
where the formal square root has constant term $2$.
The specialization $y=-1$ is exactly $Q_{-1}(u)=1+u$ and has no
inverse of the form claimed.  Over a general $\mathbb Q$-algebra the
hypothesis is that $1+y$ is a unit, and over $\mathbb Q[y]$ the
inverse formula takes place after localization at $1+y$.
\end{theorem}

\begin{proof}
Substitution of $\sinh v=(e^v-e^{-v})/2$ and
$\tanh v=(e^{2v}-1)/(e^{2v}+1)$ proves
\eqref{final:TAL}--\eqref{final:exponential-recovery}.  All denominators
in the displayed ratios are units.  The square-root recursion is
unique because its leading coefficient is $2$, so its successive
unknown coefficients have invertible coefficient $4$.  In the present
case $b=2\sinh(u/2)$ and this root is $2\cosh(u/2)$; thus $w=e^{u/2}$.
Equivalently $w(0)=1$ and $w-w^{-1}=b$, whose unit solution is unique.
The affine rescaling defining $Q_y$ gives \eqref{final:chi-inverse}
directly, and $y=-1$ gives $1+u$.  Multiplicative inversion on units
is an involution.  These calculations establish the stated reversals
on the exact series; applying the square-root algorithm to an
arbitrary even unit does not certify that its input equals $Q_A$.
\end{proof}

\begin{proposition}[Formal, analytic, and global recovery]
\label{final:F3}
The preceding reversals recover the formal exponential and its inverse
formal logarithm, hence the germ of $H(1+z)=\log(1+z)-z$ at $z=0$.
After specializing coefficients and the marked parameter $y$ to
the reals, they also construct the displayed standard real analytic functions
on their full real domains.  They identify an independently supplied
global function from its germ only in a declared uniqueness category,
for example real analytic functions on the same connected domain.
No analogous identification of arbitrary smooth global extensions is
valid.
\end{proposition}

\begin{proof}
Equation \eqref{final:exponential-recovery} recovers $e^u$, and the
series $e^u-1$ has linear coefficient one, so it has a unique
compositional inverse, $\log(1+z)$.  The real quotient definitions
extend through $u=0$ with value one.  Their other denominators do not
vanish on the real line; in the $Q_A$ reconstruction the chosen root
is $2\cosh(u/2)>0$.  Consequently the explicit formulas define the
claimed real analytic extensions.  The analytic identity theorem on
a connected interval proves uniqueness within that category.

The nearest nonremovable complex poles are $\pm2\pi i$ for $Q_T,Q_A$
and $\pm\pi i$ for $Q_L$.  Their Taylor discs therefore do not themselves
cover the real line.  For the smooth counterexample, choose a nonzero
smooth bump $\chi$ supported in a compact subinterval of $(0,1)$ away
from $1$.  Then $I+\varepsilon\chi$ has the same germ at $1$ as $I$
and differs globally for $\varepsilon\ne0$.  Since $I''=t^{-2}$ has a
positive minimum on the support, a sufficiently small
$|\varepsilon|$ preserves strict convexity.  A formal identity alone
cannot exclude this example.
\end{proof}

\begin{theorem}[The universal Thom correction, conditional on topology]
\label{final:F4}
Suppose complex topological $K$-theory and rational ordinary cohomology
are supplied with their Thom isomorphisms, compatible complex
orientations, natural multiplicative Thom classes, the Chern character
normalized by $\operatorname{ch}(L)=e^{c_1(L)}$, and an injective
splitting principle.  For a complex vector bundle $V\to B$ on a finite
CW complex, with projection $\pi$, the uniquely defined correction in
\begin{equation}\label{final:Thom-comparison}
 \operatorname{ch}(U_K(V))=\pi^*C(V)\,U_H(V)
\end{equation}
is
\[
 C(V)=\operatorname{Td}(V)^{-1}
      =\prod_i Q_T(u_i)^{-1},
\]
where the $u_i$ are formal Chern roots.  Its complete universal line
series determines $Q_T$ and conversely.  With compatible collapse and
Gysin constructions supplied as well, this yields
\begin{equation}\label{final:RR}
 \operatorname{ch}(f_!^K a)
   =f_*^H\bigl(\operatorname{ch}(a)\operatorname{Td}(T_f)\bigr)
\end{equation}
for $a\in K^0$ and a proper smooth map between manifolds with the
specified complex $K$-orientation and virtual relative tangent bundle
$T_f$, using compact supports when required.  On universal bases
the characteristic-class statement is taken in completed rational
cohomology.
\end{theorem}

\begin{proof}
The cohomological Thom isomorphism makes multiplication by $U_H(V)$
an isomorphism from the cohomology of the base to the appropriate
relative group, so it gives existence and uniqueness of $C(V)$ in
\eqref{final:Thom-comparison}.  For a line bundle put $u=c_1(L)$.
The zero-section pullbacks of the two Thom classes are respectively
$1-[L^*]$ and $u$.  Thus
\[
 uC(u)=1-e^{-u}.
\]
Naturality makes these identities compatible on the finite stages
of the universal line bundle.  Their inverse limit is
$\mathbb Q[[u]]$, in which multiplication by $u$ is injective.
Hence $C(u)=(1-e^{-u})/u=Q_T(u)^{-1}$.  Multiplicativity of Thom
classes and the Chern character gives the product formula after
splitting, and injectivity of the supplied splitting map descends
it to $B$.  For an embedding, applying the collapse map gives the
cohomological pushforward with the correction
$\operatorname{Td}(\nu)^{-1}$ for its complex normal bundle $\nu$.
Since $T_f=-\nu$, this is \eqref{final:RR}; the supplied stable
factorization and composition rules give the same formula for the
stated proper maps.  Finally, the universal line correction is a
unit, so inversion recovers $Q_T$.
\end{proof}

The scalar Euler-coordinate identity used here is itself elementary:
for $x\star y=x+y+xy$ and $z(x)=x/(1+x)$,
\[
 z(x\star y)=z(x)+z(y)-z(x)z(y).
\]
Indeed $1-z(x)=1/(1+x)$ is multiplicative under $\star$.
Interpreting $x$ as $[L]-1$ and $z$ as $1-[L^*]$ is a map into a
supplied $K$-theory ring.  The scalar identity does not construct
bundles, Thom isomorphisms, or the underlying spaces.  On a finite
base $u$ may be nilpotent and a zero divisor; cancellation in that
base would be invalid.  If $u^{N+1}=0$, adding $u^{N+1}R(u)$ to a
characteristic series does not change its evaluation there.  A single
finite-base Thom comparison or numerical index consequently cannot
determine the complete universal series.  Interpretations of the
$L$, $\widehat A$, and $\chi_y$ series retain their respective
oriented, spin, or complex/holomorphic foundations; the algebraic
formulas do not assert those foundations or their index theorems.

\begin{proposition}[Rational characteristic data lose integral data]
\label{final:F5}
On $\mathbb{RP}^2$, the complexification $L$ of the tautological real
line is nontrivial and $[L]\ne[1]$ in complex $K^0$, although
$\operatorname{ch}(L)=1$ in rational cohomology.  Thus even with the
base fixed, rational characteristic data and the complete scalar
series do not reconstruct integral bundle or $K$-theory data.
\end{proposition}

\begin{proof}
The usual cellular calculation gives
$H^2(\mathbb{RP}^2;\mathbb Z)=\mathbb Z/2$ and no positive-degree
rational cohomology.  To verify nontriviality directly, a nowhere-zero
section of $L$ would pull back to an odd map
$g:S^2\to\mathbb C\setminus\{0\}$.  Normalize it to an odd map
$S^2\to S^1$.  Since $S^2$ is simply connected, write
$g(x)=e^{i\theta(x)}$ for a continuous real function $\theta$.
Then $\theta(-x)-\theta(x)$ is a continuous odd integral multiple
of $\pi$, hence constant by connectedness.  Replacing $x$ by $-x$
changes its sign, which is impossible for a nonzero odd multiple
of $\pi$.  Thus $L$ is nontrivial.  Equivalently, under the usual
classification of complex line bundles, its first Chern class is
the nonzero element of the displayed group.

If $[L]=[1]$, the Grothendieck-group relation supplies a bundle $E$
with $L\oplus E\cong 1\oplus E$.  Taking determinant lines and
cancelling $\det E$ would make $L$ trivial, a contradiction.
On the other hand, every positive-degree rational Chern class
vanishes, so $\operatorname{ch}(L)=1$.  The transition functions of
$L$ are signs, so $L\otimes L\cong1$, as expected from the torsion
line class.  These arguments use the stated elementary topological
foundations and do not follow from formal scalar arithmetic alone.
\end{proof}

\section{The matrix lift and its Gaussian statistics}
\label{final:matrix}

Write $\mathcal P_n$ for the cone of real symmetric positive-definite
$n\times n$ matrices, $E_n$ for its identity, and use the trace pairing
$\langle U,V\rangle=\operatorname{tr}(UV)$ on symmetric matrices.
All matrix logarithms, square roots, and powers below are the positive
spectral ones.  The finite-dimensional real spectral theorem is
retained throughout this section.

\begin{theorem}[The unique scalar-block spectral lift]
\label{final:M1}
Suppose a family $F_n:\mathcal P_n\to\mathbb R$, $n\geq1$, satisfies
\[
 F_1([t])=I(t),\qquad F_n(QXQ^{\mathsf T})=F_n(X)
\]
for orthogonal $Q$, and the scalar-block recursion
\[
 F_{n+1}(X\oplus[t])=F_n(X)+I(t).
\]
Then, without any regularity assumption,
\begin{equation}\label{final:matrix-potential}
 F_n(X)=\Phi_n(X):=\sum_{i=1}^n I(\lambda_i(X))
          =\operatorname{tr}X-\log\det X-n.
\end{equation}
This family satisfies full block additivity.  Neither orthogonal
invariance nor scalar-block recursion can be removed while keeping
the scalar seed, even among nonnegative candidates.
\end{theorem}

\begin{proof}
Orthogonally diagonalize $X$ and apply scalar-block recursion until
all blocks have size one.  This forces the sum in
\eqref{final:matrix-potential}.  Trace, determinant, and their behavior
on block diagonal matrices verify every axiom for $\Phi_n$.

If orthogonal invariance is omitted,
$F_n(X)=\sum_i I(X_{ii})$ retains the seed and full block additivity.
Rotating $\operatorname{diag}(2,1)$ through angle $\pi/4$ gives diagonal
entries $3/2,3/2$, and
$I(2)-2I(3/2)=\log(9/8)\ne0$, so it is a different lift.
If scalar-block recursion is omitted, for $\varepsilon>0$ put
\[
 G_n(X)=\Phi_n(X)+\varepsilon
   \sum_{i<j}(\log\lambda_i(X)-\log\lambda_j(X))^2.
\]
This has the seed, orthogonal invariance, and nonnegativity, but on
$[a]\oplus[b]$, $a\ne b$, its extra term is nonzero whereas both
rank-one extra terms vanish.  If the seed is omitted, the zero family
already retains the two structural rules.  These are omission tests
in the stated family, rather than a claim of an absolute weakest
possible axiomatization.
\end{proof}

\begin{theorem}[Derivatives, divergence, duality, and determinant bounds]
\label{final:M2}
In the indicated affine coordinates the potential $\Phi_n$ satisfies
\begin{align}
 \nabla\Phi_n(X)&=E_n-X^{-1},\label{final:matrix-gradient}\\
 g_X(U,V):=D^2\Phi_n(X)[U,V]
   &=\operatorname{tr}(X^{-1}UX^{-1}V),\label{final:matrix-hessian}\\
 D(X,Y):=\Phi_n(X)-\Phi_n(Y)
   -\langle\nabla\Phi_n(Y),X-Y\rangle
   &=\operatorname{tr}(Y^{-1}X)-\log\det(Y^{-1}X)-n.
   \label{final:matrix-divergence}
\end{align}
The Hessian is positive definite and $D(X,Y)\geq0$, with equality
exactly when $X=Y$.  For every invertible real $C$, simultaneous
congruence preserves both $D$ and $g$.  Inversion is an isometry for
$g$ and reverses the directed divergence:
\[
 D(X^{-1},Y^{-1})=D(Y,X).
\]
Extending $\Phi_n$ by $+\infty$ outside $\mathcal P_n$ in the vector
space of symmetric matrices, its conjugate is
\begin{equation}\label{final:matrix-conjugate}
 \Phi_n^*(\Theta)=
 \begin{cases}
  -\log\det(E_n-\Theta),&E_n-\Theta\in\mathcal P_n,\\
  +\infty,&\text{otherwise}.
 \end{cases}
\end{equation}
The Legendre coordinate is $\Theta=E_n-X^{-1}$, with inverse
$X=(E_n-\Theta)^{-1}$.  Moreover
\[
 \det X\leq e^{\operatorname{tr}X-n},\qquad
 \det X\leq\left(\frac{\operatorname{tr}X}{n}\right)^n,
\]
with equality respectively at $E_n$ and at positive scalar matrices;
$\log\det$ is strictly concave.
\end{theorem}

\begin{proof}
Jacobi's formula gives $D\log\det X[U]=\operatorname{tr}(X^{-1}U)$.
Differentiating $XX^{-1}=E_n$ gives
$D(X^{-1})[U]=-X^{-1}UX^{-1}$ and proves the first two formulas.
Their positivity is the precise identity
\[
 g_X(U,U)=\operatorname{tr}
    \bigl((X^{-1/2}UX^{-1/2})^2\bigr)
          =\|X^{-1/2}UX^{-1/2}\|_F^2.
\]
It is the trace of a square, and is strictly positive for $U\ne0$.
Substitution yields \eqref{final:matrix-divergence}.  With
$A=Y^{-1/2}XY^{-1/2}$, trace and determinant show
$D(X,Y)=\sum_i I(\lambda_i(A))$.  Scalar nonnegativity and its equality
case prove the assertion, since $A=E_n$ precisely when $X=Y$.
Cyclic trace and the inverse-congruence formula prove congruence
invariance.  Substituting inverses in $D$ gives its reversal; inserting
the differential $-X^{-1}UX^{-1}$ in $g_{X^{-1}}$ gives $g_X$.

If $B=E_n-\Theta$ is positive definite, maximizing the conjugate is
equivalent to maximizing $-\operatorname{tr}(BX)+\log\det X+n$.
Substitute $A=B^{1/2}XB^{1/2}$ to obtain
$-\Phi_n(A)-\log\det B$.  Its unique maximum occurs at $A=E_n$,
which proves the finite formula and its optimizer.  If $B$ has an
eigenvalue $b\leq0$, take $X$ diagonal in an eigenbasis of $B$, with
the matching eigenvalue equal to $r\to\infty$ and all others one.
The contribution $-br+\log r$ diverges even when $b=0$.
Substituting Fenchel equality consequently gives the Bregman duality
$D_{\Phi_n}(X,Y)=D_{\Phi_n^*}(\Theta_Y,\Theta_X)$.

Finally $\Phi_n(X)\geq0$ gives the first determinant bound.  Apply
it to $nX/\operatorname{tr}X$ for the second and its equality case.
The Hessian of $\log\det$ is $-g$, which proves strict concavity on
the convex cone.
\end{proof}

The gradient in \eqref{final:matrix-gradient}, together with
$F(E_n)=0$, identifies a differentiable $F$ on the cone: the gradient
of $F-\Phi_n$ vanishes, so integration on each straight segment makes
the difference constant.  Likewise the full Hessian
\eqref{final:matrix-hessian}, together with $F(E_n)=0$ and
$\nabla F(E_n)=0$, identifies a twice differentiable $F$, since the
Hessian of the difference vanishes and its gradient is constant on
segments.  These affine marks remove exactly the affine ambiguity.
The rank-one potential or its calibrated anchored Bregman function
recovers $I$ directly.  A metric given merely up to isometry does not
provide these affine coordinates or marks.  Nor does congruence
invariance alone select this metric: the forms
\[
 \alpha g_X(U,V)+\beta\operatorname{tr}(X^{-1}U)
                             \operatorname{tr}(X^{-1}V),
 \quad \alpha>0,\quad\beta>-\alpha/n,
\]
are also invariant positive metrics.  Positivity follows by separating
$X^{-1/2}UX^{-1/2}$ into its scalar and traceless parts.  The selected
Hessian, rather than bare invariance or determinant inequalities,
removes this freedom.

\begin{theorem}[Geodesics, distance, and the symmetrization boundary]
\label{final:M3}
For $X,Y\in\mathcal P_n$ put
$K=\log(X^{-1/2}YX^{-1/2})$.  The unique constant-speed minimizing
geodesic for $g$ is
\[
 \gamma(s)=X^{1/2}e^{sK}X^{1/2},\quad0\leq s\leq1,
 \qquad d(X,Y)=\|K\|_F.
\]
If $\lambda_i$ are the eigenvalues of $X^{-1/2}YX^{-1/2}$, then
\[
 D(X,Y)+D(Y,X)=4\sum_i\sinh^2\left(\frac{\log\lambda_i}{2}\right).
\]
For $n=1$ this is $4\sinh^2(d(X,Y)/2)$.  For every $n\geq2$,
symmetrization is not a function of distance alone.
\end{theorem}

\begin{proof}
Congruence reduces the distance problem to a curve from $E_n$ to
$A=X^{-1/2}YX^{-1/2}$.  Write any piecewise continuously differentiable
such curve as $B(s)=e^{Z(s)}$ with $Z(s)$ symmetric.  The exponential
differential, obtained by differentiating its power series, is
\[
 D\exp_Z[V]=\int_0^1e^{(1-r)Z}Ve^{rZ}\,dr.
\]
In a basis diagonalizing $Z$ with eigenvalues $z_i$, its $(i,j)$ entry
is $V_{ij}(e^{z_i}-e^{z_j})/(z_i-z_j)$, with the continuous value
$e^{z_i}V_{ij}$ when $z_i=z_j$.  Consequently
\[
 g_B(B',B')=
 \sum_{i,j}\left(\frac{2\sinh((z_i-z_j)/2)}{z_i-z_j}\right)^2
                        (Z'_{ij})^2\geq\|Z'\|_F^2,
\]
where the quotient is one at zero.  Its absolute value is at least
one by the power series for $\sinh$.  Integrating speed gives
\[
 \operatorname{length}_g(B)\geq\int_0^1\|Z'(s)\|_F\,ds
                  \geq\|\log A\|_F.
\]
The curve $e^{s\log A}$ attains both inequalities and has constant
speed.  For a minimizing curve equality in the Euclidean length
inequality forces $Z(s)=a(s)\log A$ with $a$ nondecreasing from zero
to one.  On this curve metric speed is $a'(s)\|\log A\|_F$;
constant speed forces $a(s)=s$.  If $A=E_n$, zero length forces the
constant curve.  This proves uniqueness as stated.  The minimizing
curve is a geodesic by the usual first-variation theorem; directly
it satisfies $\gamma''=\gamma'\gamma^{-1}\gamma'$, the Euler--Lagrange
equation obtained from the energy
$\frac12\int\operatorname{tr}(\gamma^{-1}\gamma'
\gamma^{-1}\gamma')\,ds$.

Adding the two divergence formulas cancels logarithms and gives
$\sum_i(\lambda_i+\lambda_i^{-1}-2)$, proving symmetrization.
For $r>0$ compare $X=E_n$ with matrices whose log eigenvalue vectors
are $(r,0,\ldots,0)$ and
$(r/\sqrt2,r/\sqrt2,0,\ldots,0)$.  Both distances equal $r$, whereas
their symmetrizations are $2\cosh r-2$ and
$4\cosh(r/\sqrt2)-4$.  The difference has power series
\[
 \sum_{j\geq2}\frac{2-2^{2-j}}{(2j)!}r^{2j}>0.
\]
This proves the failure for every positive distance in rank at least two.
\end{proof}

\begin{theorem}[The supplied Gaussian and Wishart sampling model]
\label{final:M4}
Let $m\geq1$ and $n\geq1$ be specified integers and suppose
$Z_1,\ldots,Z_m$ are independent, identically distributed
$N(0,X)$ vectors with $X\in\mathcal P_n$ and known mean zero.  Put
$W=\sum_{j=1}^mZ_jZ_j^{\mathsf T}$ and $C=W/m$.  The scatter $W$ has
the Wishart law defined by this model and is a sufficient statistic
for $X$, with
\[
 \log L(X)=-\frac m2\log\det X
              -\frac12\operatorname{tr}(X^{-1}W)+\text{constant}.
\]
For symmetric $\Theta$,
\[
 \mathbb E e^{\operatorname{tr}(\Theta W)}
   =\det(E_n-2X^{1/2}\Theta X^{1/2})^{-m/2}
\]
when the matrix inside the determinant is positive definite; the
expectation is infinite otherwise.  The scatter is positive definite
almost surely exactly when $m\geq n$.  Whenever the observed $C$ is
positive definite, the maximum likelihood covariance is uniquely $C$
and, with $\ell=-\log L$,
\[
 \ell(X)-\ell(C)=\frac m2D(C,X).
\]
If $C$ is singular there is no maximum likelihood estimate in the
open cone $\mathcal P_n$.  Finally
\[
 \operatorname{KL}(N(0,X)\,\|\,N(0,Y))=\frac12D(X,Y),
\]
and the divergence for $m$ independent copies is $m$ times this value.
\end{theorem}

\begin{proof}
Multiplying the Gaussian densities proves the likelihood formula.
Its dependence on the samples factors through $W$, so the
factorization criterion for dominated statistical families proves
sufficiency.  For one sample write $Z=X^{1/2}G$ with $G\sim N(0,E_n)$.
The tilted Gaussian integral has quadratic form
$E_n-2X^{1/2}\Theta X^{1/2}$.  Orthogonal diagonalization reduces it
to one-dimensional integrals, giving its inverse square-root
determinant when all eigenvalues are positive.  A zero or negative
eigenvalue yields an infinite one-dimensional integral, so Tonelli's
theorem gives divergence otherwise.  Independence raises the finite
answer to the $m$th power.

The rank of $W$ is the dimension of the span of its $m$ sample vectors.
It is at most $m$, and the first $n$ Gaussian sample vectors are linearly
independent almost surely: the determinant-zero set has Lebesgue measure
zero, as follows, for example, by conditioning successively on the span
of the preceding vectors. Thus $W$ is positive definite almost surely
for $m\geq n$, and singular for $m<n$.

For positive-definite $C$, subtracting likelihood values gives
$mD(C,X)/2$, with the asserted minimum and equality case. For singular
$C$, diagonalize it and let the eigenvalues of $X$ along the nontrivial
kernel of $C$ tend to zero while keeping the remaining eigenvalues fixed
and positive. The trace term stays fixed and $\log\det X\to-\infty$,
so $\ell(X)\to-\infty$ and no minimizer exists in the open cone.
Finally integrating the Gaussian log
density ratio uses $\mathbb E_X ZZ^{\mathsf T}=X$ and gives $D(X,Y)/2$;
the log ratio adds over independent samples.
\end{proof}

The sample count, independence, Gaussian law, and known-mean convention
are data of this theorem.  The scalar potential alone supplies none
of those sampling choices.  Likewise the spectral or Gaussian
construction of $\Phi_n$ identifies that constructed family; an
independently designated family is identified only after the rules
of Theorem~\ref{final:M1} are imposed.

\section{Gaussian radial realization and independent spatial data}
\label{final:spatial}

\begin{theorem}[The isotropic four-dimensional Gaussian realization]
\label{final:R1}
If $Z\sim N(0,E_4)$ and $T=\|Z\|^2/2$, then $T$ has density
$p(t)=te^{-t}$ on $t>0$.  Conversely, if $T$ has this density and
$\Omega$ is an independent uniform point of $S^3$, then
$Z=\sqrt{2T}\,\Omega$ has law $N(0,E_4)$.  Equivalently the radial
law together with orthogonal invariance uniquely identifies this
Gaussian law; the radial law alone does not.

For the specified Ornstein--Uhlenbeck process in $\mathbb R^D$,
\[
 dZ_s=-\tfrac12Z_s\,ds+dB_s,
\]
the generator on radial-energy functions $F(z)=f(\|z\|^2/2)$ is
\[
 \left(\tfrac12\Delta-\tfrac12z\cdot\nabla\right)F
       =t f''(t)+(D/2-t)f'(t).
\]
Thus its radial generator is $tf''+(2-t)f'$ exactly when $D=4$.
In dimension four the radial energy satisfies
\[
 dT_s=(2-T_s)\,ds+\sqrt{2T_s}\,d\beta_s
\]
for a one-dimensional Brownian motion $\beta$.  These statements
select a dimension within this specified Gaussian construction.
\end{theorem}

\begin{proof}
Polar integration in $\mathbb R^4$, using $|S^3|=2\pi^2$, gives
\[
 (2\pi)^{-2}e^{-r^2/2}\,2\pi^2r^3\,dr
       =te^{-t}\,dt,\qquad t=r^2/2.
\]
This proves the forward law.  Conversely, integration of any bounded
Borel test function first over uniform $\Omega$ and then over this
radial density reproduces the four-dimensional Gaussian integral.
For an orthogonally invariant law, average the test function over
the normalized Haar measure on the compact group $O(4)$.  On each
nonzero sphere this average is its uniform spherical average, so
its integral is fixed by the radial law.  This proves the equivalent
uniqueness statement and the independence of radius and uniform
angle.  Choosing instead a fixed unit vector $v$ and
$Z=\sqrt{2T}\,v$ preserves the radial law and changes the vector law.

The chain rule gives
$\nabla F=f'(t)z$, $\Delta F=2tf''(t)+Df'(t)$, and
$z\cdot\nabla F=2tf'(t)$, proving the generator formula.
It\^o's formula yields
\[
 dT_s=(D/2-T_s)\,ds+Z_s\cdot dB_s,
 \qquad d\langle T\rangle_s=2T_s\,ds.
\]
Normalize the stochastic integrand to $Z_s/\|Z_s\|$ off zero and
choose a fixed unit vector at zero.  Its stochastic integral
$\beta_s$ has quadratic variation $s$, hence is Brownian by the
continuous-martingale characterization, and
$Z_s\cdot dB_s=\sqrt{2T_s}\,d\beta_s$.  At $D=4$ this gives the
displayed equation.  The stationary vector law is $N(0,E_D)$ by the
explicit Gaussian solution of the OU equation; polar integration
therefore gives the stationary energy law
$t^{D/2-1}e^{-t}/\Gamma(D/2)$.  Its shape is two precisely at $D=4$.
\end{proof}

This construction retains the Gaussian and stochastic-calculus
foundations just used.  A stationary density by itself does not
choose a diffusion: for any specified positive differentiable
coefficient $a(t)$ the formal zero-current expression
\[
 L_af=\frac1p(apf')'
     =af''+\left(a'+a\frac{p'}p\right)f'
\]
has the same formal symmetric measure, as integration by parts on
compactly supported tests verifies.  Domain and endpoint conditions
remain necessary when selecting an actual closed generator.  Choosing
$a(t)=t$ gives the radial operator in Theorem~\ref{final:R1}; it is a
selection, not a consequence of knowing only $p$.

\begin{theorem}[Nonnegative orthogonal additivity without regularity]
\label{final:R2}
Let $V$ be a real inner-product space of dimension at least two and
let $r:V\to[0,\infty)$.  Then
\[
 r(x+y)=r(x)+r(y)\quad\text{whenever }\langle x,y\rangle=0
\]
holds if and only if $r(x)=c\|x\|^2$ for some $c\geq0$.  No continuity,
measurability, homogeneity, or radiality hypothesis is required.
Consequently a function $z:V\to[0,\infty)$ satisfies
\[
 z(x+y)=z(x)+z(y)+z(x)z(y)\quad(x\perp y)
\]
if and only if $z(x)=e^{c\|x\|^2}-1$ for some $c\geq0$.
Nontriviality forces $c>0$, while a numerical value at one nonzero
vector is needed to calibrate its scale.
\end{theorem}

\begin{proof}
The equation at $(0,0)$ gives $r(0)=0$.  Define the even and odd parts
\[
 e(x)=\frac{r(x)+r(-x)}2,\qquad
 o(x)=\frac{r(x)-r(-x)}2.
\]
They are orthogonally additive, $e\geq0$, and $|o(x)|\leq e(x)$.
If $\|x\|=\|y\|$, set $a=(x+y)/2$, $b=(x-y)/2$.  Then $a\perp b$,
$x=a+b$, and $y=a-b$.  Evenness gives
$e(x)=e(a)+e(b)=e(a)+e(-b)=e(y)$.  Every nonnegative real number is
the squared norm of a vector, so $e(x)=A(\|x\|^2)$ for a function
$A:[0,\infty)\to[0,\infty)$.

There exist two orthogonal unit vectors.  Multiplying them by
$\sqrt a$ and $\sqrt b$ shows $A(a+b)=A(a)+A(b)$ for arbitrary
$a,b\geq0$.  Nonnegativity makes $A$ monotone: if $b\geq a$ then
$A(b)-A(a)=A(b-a)\geq0$.  With $c=A(1)$, additivity gives
$A(q)=cq$ for every nonnegative rational $q$.  For any real $a\geq0$,
choose rational sequences approaching $a$ from below and above.
Monotonicity sandwiches $A(a)$ between their multiples by $c$, so
$A(a)=ca$.

For $x\ne0$, choose $y\perp x$ of the same norm; the case $x=0$ is
automatic.  Then $x+y\perp x-y$, and orthogonal additivity and oddness
give
\[
 o(2x)=o(x+y)+o(x-y)
      =o(x)+o(y)+o(x)+o(-y)=2o(x).
\]
Iterating downward, $o(x)=2^ko(x/2^k)$, while
$|o(x/2^k)|\leq c\|x\|^2/4^k$.  Hence
$|o(x)|\leq c\|x\|^2/2^k$ for every $k$, forcing $o(x)=0$.
This proves $r=e=c\|x\|^2$.  Pythagoras proves the converse.

For the second statement take $r=\log(1+z)$.  Its values are
nonnegative, and the composition equation is exactly orthogonal
additivity of $r$.  Apply the first statement and exponentiate;
the converse follows by multiplication of exponentials.
\end{proof}

Both substantive hypotheses have direct omission tests.  In dimension
one, $r(x)=x^4$ is nonnegative and orthogonally additive, since every
orthogonal pair includes zero, but is not $c\|x\|^2$.  Without
nonnegativity in dimension at least two, a nonzero linear functional,
or $\|x\|^2$ plus such a functional, is orthogonally additive and
violates the conclusion.  Accordingly the condition $z>-1$ alone
also fails: $z(x)=e^{\ell(x)}-1$ for a nonzero linear functional
obeys the composition rule and takes negative values.

\begin{proposition}[The profile-independent radial residual]
\label{final:R3}
In integer Euclidean dimension $D\geq1$ the radial Laplacian is
$\Delta_D f=f''+(D-1)f'/x$ for $x>0$.  For every twice differentiable
nonvanishing profile $f$ and $\alpha=(D-1)/2$,
\begin{equation}\label{final:radial-residual}
 \frac{(x^\alpha f)''}{x^\alpha f}
    -\left(\frac{f''}{f}+\frac{D-1}{x}\frac{f'}{f}\right)
     =\frac{(D-1)(D-3)}{4x^2}.
\end{equation}
Equivalently, writing $f=x^{-\alpha}u$ gives the operator identity
\[
 -\Delta_D f=x^{-\alpha}
   \left(-u''+\frac{(D-1)(D-3)}{4x^2}u\right),
\]
which does not require $f$ to be nonvanishing.  Residual cancellation
holds exactly for $D=1$ or $D=3$; with the independently retained
restriction $D\geq2$ it selects $D=3$.  It has no reverse identifying
force for the profile $f$.
\end{proposition}

\begin{proof}
Twice applying the product rule gives
\[
 \frac{(x^\alpha f)''}{x^\alpha f}
 =\frac{f''}{f}+\frac{2\alpha}{x}\frac{f'}{f}
                           +\frac{\alpha(\alpha-1)}{x^2}.
\]
Since $2\alpha=D-1$ and
$\alpha(\alpha-1)=(D-1)(D-3)/4$, subtraction proves
\eqref{final:radial-residual}; the same product rule proves the
operator form.  The coefficient vanishes precisely at the two stated
integer dimensions.  It is independent of $f$, so for example
$f(x)=e^{-x^2}$ also gives cancellation in dimension three, despite
being a different profile from $p$ or the scalar potential.
\end{proof}

\begin{proposition}[Spatial countermodels retaining the complete intrinsic structure]
\label{final:R4}
All intrinsic scalar identities, their probability and operator
realizations, their formal characteristic series, and the canonical
four-dimensional Gaussian construction may hold while an independent
designated spatial observable fails orthogonal additivity, or while
its designated dimension fails to equal three.  Their conjunction therefore
does not imply either the hypotheses of Theorem~\ref{final:R2} or a
designated dimension conclusion from Proposition~\ref{final:R3}.
\end{proposition}

\begin{proof}
Keep the entire intrinsic structure and every canonical auxiliary
object fixed.  Adjoin a distinct Euclidean spatial sort $V$ with its
own designated observable $r$.  Three choices give the required
models.  In $V=\mathbb R^2$ take $r(x)=\|x\|^4$.  For perpendicular
unit vectors $x,y$, $r(x+y)=4$ while $r(x)+r(y)=2$, so even positivity,
smoothness, and radiality do not force the composition law.  In
$V=\mathbb R^2$ take $r(x)=\|x\|^2$.  Orthogonal additivity holds,
but the coefficient in \eqref{final:radial-residual} is $-1/(4x^2)$.
Finally in $V=\mathbb R$ take squared norm: residual cancellation
holds although $D\ne3$.  None of these additions changes any object
or equation in the original intrinsic structure.  In particular,
each coexists with the separately constructed Gaussian vector in
$\mathbb R^4$; no identification between its vector space and the
designated spatial sort was supplied.
\end{proof}

There is also a direct tensor distinction.  Pulling a one-dimensional
metric $g(r)\,dr^2$ back by a spatial scalar map produces
$g(r)\,dr\otimes dr$, of rank at most one.  The Hessian of a composite
potential is instead
\[
 D^2(I\circ r)=I''(r)\,dr\otimes dr+I'(r)D^2r.
\]
The chain rule proves this formula.  Its second term prevents an
identification of these two tensors or an inference of ambient
dimension from the rank of the first.

\section{Marked rational dynamics and transported integer arithmetic}
\label{final:arithmetic}

\begin{theorem}[Rational trees, Euclidean decoding, and the matrix monoid]
\label{final:B1}
Fix the positive rational coordinate, the root $1$, ordered edge
labels $L,R$, and the transformations
\[
 L(x)=\frac{x}{1+x},\quad R(x)=1+x,\qquad
 M_L=\begin{pmatrix}1&0\\1&1\end{pmatrix},\quad
 M_R=\begin{pmatrix}1&1\\0&1\end{pmatrix}.
\]
Successively applying the letters of a word to $1$ gives a rooted
Calkin--Wilf tree~\cite{CalkinWilf} in which every positive rational occurs exactly
once.  The matrices $M_L,M_R$ freely generate exactly the determinant-one
matrices with nonnegative integer entries, including the identity.

For a Stern--Brocot interval~\cite{Stern1858} store the upper endpoint as the first
column and the lower endpoint as the second, starting with columns
$(1,0)^{\mathsf T}$ and $(0,1)^{\mathsf T}$, representing $\infty$ and
$0$.  Its node is their mediant and a left or right step multiplies
this matrix on the right by $M_L$ or $M_R$.  With these conventions
the Stern--Brocot label of a word is the Calkin--Wilf label of the
reversed word.

The complete root path, its oriented subtraction runs, and the
finite continued fraction are mutually decodable, with the canonical
convention
\[
 [a_0;a_1,\ldots,a_k],\quad a_0\geq0,\quad
 a_j\geq1\ (j\geq1),\quad a_k\geq2\ \text{if }k\geq1.
\]
A one-term expansion $[a_0]$ has $a_0\geq1$, so the root is $[1]$.
The last run in the subtraction-to-root algorithm has length one less
than its terminal Euclidean quotient; it has length zero for the root.
\end{theorem}

\begin{proof}
For a reduced fraction $a/b<1$, the unique possible parent is
$a/(b-a)$ and the last edge is $L$; if $a/b>1$, it is $(a-b)/b$ with
last edge $R$.  Both numerators and denominators remain positive,
the gcd is unchanged, and their sum strictly decreases.  The
algorithm terminates only at $a=b$, hence at $1/1$ by coprimality.
It therefore constructs a root path for every positive rational.
At each step both the last edge and parent are forced, proving
uniqueness.  Conversely the child formulas preserve positive
coprimality, so every word has a positive reduced rational label.

For matrix generation write a nonnegative determinant-one matrix as
$M=\left(\begin{smallmatrix}a&b\\c&d\end{smallmatrix}\right)$.
Except for the identity, its two rows are comparable coordinatewise.
Indeed one pattern of incomparable rows, $a<c$ and $b>d$, would
give $ad<bc$.  In the other pattern write $a=c+u$, $d=b+v$ with
$u,v\geq1$.  The determinant equation becomes
\[
 1=cv+ub+uv,
\]
forcing $u=v=1$ and $b=c=0$, which is the identity.  If the lower
row dominates, subtract the upper row to obtain $M=M_LN$; if the
upper dominates, subtract the lower to obtain $M=M_RN$.
The resulting $N$ remains nonnegative integral with determinant one
and strictly smaller entry sum.  Equal rows are impossible.  Repeated
subtraction therefore ends at the identity and proves generation.
The dominance choice at every nonidentity step is unique; cancellation
then proves uniqueness of the generator word.  Conversely products
have nonnegative integral entries and determinant one.  The word
``nonnegative'' is essential: the generators have zero entries.

If $w=s_1\cdots s_\ell$ records chronological path steps, its
Calkin--Wilf column vector is
$M_{s_\ell}\cdots M_{s_1}(1,1)^{\mathsf T}$.  The Stern--Brocot
endpoint matrix is $M_{s_1}\cdots M_{s_\ell}$ and its mediant vector
is that matrix times $(1,1)^{\mathsf T}$.  This proves reversal and
also proves complete Stern--Brocot enumeration.  For example,
$\operatorname{CW}(LR)=3/2$ and $\operatorname{SB}(LR)=2/3$.
Right multiplication by $M_L$ replaces the upper endpoint column
by its sum with the lower, and right multiplication by $M_R$ replaces
the lower by that sum.  Determinant one ensures ordered adjacent
endpoints and a reduced mediant throughout.

Lastly, consecutive subtractions of the same smaller integer from
the larger group into Euclidean integer division.  If $a=qb+r$ with
$0<r<b$, exactly $q$ subtractions replace $(a,b)$ by $(r,b)$.
If $r=0$, coprimality makes $b=1$ and the algorithm stops after
$q-1$ subtractions at $(1,1)$; the case with $b>a$ is identical
with the roles reversed.  Thus oriented runs determine every
quotient, with precisely the stated terminal correction.  The
Euclidean algorithm terminates and its remainders are unique.
Successive identities $a/b=q+r/b$ yield the continued fraction.
If its last digit is one, the identity
$[\ldots,c,1]=[\ldots,c+1]$ removes it; the unique division
quotients give the stated canonical terminal convention.  A rational
value reconstructed from that fraction gives its unique root path
by the already proved parent algorithm, which proves mutual decoding.
\end{proof}

\begin{proposition}[Marked Farey branches and full-domain scalar recovery]
\label{final:B2}
On $[0,1]$ define the Farey map and its inverse branches~\cite{IsolaFarey}
\[
 \psi_0(x)=\frac{x}{1+x},\qquad \psi_1(x)=\frac1{1+x}.
\]
Their respective ranges are $[0,1/2]$ and $[1/2,1]$, and they are
the increasing and decreasing inverse branches of
\[
 F(y)=
 \begin{cases}
 y/(1-y),&0\leq y\leq1/2,\\
 (1-y)/y,&1/2\leq y\leq1.
 \end{cases}
\]
At the common endpoint both forward formulas equal one, and branch
labels retain the inverse-branch information.

The complete numerically labelled Calkin--Wilf tree determines $L,R$
on positive rationals.  If the candidate maps are declared to be
M\"obius transformations, these data determine their full rational
formulas, in particular $L$ on $(-1,\infty)$.  In that marked
extension category the scalar potential is recovered by
\begin{equation}\label{final:tree-integral}
 I(t)=\int_0^{t-1}L(s)\,ds=t-1-\log t,\qquad t>0.
\end{equation}
Connected real analytic continuation is another sufficient extension
category.  A bare tree, or positive-rational action without a suitable
extension category, has no such global converse.
\end{proposition}

\begin{proof}
Derivatives, endpoint values, and substitution verify the ranges,
monotonicity, and the two inverse formulas.  Notice that $\psi_1$ is
not the Calkin--Wilf map $R$.  Every positive rational occurs in the
tree, so its numerical children give both maps on all such inputs.
Two M\"obius transformations agreeing at three distinct finite inputs
where they are defined are equal: cross multiplication gives a
polynomial of degree at most two with three roots, hence the zero
polynomial.  Their fractional-linear maps therefore agree globally
away from their pole.  This supplies the extension $s/(1+s)$ for
$-1<s<0$ as well as for positive $s$, and direct integration proves
\eqref{final:tree-integral}.  Within the real analytic category on
$(-1,\infty)$, equality on positive rationals first gives equality
on $(0,\infty)$ by continuity and then everywhere by the identity
theorem.

Without that category, choose a nonzero smooth bump $\chi$ supported
in $(1/4,1/2)$ and put $\widetilde I=I+\varepsilon\chi$.
The induced generator $\widetilde L(s)=\widetilde I'(1+s)$ agrees
with $L(s)$ for every $s\geq0$, and hence on all positive-rational
inputs, but the potential differs on $(0,1)$.  Small nonzero
$\varepsilon$ even preserves strict convexity, as in
Proposition~\ref{final:F3}.  Finally an abstract unlabelled binary
tree contains no numerical input/output pairs at all, so it cannot
supply the coordinate markings used in this proof.
\end{proof}

The scalar formal-group inverse itself constructs the canonical
global rational expression: $\operatorname{inv}_\star(x)=-x/(1+x)$,
so $L=-\operatorname{inv}_\star$ and $R(x)=1+x$ is the marked shift.
This constructs a labelled rational model.  It does not attach those
labels to an independently supplied abstract tree.

\begin{theorem}[Exact collisions and complete labelled arithmetic transport]
\label{final:B3}
Fix $\mu,\gamma>0$ and a real constant $c_P$, and consider the placed
integer code
\[
 e(n)=\log(1+\gamma n)-\mu n+c_P,\qquad n\geq2.
\]
For distinct $n>m\geq2$,
\begin{equation}\label{final:integer-collision}
 e(n)=e(m)\quad\Longleftrightarrow\quad
 \mu=\frac{1}{n-m}\log\frac{1+\gamma n}{1+\gamma m}.
\end{equation}
The rule $e(a)\odot e(b)=e(ab)$ defines a well-defined binary
operation on the image if and only if $e$ is injective.  Under that
hypothesis, write $S$ for the image with an adjoined abstract unit $1_S$.
Extending $e(1)=1_S$ makes $e:\mathbb N_{>0}\to S$ an isomorphism of labelled multiplicative
monoids.  It transports the full prime, valuation, divisor,
gcd/lcm, M\"obius, and Dirichlet-convolution inventory stated below.

The different intrinsic code
\[
 e_0(n)=H(n+1),\qquad n\geq1,
\]
is injective for every placement and already includes the unit label.
It does not remove any collision of the original code $e$.
\end{theorem}

\begin{proof}
Subtracting the two values of $e$ proves
\eqref{final:integer-collision}; thus avoiding every displayed
pairwise parameter curve is exactly injectivity on the specified
carrier.  If injective, $e$ is a bijection onto its image and ordinary
multiplication transports through that bijection, proving
well-definedness.

For the reverse implication suppose $n>m$ and $e(n)=e(m)$.  Then
\[
 \mu(n-m)=\int_m^n\frac{\gamma}{1+\gamma u}\,du.
\]
For every real $k>1$, the same difference calculation and substitution
in the integral give
\begin{align*}
 e(kn)-e(km)
  &=k\int_m^n\left(\frac{\gamma}{1+k\gamma u}
                        -\frac{\gamma}{1+\gamma u}\right)du<0.
\end{align*}
Here the formula for $e$ is evaluated on the positive real ray in
this calculation.  The strict inequality follows pointwise from
$u>0$, $\gamma>0$, and $k>1$.  Taking an integer $k\geq2$ shows that
multiplication of the common image $e(n)=e(m)$ by $e(k)$ would
require two distinct outputs.  Hence no operation with the prescribed
rule exists.  This proves necessity for every collision, not only
for a selected example.  For instance, $\gamma=1$ and
$\mu=\log(5/4)$ give $e(3)=e(4)$ but
$e(6)-e(8)=\log(175/144)>0$.

When injectivity holds, adjoin an element $1_S$ distinct from every
point of the original image and set $1_S\odot x=x\odot1_S=x$.
The extended map is a bijection preserving multiplication and unit.
The integer fundamental theorem of arithmetic consequently gives a
unique finite factorization of each $x\in S$ into irreducibles,
with the empty product for the unit.  All the inventory formulas
below follow from that factorization, as verified after the table.
Finally $H'(t)=1/t-1<0$ for $t>1$, so $H(n+1)$ is strictly decreasing
on positive integers.  This proves the alternative's injectivity.
For the placed profile with arbitrary additive level, write
$\widetilde\Sigma_P(r)=\log(1+\gamma r)-\mu r+c_P$;
its parameters here are $\mu,\gamma>0$ and $c_P\in\R$.
The critical point is
$r_*=1/\mu-1/\gamma$ and
\[
 \widetilde\Sigma_P(r_*+n/\mu)=H(n+1)
                   +\log(\gamma/\mu)+\mu/\gamma-1+c_P.
\]
Indeed its intrinsic coordinate $\mu(r+1/\gamma)$ is $n+1$ there.
This is a shifted sampling code, with a fixed additive level, rather
than the original sampling at $r=n$.
\end{proof}

Write $x\mid y$ for divisibility in the unit-adjoined monoid, let
$q$ range over its irreducibles, and write $\mu_S$ for the arithmetic
M\"obius function, to distinguish it from the placement parameter.
The complete inventory is
\[
\begin{array}{ll}
 \text{Irreducibles}&q=e(p),\quad p\text{ a numerical prime};\\[2pt]
 \text{Valuations}&v_q(x)=\max\{j\geq0:q^j\mid x\};\\[2pt]
 \text{Factor counts}&\Omega(x)=\sum_qv_q(x),\quad
             \omega(x)=\#\{q:v_q(x)>0\};\\[2pt]
 \text{Greatest common divisor}&
       v_q(\gcd(x,y))=\min(v_q(x),v_q(y));\\[2pt]
 \text{Least common multiple}&
       v_q(\operatorname{lcm}(x,y))=\max(v_q(x),v_q(y));\\[2pt]
 \text{Divisor lattice}&
       \{d:d\mid x\}\cong\prod_{q\mid x}\{0,\ldots,v_q(x)\};\\[2pt]
 \text{Divisor count}&\tau(x)=\prod_{q\mid x}(v_q(x)+1);\\[2pt]
 \text{M\"obius function}&
       \mu_S(x)=0\ \text{if some }v_q(x)\geq2,
       \quad\mu_S(x)=(-1)^{\omega(x)}\ \text{otherwise};\\[2pt]
 \text{Dirichlet convolution}&
       (f*g)(x)=\sum_{d\mid x}f(d)g(x/d).
\end{array}
\]
All sums and products over the support of a single element are finite,
and $\mu_S(1_S)=1$.  To verify the inventory, divisibility of
$x=\prod_q q^{v_q(x)}$ by $d$ is precisely the componentwise condition
$0\leq v_q(d)\leq v_q(x)$.  This proves the valuation maximum,
gcd/lcm formulas, lattice product, and divisor count.  Irreducibility
means exactly that the exponent vector has one coordinate equal to
one and all others zero; it therefore identifies the numerical
primes under $e$.  Summing $\mu_S$ over divisors leaves only the
square-free choices, giving
\[
 \sum_{d\mid x}\mu_S(d)=(1-1)^{\omega(x)}
 =\begin{cases}1,&x=1_S,\\0,&x\ne1_S.\end{cases}
\]
In any fixed commutative coefficient ring, finite reindexing of
divisor pairs makes convolution associative and commutative with
identity $\delta_{1_S}$.  The last displayed identity says
$\mu_S*\mathbf1=\delta_{1_S}$.  Thus if
$F(x)=\sum_{d\mid x}f(d)$, convolution with $\mu_S$ gives
\[
 f(x)=\sum_{d\mid x}\mu_S(d)F(x/d),
\]
proving M\"obius inversion.  A normalized multiplicative function
satisfies $f(1_S)=1$ and $f(xy)=f(x)f(y)$ for coprime $x,y$.
It is exactly specified by arbitrary values at the marked prime
powers, followed by
$f(x)=\prod_{q\mid x}f(q^{v_q(x)})$.  The operation does not choose
those values.  Over a field, every nonzero function satisfying the
unnormalized multiplicativity equation automatically has unit
value one; the identically zero function is a separate convention.

On the original image omitting one, reflexive divisibility means
$x=y$ or $y=x\odot z$ with $z$ in that image.  Coprime gcds and the
bottom of a full divisor lattice require the adjoined $1_S$.  It
cannot be identified with the real number $\widetilde\Sigma_P(1)$ unless
injectivity has also been checked after adding that index.
Prime permutations preserve the abstract free commutative monoid,
so numerical prime names and addition are extra labels.  With the
integer labels supplied, the totient is also transported by
\[
 \varphi(e(n))=\prod_{p^a\parallel n}p^{a-1}(p-1),
 \qquad\varphi(1_S)=1.
\]
The formula counts residues not divisible by any prime divisor by
inclusion--exclusion: for each square-free divisor $d$ of $n$ there
are $n/d$ multiples of $d$ among $1,\ldots,n$, giving
$n\prod_{p\mid n}(1-1/p)$.  Its numerical factors show why this
function requires the integer labels and their sizes.  No step in
this arithmetic transport identifies a real embedding or the
analytic scalar shape from an unlabelled monoid.

\begin{theorem}[Numerical Euler products and the full multiset inverse]
\label{final:B4}
With the numerical positive-integer labels and their usual
multiplication,
\[
 \zeta(s)=\sum_{n\geq1}n^{-s}
          =\prod_{p\text{ prime}}(1-p^{-s})^{-1},
 \qquad\operatorname{Re}s>1.
\]
If a supplied nonnegative self-adjoint operator $A$ has a complete
simple eigenbasis with eigenvalues $0,1,2,\ldots$, then on that same
domain $\operatorname{Tr}(1+A)^{-s}=\zeta(s)$.

Conversely, let a multiset of real generators $q>1$ have positive
integer multiplicities $m_q$.  Suppose its Euler product
\[
 Z(s)=\prod_q(1-q^{-s})^{-m_q}
\]
converges to a finite positive value for all sufficiently large real
$s$.  The full real function on any such terminal ray uniquely
determines the numerical multiset, including multiplicities.
For a nonempty multiset its least generator $q_0$ and multiplicity
$m_0$ are recovered by
\[
 q_0=\lim_{s\to\infty}(\log Z(s))^{-1/s},\qquad
 m_0=\lim_{s\to\infty}q_0^s\log Z(s).
\]
Removing the entire factor $(1-q_0^{-s})^{-m_0}$ and repeating
recovers every remaining generator.  The empty multiset is exactly
$Z\equiv1$.
\end{theorem}

\begin{proof}
Expand the geometric factors over any finite set of primes.  Unique
factorization gives the sum over exactly the integers supported on
those primes.  Absolute convergence of
$\sum_{n\geq1}n^{-\operatorname{Re}s}$ permits the limit over all
primes and proves the identity.  In the indicated eigenbasis the
operator $(1+A)^{-s}$ is diagonal with entries $(n+1)^{-s}$ and
singular values $(n+1)^{-\operatorname{Re}s}$.  It is trace class
precisely for $\operatorname{Re}s>1$, with the asserted trace.

For the inverse choose a sufficiently large $s_0>0$ at which the
product converges.  Positivity of its factors gives
\[
 \log Z(s_0)=\sum_qm_q[-\log(1-q^{-s_0})]
           =\sum_qm_q\sum_{k\geq1}\frac{q^{-ks_0}}k<\infty.
\]
Tonelli's theorem justifies the expansion.  For each finite $M>1$
every generator $q\leq M$ contributes at least $M^{-s_0}$, counted
with multiplicity.  There are therefore only finitely many such
generators.  In particular the multiset is locally finite down to
one, is countable, and if nonempty has a least $q_0>1$ of finite
multiplicity $m_0$.

For $s\geq s_0$ multiply the logarithmic series by $q_0^s$.
Every term is
$(m_q/k)(q_0/q^k)^s$ and is bounded by its value at $s_0$.
Those bounds sum to $q_0^{s_0}\log Z(s_0)<\infty$.  The ratio
$q_0/q^k$ equals one exactly when $q=q_0,k=1$, and is otherwise
strictly less than one.  Dominated convergence consequently gives
\[
 q_0^s\log Z(s)\longrightarrow m_0,
 \qquad\log Z(s)=m_0q_0^{-s}(1+o(1)).
\]
Taking the negative $1/s$ power gives the first limit and the second
is already proved.  Divide $Z(s)$ by the entire recovered factor,
or equivalently subtract
$-m_0\log(1-q_0^{-s})$ from its logarithm.  The remaining product
still satisfies the same convergence hypothesis.  This removes all
powers of $q_0$ at once, including those that coincide numerically
with another generator.  Repeating is therefore valid even in the
presence of such coincidences.  Every specified generator has only
finitely many predecessors by the bounded-interval argument, so
the iteration reaches all of them.  If none remains the logarithm
is zero; if one remains it is strictly positive.  This proves both
uniqueness and the empty-product assertion.
\end{proof}

The inverse in Theorem~\ref{final:B4} retains the positive Euler-factor
category and the numerical exponential coordinate $q^{-s}$.  It
does not reconstruct an addition law from abstract factors or
resolve any collision in Theorem~\ref{final:B3}.  Across these
sections, every construction and converse has its mathematical
proof and its retained foundations stated explicitly. The series,
topology, matrix, stochastic, and arithmetic development is included
in the complete Lean~4 formalization with no admitted proofs.

\section{Global reconstruction closure and irredundancy}
\label{final:closure}

The complete identifying presentations below reconstruct the same
mathematical object $\Sigma$ exactly. Their different forms do not
make them intrinsically distinct mathematical objects. The global
reconstruction theorem retains each presentation's candidate class,
calibrations, and any specified context.

\begin{definition}[Identification and a retained context]
\label{final:identification}
An identifying presentation consists of a candidate class, its marked
observation map, the specified observed value, and a proved inverse on
that class. A reconstruction recovers the object whose observation was
taken. A canonical construction in a context $C$ is a specified map
$F_C$ giving a presentation of $\Sigma$ in $C$; it identifies an arbitrary supplied
object only if that object satisfies the construction's identifying
conditions. Contexts include the coordinate, ambient category and any
designated argument. The reconstruction of an observed object and the
identification of an arbitrary supplied object therefore have distinct
quantifier structures.
\end{definition}

Let $E_\Sigma$ be the following finite collection of complete presentations,
with every candidate class and calibration stated in the referenced
theorem retained. Grouping presentations in a table row does not discard
their distinct observation maps.
\begin{longtable}{@{}>{\raggedright\arraybackslash}p{.24\textwidth}
>{\raggedright\arraybackslash}p{.49\textwidth}
>{\raggedright\arraybackslash}p{.19\textwidth}@{}}
\toprule
Presentation family&Complete identifying data&Results\\\midrule
\endhead
Intrinsic entries&$H$, $I$, or the positive normalized $p$, in the marked
coordinate&\ref{final:C0-placement}\\
Differential core&Calibrated curvature, Riccati/exact flow, recentering&
\ref{final:C1}--\ref{final:C2}\\
Group core&Normalized logarithm, calibrated cocycle, typed completion&
\ref{final:C3}, \ref{final:C3-completion}\\
Convex/projective core&Anchored Bregman/symmetric slice, exact conjugate,
calibrated self-concordance/Hessian metric, projective/Schwarzian data&
\ref{final:C4}--\ref{final:C7}\\
Gamma transforms&Complete law, Laplace, moments, Mellin line, characteristic
function; entropy-maximizing law&\ref{final:P1}, \ref{final:P7}\\
Poisson&Normalized recurrence, centered CGF, exact oriented KL slice&
\ref{final:P2}\\
Gumbel/Haar&Two marked max identities and anchor, or complete marked
Haar-weighted density&\ref{final:P3}, \ref{final:P3-haar}\\
Linked deficit&Both prescribed observations: deficit law and coordinate
involution, linked to the same candidate potential&\ref{final:P4}\\
Other exact probability nodes&Normalized survival/logistic/Green;
complete Stieltjes data; atom-free size-bias/equilibrium inverse&
\ref{final:P5}, \ref{final:P5-stieltjes},
\ref{final:P5-transforms}\\
Formal/analytic series&Todd unit tower with constant one; complete
Todd, A-hat, L and marked normalized $\chi_y$, $y\ne-1$&
\ref{final:F1}--\ref{final:F3}\\
Tree inverse&Complete marked rooted labelled EGF/Borel PGF and its
analytic inverse germ&\ref{final:P6}\\
Marked rational presentation&Calkin--Wilf/Farey numerical labels and
the specified full-domain M\"obius extension&\ref{final:B1}--\ref{final:B2}\\
Measure/operator observations&Full-line Stein identity; linked mixing
measure samples; complete marked polynomial orthogonality; full integer
tail in the Bernstein class&\ref{final:O1}, \ref{final:O5},
\ref{final:O7}, \ref{final:O6}\\
\bottomrule
\end{longtable}
Formal/local nodes identify their formal germs directly. Identifying a
supplied global function from a germ retains the connected real-analytic
candidate class. The explicit expression also constructs its canonical
global representative. Neither assertion identifies an arbitrary smooth
continuation after that class has been forgotten.

\begin{theorem}[Intrinsic uniqueness and reconstruction closure]
\label{final:global-A}
\label{final:global-B}

There exists a unique calibrated intrinsic analytic object $\Sigma$ on the
marked positive coordinate represented by every complete identifying
presentation in $E_\Sigma$. Every presentation in $E_\Sigma$ reconstructs
$\Sigma$ uniquely within its stated candidate class, and $\Sigma$
reconstructs that presentation exactly.

Hence the complete presentations in $E_\Sigma$ are not intrinsically
distinct mathematical objects: they are exact presentations of the same
intrinsic mathematical object $\Sigma$. In particular, every two
presentations $A,B\in E_\Sigma$ are mutually reconstructible through
$\Sigma$.
\end{theorem}

\begin{proof}
Existence is established by the explicit constructions in the referenced
theorems, and their identifying inverses establish uniqueness in the
corresponding candidate classes. In the formal cases, the retained
continuation class is exactly that of Proposition~\ref{final:F3}.

For each $A\in E_\Sigma$, write
\[
e_A:\Sigma\longrightarrow A,
\qquad
d_A:A\longrightarrow\Sigma
\]
for its encoding and identifying inverse. On their stated solution classes,
\[
d_Ae_A=\mathrm{id}_\Sigma,
\qquad
e_Ad_A=\mathrm{id}_A.
\]
For $A,B\in E_\Sigma$, the reconstruction from $A$ to $B$ is
\[
e_Bd_A:A\longrightarrow B,
\]
with inverse
\[
e_Ad_B:B\longrightarrow A.
\]
Indeed,
\[
e_Ad_Be_Bd_A
=
e_Ad_A
=
\mathrm{id}_A,
\]
and, exchanging $A$ and $B$,
\[
e_Bd_Ae_Ad_B
=
e_Bd_B
=
\mathrm{id}_B.
\]

Thus this is one global reconstruction theorem. Pairwise equivalences are
consequences of the reconstruction through $\Sigma$, not additional
hypotheses.
\end{proof}

\begin{theorem}[Canonical realization fibres]
\label{final:global-E}
For each of the following supplied contexts, $\Sigma$ canonically determines
the stated object uniquely under the indicated identifying conditions:
\begin{enumerate}
\item the probability-convolution completion with marked time one:
      $\mu_r=\operatorname{Gamma}(2r,1)$;
\item the normalized local conservative Sturm--Liouville context:
      the unique marked Laguerre self-adjoint closure;
\item the finite-dimensional spectral/scalar-block matrix context:
      $\Phi_n(X)=\operatorname{tr}X-\log\det X-n$;
\item supplied isotropic $\R^4$ with $t=\|z\|^2/2$:
      the standard Gaussian radial realization;
\item the one-ancestor iid Galton--Watson context:
      the critical Poisson(1) offspring/tree law;
\item the supplied universal complex Thom/Chern-character context:
      the correction $\operatorname{Td}^{-1}$.
\end{enumerate}
The complete time-one, coordinate, rank-one, radial, total-progeny or
universal-line data, respectively, recover the intrinsic object under
the corresponding identifying conditions. These are equivalences in a
fixed context, not identifications after forgetting that context.
\end{theorem}
\begin{proof}
The existence, uniqueness and reverse statements are precisely
Theorems~\ref{final:P8}, \ref{final:O3}, \ref{final:M1},
\ref{final:R1}, \ref{final:P6-gw} and \ref{final:F4}.
Their proofs retain the indicated contextual data. In the last case only
full universal line data identify the series, not a single finite-base
evaluation. The operator spectrum and the fixed squared-resolvent
mixing law likewise do not identify an unmarked stationary coordinate
law, by the explicit isospectral countermodel in the operator section.
\end{proof}

\begin{theorem}[Relative irredundancy and its deletion witnesses]
\label{final:global-C}
\label{final:global-D}
Fix real analytic foundations and the following target convention:
the output is the intrinsic shape, a placed function on a marked $r$
coordinate if requested, and any requested externally selected
realization packets. A packet includes its own type and its selected
object, so that deleting it does not leave an ill-typed datum such as
a bundle without its base. Canonical recipes are fixed, and different
external packets have no unstated identifying map between them.

The sufficient decomposition is one intrinsic certificate, the optional
placement $(\mu,a)$, and the requested external selection packets.
The components $(H,I,p)$ of the intrinsic decomposition of $\Sigma$ are
mutually reconstructible, with no free shape parameter after calibration.
The two free placed-family scalar
coordinates are exactly $\mu>0$ and $0<a<1$.

For a deletion test, remove all alternative complete certificates of the
same intrinsic shape; in particular do not retain a linked realization
whose scalar restriction already identifies $I$. Relax the intrinsic
candidate to anchored real-analytic strictly convex $J$ with
$\int_0^\infty e^{-1-J}=1$. In this declared deletion universe the
decomposition is blockwise irredundant whenever the target distinguishes the
alternatives described in the proof. No assertion of shortest encoding
across arbitrary languages is made.
\end{theorem}
\begin{proof}
Sufficiency is the closure theorem, the exact placed reconstruction,
and the fixed-context construction theorem. For intrinsic irredundancy
put $g(t)=(t-1)^4e^{-t}$ and $h(t)=(t-1)^4e^{-2t}$. They are linearly
independent analytic functions vanishing to order four at one.
For
\[
 M(\epsilon,\delta)=\int_0^\infty
 e^{-1-I(t)-\epsilon g(t)-\delta h(t)}dt
\]
bounded perturbations justify differentiation under the integral,
$M(0,0)=1$, and
$\partial_\delta M(0,0)=-\int p(t)h(t)dt<0$.
The implicit-function theorem gives $\delta(\epsilon)$ preserving
mass. Since $t^2g''$ and $t^2h''$ are bounded, sufficiently small
parameters preserve positive second derivative. The affine anchors
and even unit curvature at one remain fixed, endpoint divergence is
unchanged, and linear independence gives a different shape for
$\epsilon\ne0$. Thus the weakened generic class is not already
another complete certificate of $\Sigma$.

At fixed $a=1/2$, take $\mu=1$ and $2$; at fixed $\mu=1$, take
$a=1/3$ and $2/3$. The intrinsic data agree but the closure points
or curvatures and placed functions differ. This proves both placement
deletion witnesses.

For external packets, keep the whole scalar model and every other
packet fixed. A coordinate relabelling changes a marked presentation
without changing its underlying abstract object. Gamma shapes two and
three have the same integer spectral package but distinct stationary
coordinate laws. Uniform sphere angle and a fixed direction have the
same radius; a Gamma subordinator and $X_r=rT$ have the same time-one
law and different increments. Sample a Borel size and choose a path
or a star to obtain distinct tree laws. Spectral matrix perturbations
preserve the rank-one seed while changing the extension. The two
complex lines on $\mathbb{RP}^2$ in Proposition~\ref{final:F5} retain the
same rational characteristic image and differ integrally. Permuting
prime generators preserves the abstract multiplicative monoid and
changes its numerical labels. Finally $\R^2$ and $\R^3$ separate
spatial dimension; on fixed $\R^2$, squared norm and its square
separate orthogonal composition, and squared norm and twice squared
norm separate the remaining radius scale even inside the additive
class. Each pair is an independent expansion of the same remaining
decomposition data and changes its named target.
\end{proof}

\begin{corollary}[Boundary of scalar reconstruction]
\label{final:global-F}
The complete scalar component alone does not identify a placed
coordinate choice, an arbitrary stationary/vector/process/tree/matrix
realization, actual spaces or bundles and integral characteristic
classes, numerical arithmetic labels on a bare carrier, a designated
spatial observable, spatial dimension, or the radial cancellation rule.
Under supplied global nonnegative orthogonal additivity in dimension
at least two the spatial observable is $c\|x\|^2$. Under the supplied
radial operator and residual cancellation one obtains $D\in\{1,3\}$,
and $D\ge2$ selects three. These are exact boundaries of the stated scalar
reconstruction; additional identifying observations may supply further inverses.
\end{corollary}
\begin{proof}
The deletion witnesses in Theorem~\ref{final:global-D}, together with
Theorem~\ref{final:R2} and Propositions~\ref{final:R3}--\ref{final:R4},
preserve the entire scalar component while varying the designated targets.
A profile different from that of $\Sigma$ also
cancels the radial residual in dimension three. The auxiliary
four-dimensional Gaussian remains fixed in these spatial expansions.
\end{proof}

\subsection{Exact reconstruction and formalization}

Theorem~\ref{final:global-A} proves the global reconstruction of
$\Sigma$ across all complete identifying presentations in $E_\Sigma$.
The typed presentation maps encode this exact reconstruction and its
consequences. Ordinary spectral data retain unitary information, and
canonical realization fibres retain their stated contexts. Forgetting a
coordinate and then selecting a canonical realization does not provide an
identifying inverse for an arbitrary candidate. The countermodels above
establish the exact boundaries of the proved reconstruction theorem.

The complete mathematical development has been formalized in Lean~4 with
no admitted proofs. The formal statements preserve the stated candidate
classes, domains, calibrations, retained data, and identifying conditions
of the corresponding mathematical theorems.

\section{Further Reconstructions}
\label{final:further}

The global reconstruction theorem is proved; the presently known complete
presentations of $\Sigma$ need not be exhaustive. Additional complete
mathematical presentations may remain to be discovered.

The criterion for inclusion is exact bidirectional reconstruction. Within its
stated candidate class, an additional presentation is complete only when it
provides an explicit reconstruction from $\Sigma$, an exact inverse recovering
$\Sigma$, a uniqueness theorem, and boundary or counterexample data identifying
which weaker information fails to reconstruct it. Visual or formal resemblance
alone does not satisfy this criterion.

Every additional complete presentation must reconstruct the same $\Sigma$
exactly and be reconstructed from $\Sigma$ exactly within its stated candidate
class.

\section*{Acknowledgments}

I am deeply grateful to my wife, Malaithong Phommasone, who stood beside me
throughout the years in which this work was developed. Her patience, strength,
and unwavering support accompanied me through the long course of this work.

I am also deeply grateful to my family for their continued encouragement,
understanding, and support throughout this journey. I am grateful as well to
P\'eter T\'im\'ar and his wife, Rebeka, and to the friends who stood beside me along
the way. Their friendship, encouragement, and support meant more than can be
expressed here.


\begin{thebibliography}{9}

\bibitem{SSV}
R.~L.~Schilling, R.~Song, and Z.~Vondra\v{c}ek,
\emph{Bernstein Functions: Theory and Applications},
2nd revised and extended ed., De Gruyter Studies in Mathematics~37,
De Gruyter, Berlin/Boston, 2012.
\href{https://doi.org/10.1515/9783110269338}
     {doi:10.1515/9783110269338}.

\bibitem{AguechJedidi}
R.~Aguech and W.~Jedidi,
\emph{New characterizations of completely monotone functions and Bernstein
functions, a converse to Hausdorff's moment characterization theorem},
Arab. J. Math. Sci. \textbf{25} (2019), no.~1, 57--82.
\href{https://doi.org/10.1016/j.ajmsc.2018.03.001}
     {doi:10.1016/j.ajmsc.2018.03.001}.

\bibitem{CalkinWilf}
N.~Calkin and H.~S.~Wilf,
\emph{Recounting the Rationals},
Amer. Math. Monthly \textbf{107} (2000), no.~4, 360--363.
\href{https://doi.org/10.1080/00029890.2000.12005205}
     {doi:10.1080/00029890.2000.12005205}.

\bibitem{Stern1858}
M.~Stern,
\emph{Ueber eine zahlentheoretische Funktion},
J. Reine Angew. Math. \textbf{55} (1858), 193--220.
\href{https://doi.org/10.1515/crll.1858.55.193}
     {doi:10.1515/crll.1858.55.193}.

\bibitem{IsolaFarey}
S.~Isola,
\emph{On the spectrum of Farey and Gauss maps},
Nonlinearity \textbf{15} (2002), no.~5, 1521--1539.
\href{https://doi.org/10.1088/0951-7715/15/5/310}
     {doi:10.1088/0951-7715/15/5/310};
\href{https://arxiv.org/abs/math/0308017}{arXiv:math/0308017}.

\bibitem{GLN}
F.~Gesztesy, L.~L.~Littlejohn, and R.~Nichols,
\emph{On self-adjoint boundary conditions for singular Sturm--Liouville
operators bounded from below},
J. Differential Equations \textbf{269} (2020), no.~9, 6448--6491.
\href{https://doi.org/10.1016/j.jde.2020.05.005}
     {doi:10.1016/j.jde.2020.05.005};
\href{https://arxiv.org/abs/1910.13117}{arXiv:1910.13117}.

\bibitem{DLMF}
NIST Digital Library of Mathematical Functions,
Release 1.2.7 of 2026-06-15,
F.~W.~J.~Olver, A.~B.~Olde Daalhuis, D.~W.~Lozier,
B.~I.~Schneider, R.~F.~Boisvert, C.~W.~Clark,
B.~R.~Miller, B.~V.~Saunders, H.~S.~Cohl,
and M.~A.~McClain, eds.,
\url{https://dlmf.nist.gov/}.

\end{thebibliography}

\end{document}
